\documentclass[a4paper,12pt]{article}
\usepackage{calc,amsmath,amssymb,amsfonts}
\usepackage[T1]{fontenc}
\usepackage[english]{babel}
\usepackage{xcolor,fancyhdr}
\usepackage[margin=0.7874in,noheadfoot]{geometry}
\usepackage[style=numeric,backend=biber]{biblatex}
\usepackage{enumitem,
hyperref}
\hypersetup{colorlinks=true,allcolors=blue}
% Outline numbering
\setcounter{secnumdepth}{5}
\renewcommand\thesection{\arabic{section}}
\renewcommand\thesubsection{\arabic{section}.\arabic{subsection}}
\renewcommand\thesubsubsection{\arabic{section}.\arabic{subsection}.\arabic{subsubsection}}
\renewcommand\theparagraph{\arabic{section}.\arabic{subsection}.\arabic{subsubsection}.\arabic{paragraph}}
\renewcommand\thesubparagraph{\arabic{section}.\arabic{subsection}.\arabic{subsubsection}.\arabic{paragraph}.\arabic{subparagraph}}
% Pages
\fancypagestyle{Standard}{\fancyhf{}
  \fancyhead[L]{}
  \fancyfoot[L]{}
  \renewcommand\headrulewidth{0pt}
  \renewcommand\footrulewidth{0pt}
  \renewcommand\thepage{\arabic{page}}
}
\pagestyle{Standard}
\title{\bigskip\bigskip\bigskip\bigskip\bigskip\bigskip\bigskip\bigskip\bigskip\bigskip\bigskip\bigskip\centering\Large Linux Commands Summary }
\author{Prepared by : Amit}
\date{29/06/2026}
\begin{document}


\maketitle
\pagebreak


\setcounter{tocdepth}{10}

\tableofcontents
\pagebreak

\section{Miscellaneous commands}

\subsection{py : evaluate python command from command line}
\begin{verbatim}
   -c : python command run
   
     Example : 
     $ py -c "print(2**32-1)"
      4294967295
 
\end{verbatim}


\subsection{shutdown : shutdown the system in a safe way}
\begin{verbatim}

        Shutdown the system in a safe way. You can shutdown
        the machine immediately, or schedule a shutdown using 24
        hour format.It brings the system down in a secure way. When
        the shutdown is initiated, all logged-in users and processes
        are notified that the system is going down, and no further
        logins are allowed.
        Only root user can execute shutdown command
           -r : Requests that the system be rebooted after it has 
                 been brought down.
           -P : Requests that the system be powered off after it has 
                  been brought down.
           -k : Only send out the warning messages and disable logins, 
                 do not actually bring the system down.

        Shut down :
          $ sudo shutdown
        
        How to shutdown the system at a specified time
          $ sudo shutdown 05:00

        System shutdown in 20 minutes from now:
          $ sudo shutdown +20

        How to shutdown the system immediately
          $ sudo shutdown now

        How to make shutdown power-off machine
        Although this is by default, you can still use the -P option
        to explicitly specify that you want shutdown to power off the system.
          $ shutdown -P

        How to broadcast a custom message
          $ sudo shutdown +10 "System upgrade"

\end{verbatim}


\subsection{minicom : RS232 communication on serial devices }
\begin{verbatim}

      minicom is a communication on serial devices, supports script language
      interpreting, capturing to file, multiple users with individual    
      configurations, and more.
    
      $ sudo apt install minicom
      $ minicom -s # for settings, will edit /etc/minirc.dfl file
        serial port settings setup:

    +-----------------------------------------------------------------------+
    | A -    Serial Device      : /dev/modem   # device name                |
    | B - Lockfile Location     : /var/lock                                 |
    | C -   Callin Program      :                                           |
    | D -  Callout Program      :                                           |
    | E -    Bps/Par/Bits       : 115200 8N1                                |
    | F - Hardware Flow Control : Yes                                       |
    | G - Software Flow Control : No                                        |
    | H -     RS485 Enable      : No                                        |
    | I -   RS485 Rts On Send   : No                                        |
    | J -  RS485 Rts After Send : No                                        |
    | K -  RS485 Rx During Tx   : No                                        |
    | L -  RS485 Terminate Bus  : No                                        |
    | M - RS485 Delay Rts Before: 0                                         |
    | N - RS485 Delay Rts After : 0                                         |
    |                                                                       |
    |    Change which setting?                                              |
    +-----------------------------------------------------------------------+
        
        after end of settings setuo we can choose
        "save as dflt"
     
     minicom will control the terminal.
     to choose options : ctrl+a+o : options
                         ctrl+a+q : exit  
     
\end{verbatim}





\subsection{google-chrome : the web browser from Google }
\begin{verbatim}

         $ google-chrome --new-window https://rt-ed.co.il/
         
     --new-window
              If PATH or URL is given, open it in a new window.
    


\end{verbatim}



\subsection{mktemp : Create a temporary file or directory }
\begin{verbatim}

         - Create an empty temporary directory and print its absolute path:
               $ mktemp -d
                  /tmp/tmp.Ubw3bJBvR2


\end{verbatim}

\subsection{clamscan : antivirus }
\begin{verbatim}
     Scan home directory ("~"):
         $ clamscan -r --bell -i -v -o -a ~
     
     -v: Be verbose.
     -a: Show filenames while scan
     -r: Recursive
     -i: Only print infected files.
     -o: Skip printing OK files

     To view the version of the scan engine
         $ clamscan -V # Capital V
           ClamAV 0.103.11/27191/Tue Feb 20 11:25:13 2024
     
     # Base Format
         $ clamscan [options] [directory or file to scan]
     # To scan the current directory  
         $ clamscan 
     # To view the version of the scan engine
         $ clamscan -V # Capital V
     # To scan a specific folder or file
         $ clamscan /path or file
     # To scan a specific path and its sub-folders
         $ clamscan -r /path
     # To scan and display only the infected files
         $ clamscan -r -i /path
     # To save the scan report
         $ clamscan -l /path/log-file-name /path
     # Move the infected files to a different folder
         $ clamscan --move /directory /folder-or-file-to-scan
     # Delete the infected files
         $ clamscan --remove yes /folder-or-file-to-scan
     # Notify the user with a bell sound on identifying infected files
         $ clamscan --bell /folder-or-file-to-scan
     # Different options in single go
         $ clamscan  -r --bell --remove yes -l /path/log-file-name  \
         /folder-or-file-to-scan 

     Freshclam runs as a background service and updates the virus signature
     database in regular intervals. If you would like to manually perform an update,
     execute the below commands in the terminal. 
     This would stop the Freshclam service and updates the db and finally starts
     the service back.

        $ sudo systemctl stop clamav-freshclam.service
        $ sudo freshclam
        $ sudo systemctl start clamav-freshclam.service
     See : man clamscan 
     
\end{verbatim}


\subsection{reboot : Restart or reboot the system.}
\begin{verbatim}

        Using 'reboot' command to restart our Linux system
          $ sudo reboot

        Force Immediate reboot in Linux system:
          $ sudo reboot -f

\end{verbatim}



\subsection{Zeal : Offline documentation browser}
\begin{verbatim}

     Zeal : Offline documentation browser
       for C,C++,Python,BASH etc.

\end{verbatim}

\subsection{wine     : run microsoft software in Linux. }
\begin{verbatim}

     Wine is a free and open-source compatibility
     layer that aims to allow application software and computer
     games developed for Microsoft Windows to run on Unix-like
     operating systems. 
     Wine also provides a software library, named Winelib,
     against which developers can compile Windows applications
     to help port them to Unix-like systems


\end{verbatim}


\subsection{yes     : accept everything... }
\begin{verbatim}

      - Repeatedly output "message":
           yes {{message}}

      - Repeatedly output "y":
           yes

      - Accept everything prompted by the apt-get command:
           yes | sudo apt-get install {{program}}

\end{verbatim}


\subsection{uptime     : tells how long the system is running }
\begin{verbatim}
      Tell how long the system has been running
      - Print current time, uptime, number of logged-in users and other information:
         $ uptime

      - Show only the amount of time the system has been booted for:
         $ uptime --pretty

      - Print the date and time the system booted up at:
         $ uptime --since
\end{verbatim}








\subsection{sleep   : delay for a specified amount of time}
\begin{verbatim}
       sleep NUMBER[SUFFIX]...
       Pause for NUMBER seconds.  
       SUFFIX may be 
          's' for seconds (the default),
          'm' for minutes, 
          'h' for hours or 
          'd' for days.  
        NUMBER need not be an integer. 
        Given two or more arguments, pause for the amount of time
        specified by the sum of their values.
          - Delay in seconds:
              sleep {{seconds}}
              Example: sleep 3
          - Execute a specific command after 20 seconds delay:
              sleep 20 && {{command}}
\end{verbatim}


\subsection{Linux..O.S..version..check }
\begin{verbatim}
     
          • lsb_release -a       #lsb is for LinuxStandardBase: 
          • cat /etc/os-release
          • hostnamectl
          • cat /proc/version
          
\end{verbatim}



\subsection{bc       : "Basic Calculator" - simplified man pages}
\begin{verbatim}
    https://en.wikipedia.org/wiki/Bc_(programming_language)
     
    bc, for basic calculator, is "an arbitrary-precision calculator 
    language" with syntax similar to the C programming language. 
    bc is typically used as either a mathematical scripting language or as
    an interactive mathematical shell. 

         scale = 5 : set precision to 5 numbers after the ".".
    
    bc can be used non-interactively, with input through a pipe. 
    This is useful inside shell scripts:
         $ result=$(echo "scale=2; 5 * 7 /3;" | bc)
         $ echo $result
\end{verbatim}


\subsection{cal       : "CALendar" - basic calendar}
\begin{verbatim}
    https://en.wikipedia.org/wiki/Cal_(command)
    
    see ncal command with sligtly advanced options    
     
     Prints calendar information, with the current day highlighted
       - Display a calendar for the current month:
           cal
       - Display previous, current and next month:
           cal -3
       - Display a calendar for a specific year (4 digits):
           cal 2024
       - Display a calendar for a specific month and year:
           cal 4 2024
       - Display calendar with working list numbers (example):
           cal -W 2024
\end{verbatim}

\subsection{date       : show current date with format}
\begin{verbatim}
       Return current date
       
       $ date          #return current date
          Wed 13 Sep 2023 15:31:52 IDT

       $ date "+%Y"    # year long format
         2023
       $ date "+%y"    # year short format   
         23
       $ date "+%m"    # month
         03
       $ date "+%b"    # 
         Mar  
       $ date "+%d"    # day
         13
       $ date "+%S"    # seconds ?
         33
       $ date "+%H"    # hours
         15
       $ date "+%M"    # minutes
         35
       $ date "+%Y-%m-%d %H:%M"
         2023-09-13 15:36
       
       Example:
       $ x=`date +%Y` # command substitution
       $ echo $x
         2024
       $ x=$(date +%Y)
       $ x=$(date +%Y)
       $ echo $x
         2024

       
\end{verbatim}



\subsection{hdate    : Hebrew DATE}
\begin{verbatim}
     
     Display Hebrew date.
     
     -r --parasha      print weekly reading for Shabbat.
     -h --holidays     print holidays.
     -c                print Shabbat start/end times.

     
     
     Location of Haifa : Haifa       32, 34, 2
        hdate -sRH 02 2024 -l32 -L34 -z2  :  sun rize+sunset (Haifa)+ 
                                             weekly reading+holidays
\end{verbatim}

\section{System administration basics}

\subsection{lspci : List all PCI devices }
\begin{verbatim}
    
     display information about PCI buses in the system and devices connected 
     to them.
    
    -v – Be verbose, Include the Kernel Drivers for each device and Modules
         Modules - drivers in Linux.
    -vv – Be more verbose. Show a more detailed version of the last command.
    -vvv – Be as verbose as possible. Show all the details.

    -t – Show a tree-like diagram containing all buses, bridges, devices
         and connections between them.

    -b – Bus-centric view. Show all IRQ numbers and addresses as seen by the
         cards on the PCI bus instead of as seen by the kernel

   Example : 
   $ lspci -t -v
-+-[10000:e0]-+-17.0  Intel Corporation Alder Lake-P SATA AHCI Controller
 |            +-1c.0  Intel Corporation Device 09ab
 |            \-1c.4-[e1]----00.0  Intel Corporation Device f1aa
 \-[0000:00]-+-00.0  Intel Corporation Device 4601
             +-02.0  Intel Corporation Device 46a8
             +-04.0  Intel Corporation Alder Lake Innovation Platform
              Framework Processor Participant
             +-08.0  Intel Corporation 12th Gen Core Processor Gaussian &
              Neural Accelerator
             +-0e.0  Intel Corporation Volume Management Device NVMe RAID
              Controller
             +-14.0  Intel Corporation Alder Lake PCH USB 3.2 xHCI Host
              Controller
             +-14.2  Intel Corporation Alder Lake PCH Shared SRAM
             +-14.3  Intel Corporation Alder Lake-P PCH CNVi WiFi
             +-15.0  Intel Corporation Alder Lake PCH Serial IO I2C Controller
              #0
             +-15.1  Intel Corporation Alder Lake PCH Serial IO I2C Controller
              #1
             +-16.0  Intel Corporation Alder Lake PCH HECI Controller
             +-17.0  Intel Corporation Device 09ab
             +-1d.0-[01]----00.0  Realtek Semiconductor Co., Ltd.
              RTL8111/8168/8411 PCI Express Gigabit Ethernet Controller
\end{verbatim}

\subsection{kernel..modules : Linux drivers }
\begin{verbatim}
      module - Linux drivers
      drivers are in files named *.ko (Kernel Object)
      Add Code \Functionality to the Linux kernel while it is running.
        $ lsmod                 # list kernel modules (drivers)
        $ insmod module_name.ko # load kernel modules
        $ rmmod module_name     # unload kernel modules

     Example
     
     $ lsmod
      Module                  Size  Used by
      rfcomm                 98304  4
      ccm                    20480  6
      cmac                   12288  3
      algif_hash             12288  1
      algif_skcipher         12288  1
      spi_intel_pci          12288  0
      i2c_hid                40960  1 i2c_hid_acpi
      realtek                36864  1
      i2c_smbus              16384  1 i2c_i801
      idma64                 20480  0
      xhci_pci_renesas       20480  1 xhci_pci
      vmd                    24576  0
      spi_intel              32768  1 spi_intel_pci

\end{verbatim}

\subsection{dmesg..modules : Show kernel messages }
\begin{verbatim}
       dmesg : Show kernel drivers messages
       dmesg is used to examine or control the kernel ring buffer.

       $ dmesg [ -c ] [ -n level ] [ -s bufsize ]

       $ sudo dmsg -c   # -c is for clear after previous messages.
         after USB flash/ USB developement connection we will see 
         USB messages. 

        - Show kernel messages:
          $ dmesg

      - Show kernel error messages:
           $ dmesg --level err

      - Show kernel messages and keep reading new ones, similar to tail -f 
        (available in kernels 3.5.0 and newer):
           $ dmesg -w
   
      - Show how much physical memory is available on this system:
           $ dmesg | grep -i memory

       - Show kernel messages 1 page at a time:
           $ dmesg | less

       - Show kernel messages with a timestamp (available in kernels 3.5.0 
         and newer):
           $ dmesg -T

       - Show kernel messages in human-readable form (available in kernels 
         3.5.0 and newer):
           $ dmesg -H

       - Colorize output (available in kernels 3.5.0 and newer):
           $ dmesg -L


       The program helps users to print out their bootup/device 
       driver messages.
       
       - c: Clear the ring buffer contents after printing.
       -sbufsize: Use a buffer of size bufsize to query the kernel ring 
                  buffer. This is 16392 by default.
       -nlevel : Set the level at which logging of messages is done
                  to the console.
                  For example, -n 1 prevents all messages, expect
                  panic messages, from appearing on the console.
        
        All levels of messages are still written to /proc/kmsg


\end{verbatim}

\subsection{strace  :  trace system calls of a program }
\begin{verbatim}
                           RAM
    |------|           |------------|
    | CPU  | <--|      |  O.S.      |
    |      |    |      |    | system| calls
    |------|    |      |    |       | 
                |----->|  Program   |
              opcode   |------------|
     
     Program communicates with CPU trough opcode and with
     O.S. with system calls , within program there are 
     opcodes and system calls.
     
     Withis assembly code (file *.s) there are opcodes
     and system calls, for example program hello_world.c:
     ----------------------------     
     # include <stdio.h>

       int main(void) {
          printf ("Hello world");
       }
     ---------------------------
     present assembly source file:     
     $gcc -S hello_world.c
     
     $ cat hello_world.s
	.file	"hello_world.c"
	.text
	.section	.rodata
.LC0:
	.string	"Hello world"
	.text
	.globl	main
	...
	movl	$0, %eax
	call	printf@PLT  # here call is SYSTE CALL for Linux Kernell
	movl	$0, %eax
	popq	%rbp
     --------------------------- 
     $ gcc hello_world.c -o hello
     $ ./ hello
     
     $ strace ./hello
      execve("./hello", ["./hello"], 0x7ffd79c58fc0 /* 46 vars */) = 0
      brk(NULL)                               = 0x64c5259e8000
      arch_prctl(0x3001 /* ARCH_??? */, 0x7ffe171e4d90) = -1 EINVAL 
      (Invalid argument)
      mmap(NULL, 8192, PROT_READ|PROT_WRITE, MAP_PRIVATE|MAP_ANONYMOUS, -1, 0)    
      = 0x743fb8c0c000
      access("/etc/ld.so.preload", R_OK)      = -1 ENOENT (No such file or    
      directory)
      openat(AT_FDCWD, "/etc/ld.so.cache", O_RDONLY|O_CLOEXEC) = 3
      newfstatat(3, "", {st_mode=S_IFREG|0644, st_size=93767, ...}, 
      AT_EMPTY_PATH) = 0
      mmap(NULL, 93767, PROT_READ, MAP_PRIVATE, 3, 0) = 0x743fb8bf5000
      close(3)                                = 0

     ---------------------------
     Used for tracing the system calls of a program.
     when it is run in conjunction with a program, 
     it outputs all the calls made to the kernel by the program.

     In many cases, a program may fail because it is
     unable to open a file or because of insufficient memory. 
     And tracing the output of the program will clearly show the 
     cause of either problem.
     
     Usage: 
     $ strace <name of the program>
     
     Example:
     $ strace ls

     this will output a great amount of data on to the screen. 
     If it is hard to keep track of the scrolling mass of data, then there is 
     an option to write the output of strace to a file instead which is done 
     using the -o option:
     
     $ strace -o strace_ls_output.txt ls

      write all the tracing output of 'ls' to the 'strace_ls_output.txt' file.


     $ strace echo "Hello World"
       execve("/usr/bin/echo", ["echo", "Hello World"], 0x7ffc1265a308 
       /* 45 vars */) = 0
       brk(NULL)                               = 0x643239cd2000
       arch_prctl(0x3001 /* ARCH_??? */, 0x7ffc5285b970) = -1 EINVAL 
       (Invalid argument)
       mmap(NULL, 8192, PROT_READ|PROT_WRITE, MAP_PRIVATE|MAP_ANONYMOUS, 
       -1, 0) = 0x75c5d1538000
       access("/etc/ld.so.preload", R_OK)      = -1 ENOENT 
       (No such file or directory)
       openat(AT_FDCWD, "/etc/ld.so.cache", O_RDONLY|O_CLOEXEC) = 3
       newfstatat(3, "", {st_mode=S_IFREG|0644, st_size=93767, ...},    
       AT_EMPTY_PATH) = 0
       mmap(NULL, 93767, PROT_READ, MAP_PRIVATE, 3, 0) = 
       0x75c5d1521000
       close(3)                                = 0
       ....
       ....
     

\end{verbatim}

\subsection{uname : print system information }
\begin{verbatim}
      print all information:
      $ uname -a
        Linux amits-vostro-3520 6.5.0-25-generic #25~22.04.1-Ubuntu SMP 
        PREEMPT_DYNAMIC Tue Feb 20 16:09:15 UTC 2 x86_64 x86_64 x86_64 
        GNU/Linux
      
      -a | all info
      -s | kernel name
      -n | network node hostname
      -r | kernel release
      -v | kernel version
      -m | machine hw name
      -p | processor type
      -i | hardware platform
      -o | operating system
      
      
      
      print the kernel release  
      $ uname -r
        6.5.0-25-generic
    
    
\end{verbatim}





\subsection{service..systemctl : service / deamon }
\begin{verbatim}

     A "service" is a program that starts automatically when you start your  
     computer, and runs in the background.
     
     For example, the "network" service sets up your connection to
     the Internet and keeps it running correctly.
     
     service <service-name> status - Check if a service is running:
     service <service-name> start  - Starting a service:
     service <service-name> stop   - Stopping a service:
     
     For example:
     $ service network [start |stop| restart]
     or
     $ sytemctl [start |stop| restart] network)
     
        $ sudo systemctl status ssh  #check the status
        $ sudo systemctl stop ssh    #stop ssh for this session
        $ sudo systemctl start ssh   #start ssh for this session
        $ sudo systemctl disable ssh #stop ssh from boot to computer
        $ sudo systemctl enable ssh  #start ssh from boot to computer   
     
     
     
     you can also modify services by /etc/init.d/service-name action_taken    
     (status, stop and start too).
     
     look in /etc/init.d/ for the services available.
\end{verbatim}

\section{Bash Features and Scripting}


\subsection{only..bash..features }
\begin{verbatim}
      # ((  ))  $((  ))  $[   ]
      above are features of BASH only 
      not other terminals as ZASH ... ("bashizm")

\end{verbatim}


\subsection{run..BASH..script../bin/bash }
\begin{verbatim}
      #!  ----> named "hash bang" or in short form "shebang"   
       which bash
          /usr/bin/bash
          
       
      
      By default, every text file can be run as BASH script 
      (file permissions must be executable):
      $ cat hi555
         echo hi
      $ ./hi555
         hi
      $ ls -l hi555
        -rwxrw-r-- 1 amit amit 8 Oct  8 22:01 hi555
      
      What happens if we add our own command "cat" that prints “miay” at the 
      top of the PATH: our "cat" now is the first in search order.
      Prepare following script named "cat" : 
            $ nano cat 
            --------------            
            #! /bin/bash
            echo miay !!!
            --------------
       $ PATH=$PWD:$PATH
       $ cat
         miay!!!!

\end{verbatim}

\subsection{"hashbang"..."shebang"    }
\begin{verbatim}
      
      https://en.wikipedia.org/wiki/Shebang_(Unix)
      
      For bash script:
      ---------------- 
      $ which bash      
        /usr/bin/bash
      
      #! symbol in named "hash bang" or shortly "shebang"
      #! /usr/bin/bash   


      For Python:
      -----------      
      https://en.wikipedia.org/wiki/Env
      
      env may also be used in the hashbang line of a script to allow
      the interpreter to be looked up via the $PATH.
      For example, here is the code of a Python script:

            #!/usr/bin/env python3
            print("Hello, World!")
            
      In this example, /usr/bin/env is the full path of the env command.
      The environment is not altered, i.e. the original $PATH will be used.
      Note that it is possible to specify the interpreter without using env, by giving             
      the full path of the python interpreter. A problem with that approach is that on 
      different computer systems, the exact path may be different. 
      By instead using env as in the example, the interpreter is searched for
      and located at the time the script is run
\end{verbatim}

\subsection{stop..bash..script..at...error }
\begin{verbatim}
      
      https://ioflood.com/blog/bash-exit-on-error/
      
      Using ‘set -e’
      The simplest way to make a bash script exit when an error occurs is 
      by using the 'set -e' command at the beginning of your bash script. 
      This command instructs the bash script to stop execution if any 
      command it runs returns a non-zero exit status, which is the 
      convention for indicating an error in bash.
         ----------------
         #!/bin/bash         
         set -e
         ----------------
      
      Using ‘set -o pipefail’
      The ‘set -o pipefail’ command ensures that your script exits
      if any command in a pipeline fails, not just the last one.
      This can be useful in scripts where you’re chaining together 
      multiple commands using pipes.
         ----------------
         #!/bin/bash
         set -o pipefail
         ---------------
      
      Using ‘set -u’
      The ‘set -u’ command causes your script to exit if it tries
      to use an uninitialized variable. This can help catch typos
      and other bugs related to variables.
        -------------
        #!/bin/bash
        set -u 
       --------------
       
     - Execute a specific script and stop at the first [e]rror:
        $ bash -e {{path/to/script.sh}}          

\end{verbatim}

\subsection{bash..debug }
\begin{verbatim}
        
        Execute a specific script while printing each command before executing it:
        Usefull for debug !
           $ bash -x {{path/to/script.sh}}      


        Print verbose information for each step (For Debug).
           $ env -v echo hello world
             executing: echo
             arg[0]= ‘echo’
             arg[1]= ‘hello’
             arg[2]= ‘world’
             hello world
  
        Get warning on undefined variable use :
           $ set -u  # To 
           $ echo [$b]
             []
           $ set -u
           $ echo [$b]
             bash: b: unbound variable     # warning on undefined var.
           $ set +u
           $ echo [$b]
             []

\end{verbatim}




\subsection{shopt     : SHell OPTions manage}
\begin{verbatim}
     
     List shell definitions.
     set/unset options

     Options:
      -p	print each shell option with an indication of its status
      -s	enable (set) each OPTNAME
      -u	disable (unset) each OPTNAME

         • s : set
         • u : unset
         Examples : 
            $ shopt -s globstar           
            
            
         to return failure if glob did not get result::

            $ shopt -s failglob
            $ echo [j-z].txt
                 bash: no match: [j-z].txt
            $ shopt -u failglob
            $ echo [j-z].txt
                 [j-z].txt
           
         to enable ** recursive glob :
            $ shopt -s globstar  // set recursive glob on
            $ shopt globstar
               globstar       	on
            $ echo **/*.txt   # recursive glob: 
                              # find in the current dir (!)
                              # and all sub dirs under it 
              1.txt 2.txt 3.txt 4.txt 5.txt   # files from current dir also
              temp1/a.txt temp1/b.txt temp1/c.txt temp1/d.txt
              temp2/10.txt temp2/1.txt temp2/2.txt temp2/3.txt
              temp2/4.txt temp2/5.txt temp2/6.txt temp2/7.txt
              temp2/8.txt temp2/9.txt
            $ shopt -u globstar
            $ shopt globstar
               globstar       	off
            $ echo **/*.*    # now it will search only in sub dirs
              temp1/a.txt temp1/b.txt temp1/c.txt temp1/d.txt 
              temp2/10.txt temp2/1.txt temp2/2.txt temp2/3.txt 
              temp2/4.txt temp2/5.txt temp2/6.txt temp2/7.txt 
              temp2/8.txt temp2/9.txt
\end{verbatim}



\subsection{set  : Change shell attributes / display shell vars}
\begin{verbatim}
      
      https://phoenixnap.com/kb/linux-set

      $ type set
        set is a shell builtin
      
      Change the value of shell attributes and positional parameters, or
    	display the names and values of shell variables.
      
      set command: warnings for undefined variables.
      set command is similar to shopt.
      
      The set command is a built-in Linux shell command that displays
      and sets the names and values of shell and Linux environment
      variables.
      On Unix-like operating systems, the set command functions within the 
      Bourne shell (sh), C shell (csh), and Korn shell (ksh).     
     
      If we type set without any additional parameters, we will get a list of       
      all shell variables, environmental variables, local variables, 
      and shell functions:
        $ set
      This is usually a huge list. You probably want to pipe it into a pager              
      program to more easily deal with the amount of output:
        $ set | less
      ----------------------------------------------------------------  
      Options to set/unset warning on undefined variable.
      -u  Treat unset variables as an error when substituting.
      +u  Return warning on undefined variable. 
      ----------------------------------------------------------------  
      set -e to exit immediately when there is an error in script
             and don't try to continue run  
      ----------------------------------------------------------------
      
      $ set -u  # To get warning on undefined variable use 
      $ echo [$b]
        []
      $ set -u
      $ echo [$b]
        bash: b: unbound variable     # warning on undefined var.
      $ set +u
      $ echo [$b]
        []
      
      $ help set
	      set: set [-abefhkmnptuvxBCHP] [-o option-name] [--] [arg ...]
    	      Set or unset values of shell options and positional parameters.
    	      Change the value of shell attributes and positional parameters, or
    	      display the names and values of shell variables.
            -u  Treat unset variables as an error when substituting.
            +u  Return warning on undefined variable.  


      Split Strings
      Use the set command to split strings based on spaces into separate 
      variables. For example, split the strings in a variable called myvar, 
      which says "This is a test".
            $ myvar="This is a test"
            $ set -- $myvar
            $ echo $1
            $ echo $2
            $ echo $3
            $ echo $4
             This
             is
             a
             test
      
\end{verbatim}


\subsection{chaining..operators explanation}
\begin{verbatim}
      
      https://www.geeksforgeeks.org/chaining-commands-in-linux/     
      
      Chaining commands in Linux allows us to execute multiple commands 
      at the same time and directly through the terminal. 
      It’s like short shell scripts that can be executed through the terminal 
      directly. Linux command chaining is a technique of merging several 
      commands such that each of them can execute in sequence depending on the    
      operator that separates them and these operators decide how the
      commands will get executed. It allows us to run multiple commands in 
      succession or simultaneously.


 
      Some commonly used chaining operators are as follows:

        Operators           Function

        & (Ampersand)       This command sends a process/script/command to the  
                            background. 
        && (Logical AND)    The command following this operator will only 
                            execute if the command preceding this operator has been  
                            successfully executed. 
        ; (Semi-colon)      The command following this operator will execute even if  
                            the command preceding this operator is not successfully  
                            executed.
        || ( Logical OR)	   The command succeeding this operator is only executed if   
                            the command preceding it has failed.
        | (Pipe)            The output of the first command acts as input to the 
                            second command.
        ! (NOT)             Negates an expression within a command. It is 
                            much like an  “except” statement.
        >,>>,<(Redirection) Redirects the output of a command or a group of commands   
                            to a file or stream.
        &&-|| (AND-OR)      It is a combination of AND OR operator and is similar to 
                            the if-else statement.
        \ (Concatenation)   Used to concatenate large commands over several lines in    
                            the shell.
        () (Precedence)     Allows command to execute in precedence order.
        {} (Combination)    The execution of the command succeeding this operator   
                            depends on the execution of the first command.
\end{verbatim}

\subsection{ |     : pipe - redirect stdout of one cmd to stdin of another cmd.}
\begin{verbatim}


    • “Pipe” redirects standard output of one command to standard 
       input of another command
    • t1, t2, tL are “temporary files” in this explanation.
    • “|” acts as temporary file.
    • Pipe redirects only REGULAR input/output, not error


    keyboard --> command1 ---> | ---> command2 ---> |--> command3---> display
    command1 > t1
    command2 <t1 >t2
    command3 > tL
    
    • Example 1: 
      Redirect output of (print of all files *.log)
      to (grep -i (insensitive) that search lines with “error”)
      and sort all this lines:   
        $ cat *log|grep -i error|sort  
        


    • Example 2:
      Redirect (result of recursive insensitive search ,
      starting with current directory) to (search of lines that not    
      include word “ignored”), the (sort the results with “unic”)
      and send the result to file serious_errors.log:  
        $ grep -ri error .|grep -v “ignored”|sort u> serious_errors.log
        
    • Example 3 (YES command usage):          
      $ cat dir1_rm_inputs.txt
        y
        y
        n
        n
        y
        y
        y
      $ cat dir1_rm_inputs.txt | rm ./dir1/*.* -i

      $ rm -i dir2/*
         rm: remove regular empty file &apos;dir2/6.txt? n
         rm: remove regular empty file &apos;dir2/7.txt? n
         rm: remove regular file &apos;dir2/errors2.log? n
         rm: remove regular file &apos;dir2/errors3.log? n
         rm: remove regular file &apos;dir2/errors.log? n
      $ yes | rm -i dir2/*
         rm: remove regular empty file ‘dir2/6.txt’? rm: remove regular                             
         empty file ‘dir2/7.txt’? rm: remove regular file ‘dir2/
         errors2.log’? rm: remove regular file ‘dir2/errors3.log’? rm: 
         remove regular file ‘dir2/errors.log’?
         
    • Example: automation to enter “yes”  when install: 

      $ yes | sudo apt install some_programs

    • apt has -flag to allow install without interrupts.          
         
         
\end{verbatim}

\subsection{;  : send several commands in BASH }
\begin{verbatim}
     run several commands one after another:
     
     $ echo hello ; sleep 1;echo world
        hello
        world
        
     $ echo A; echo B; echo C
        A
        B
        C   
    This will output A, wait 3 sec., clear line and output B:
    $ printf "A"; sleep 3; printf "\rB\n"
        B 
     
\end{verbatim}

\subsection{||  : logical OR }
\begin{verbatim}
        The command succeeding this operator will only execute if the   
        execution of the command preceding it fails.
    
        Exit status: success: 0, others – fail 
         
        Exit status of grep command:
           Normally the exit status is 
             0 if a line is selected, 
             1 if no lines were selected, 
             2 if an  error occurred.  
	         However, if the -q or --quiet or --silent is used and a line is 
	         selected, the exit status is 0 even if an error occurred. 
	         
	    Short circuiting:   
           x || x || executed    :    stops on first success;
           Example: Do the same task with several alternatives automatically.
           Stop in case of first success in one of tasks.
        Example: 
            $ ls
               latex-document-prototype.tex  try.docx  try.html  try.pdf  try.tex
            $ ll try.pdf || echo "2-nd" || echo "3 - rd"
               -rw-rw-r-- 1 amits amits 256176 Feb 28 12:11 try.pdf
            $ ll 222.txt || echo "2-nd" || echo "3 -rd"
               ls: cannot access ‘222.txt’: No such file or directory
               2-nd

        Example of "turnery if":      
             $ ls 
               latex-document-prototype.tex  try.docx  try.html  try.pdf  try.tex    
	         $ ll try.pdf && echo "OK !" || echo "FAILED !" 
               -rw-rw-r-- 1 amits amits 256176 Feb 28 12:11 try.pdf
               OK !
             $ ll 222.txt && echo "OK !" || echo "FAILED !" 
               ls: cannot access '222.txt': No such file or directory
               FAILED !

\end{verbatim}

\subsection{ $\&\&$ : logical AND}
\begin{verbatim}
      Exit status: success: 0, others – fail.
      
      The command succeeding this operator will only execute if the 
      command preceding it executed successfully

      v && v && executed :   stops on first fail
      Example: run chain of commands automatically, where we interested that 
      all the commands in the chain will run successfully. 
      Stop at first fail.

      Example:
      $ echo "1-st" && ll 222.txt && echo "3-rd"
        1-st
        ls: cannot access ‘222.txt’: No such file or directory
\end{verbatim}


\subsection{backslash : Split commands in BASH }
\begin{verbatim}
     To split the command in command line into 
     number of lines we use back slash  symbol \ :
     $ ls\
     >  -l   
           
\end{verbatim}

\subsection{\& symbol   : start process in the background immediately  }
\begin{verbatim}
     & Symbol
     Another method to run a process in the background is by appending the
     ampersand symbol & to the command you execute. 
     This symbol tells the shell to start the process
     in the background immediately. 
     To just start the command in the background at the beginning: 
     all you need to do is put an ampersand at the end of any Linux command.
     This will allow to use terminal along with running the command:

     Example:

       $ gzip largefile.txt &
      
     
     In this example, the gzip command starts running in the background 
     as soon as you execute it. 
     The process ID (PID) will be displayed on the terminal, 
     allowing you to continue using the shell while the process completes silently.

     It's important to note that when using the & symbol, 
     the process may still produce output, which can interfere with
     the prompt on the terminal. 
     
     Run several commands in background : 
       $ xed & nemo & geany &
         [1] 19420
         [2] 19421
         [3] 19422
     
     
     
     Note: 
     -----
     To avoid this, it's recommended to redirect the output
     to a file or to /dev/null, which discards the output. 
     You can do this by appending > filename or > /dev/null
     to the command, respectively.
     
     
\end{verbatim}



\subsection{bash   : Start an interactive shell session }
\begin{verbatim}
      -------------------------------------------------
      when the variable is defined in following form
      it is used in internal bash !       
      $ x=123 bash # x will be defined in internal bash
                   
      $ echo $x
        123
      $ exit
        exit
      $ echo $x
        2024
      --------------------------------------------------
      - Execute a specific script: 
        Allows to run script without execution permissions !
         
            $ bash {{path/to/script.sh}}      
      
      bash command allows to run additional bash session within parent bash window. 
      To exit it use exit command.

      - Start an interactive shell session:
        $ bash

      - Start an interactive shell session without loading startup   
        configs:
        $ bash --norc


      - Execute a specific script while printing each command before executing it:
        Usefull for debug !
        
        $ bash -x {{path/to/script.sh}}

      - Execute a specific script and stop at the first [e]rror:
        $ bash -e {{path/to/script.sh}}

        Exit bash session : 
        $ exit      
      
      --------------------------------------------
      
      $SHLVL  presents shell level:

      $ bash
      $ echo $SHLVL
        2
      $ bash
      $ echo $SHLVL
        3
      $ exit
        exit
      $ echo $SHLVL
        2
      $ exit
        exit
      $ echo $SHLVL
        1

\end{verbatim}





\subsection{SHELL.variables..vs..Env.varables}
\begin{verbatim}

     https://www.digitalocean.com/community/tutorials/how-to-read-and-set-environmental-and-shell-variables-on-linux

     https://www.digitalocean.com/community/tutorials/how-to-read-and-set-             
     environmental-and-shell-variables-on-linux

      By convention, Env. and Shell variables are  
      usually defined using all capital 
      letters. This helps users distinguish 
      environmental variables within other contexts

      for example command grep uses environmental variables:
       $ man grep
       ...
       ENVIRONMENT
       The behavior of grep is affected by the following environment
       variables.... LC_ALL ... LANG ... GREP_COLOR ...
       ...


      Shell variables :
      -----------------
      variables that are contained exclusively within
      the shell in which they were set or defined. They 
      are often used to keep track of ephemeral data, 
      like the current working directory. Shell variable 
      shouldn’t be passed on to any child processes

      Environmental variables : 
      -------------------------
      variables that are defined for the current shell 
      and are inherited by any child shells or 
      processes. Environmental variables are used to 
      pass information into processes that are spawned 
      from the shell.
      
      Env variables are local to the program/process that use
      by them and they children and are deleted after the parent 
      program run finish.
      For example, environ  for C programs allows to use env. variables. 
      Java uses CLASS PATH and JAVA HOME variables that define where
      to find executable files to run them.

             Creating shell variable.
             ------------------------
      $ TEST_VAR='Hello World!'
      $ echo $TEST_VAR
           Hello World!

      We can see this by grepping for our new variable within  
      the set output:
      $ set | grep TEST_VAR
           TEST_VAR='Hello World!'

      We can verify that this is not an environmental variable
      by trying the same thing with printenv:
        $ printenv | grep TEST_VAR
           # No output should be returned.
        
      Accessing the value of shell or environmental variable:
      $ echo $TEST_VAR
        Hello World!
      As you can see, reference the value of a variable by 
      preceding it with a $ sign. The shell takes this to mean that 
      it should substitute the value of the variable when it comes 
      across this.

      So now we have a shell variable TEST_VAR. 
      It shouldn’t be passed on to any child processes. 
      We can spawn a new bash shell from within our current one to 
      demonstrate:
        $ bash # open child bash session inside the existing one
        $ echo $TEST_VAR
      If we type bash to spawn a child shell, and then try to 
      access the contents of the variable, nothing will be 
      returned. This is what we expected.

      Get back to our original shell by typing exit:
        $ exit
      
      $ TEST_VAR="current BASH var"  # must be without spaces 
                                       around !!!!
      $ echo $TEST_VAR
        current BASH var
      $ bash
      $ echo $TEST_VAR
         # no output
      $ exit  # returen to parent BASH shell
        exit
      $ echo $TEST_VAR
       current BASH var

                Creating Environmental Variables.
                ---------------------------------
      Now, let’s turn our shell variable into an environmental variable. 
      We can do this by exporting the variable. The command to do so 
      is appropriately named:
        $ export TEST_VAR
      This will change our variable into an environmental variable. 
      We can check this by checking our environmental listing again:
        $ printenv | grep TEST_VAR
          TEST_VAR=Hello World!
      Or
        $ export | grep TEST_VAR
          declare -x TEST_VAR="current BASH var"
      This time, our variable shows up. 
      Let’s try our experiment with our child shell again:
        $ bash
        $ echo $TEST_VAR
          Hello World!
      Great! Our child shell has received the variable set by its
      parent. 
      Before we exit this child shell, let’s try to export another       
      variable. 
      We can set environmental variables in a single step like this:
        $ export NEW_VAR="Testing export"
      Test that it’s exported as an environmental variable:
        $ printenv | grep NEW_VAR
         NEW_VAR=Testing export
      Now, let’s exit back into our original shell:
        $ exit
      Let’s see if our new variable is available:
        $ echo $NEW_VAR
        $      
      Nothing is returned.
      This is because environmental variables are only passed
      to child processes. There isn’t a built-in way of setting       
      environmental variables of the parent shell. This is good 
      in most cases and prevents programs from affecting the 
      operating environment from which they were called.

      
      Example of shell vars vs. env. vars: 

      $ a=5
      $ b=7
      $ export b
      $ bash
      $ echo $a
      $
      $ echo $b
        7
      $ exit
       exit
      $ echo $a
        5
      $ echo $b
        7

\end{verbatim}

\subsection{"="      : assign value for Shell/Env. variables. }
\begin{verbatim}
      
      Every shell/env variable is string !
            
      Example of define:
        $ TEST_VAR1 = "current BASH var"   
          TEST_VAR1: command not found  # because of spaces
                                          around "=" !!!    
        $ TEST_VAR1="current BASH var"  # must be entered without
                                          spaces around "=" !!!!
        $ echo $TEST_VAR1
          current BASH var
        $ TEST_VAR2="previous and $TEST_VAR1"
        $ echo $TEST_VAR2
          previous and current BASH var
        $ TEST_VAR3='previous and $TEST_VAR1'
        $ echo $TEST_VAR3
          previous and $TEST_VAR1
        $ a=5
        $ echo $a
          5

\end{verbatim}

\subsection{examples..of..system..vars }
\begin{verbatim}
      
      $ echo `whoami`
       amit
      $ echo $(whoami)
       amit
      $ echo aaa $(whoami) bbb
       aaa amit bbb
      $ echo aaa $USER bbb
       aaa amit bbb
      $ echo [$(pwd)]
       [/home/amit]
      $ echo [$PWD]
       [/home/amit]
      $ echo [$HOME]
       [/home/amit]
      $ echo $HOSTNAME
       amits-vostro-3520
      $ echo $SHELL # current shell name
       /bin/bash

\end{verbatim}

\subsection{\$PATH : path system variable }
\begin{verbatim}
      If the command is not built in command in bash, bash performs the
      search for the command by using path in $PATH variable.
      The search is done according to the order of directories in $PATH
      variable.
      
      It might be a good idea to create a directory ~/scripts to hold your
      scripts. Add the directory to the contents of the PATH variable:
        $ export PATH="$PATH:~/scripts"
\end{verbatim}


\subsection{\$SHELL   : current shell type) }
\begin{verbatim}
      https://unix.stackexchange.com/questions/43499/difference-between-
      echo-shell-and-which-bash
      
      $ echo $SHELL
       /bin/zsh
      $ which bash
       /bin/bash
       The first command tells me which shell will be executed by login
       when you log in. In my case, /bin/zsh.
       The second command tells me the first occurrence in my
       $PATH the bash command can be found.


\end{verbatim}

\subsection{\$SHLVL : SHell LeVeL check shell level }
\begin{verbatim}
      Allows to check the shell's level
      bash command runs additional bash session within parent bash window. 
      To exit it use exit command.

      $SHLVL  presents shell level:

      $ bash
      $ echo $SHLVL
        2
      $ bash
      $ echo $SHLVL
        3
      $ exit
        exit
      $ echo $SHLVL
        2
      $ exit
        exit
      $ echo $SHLVL
        1
\end{verbatim}

\subsection{\$PS1 : Prompt String 1 }
\begin{verbatim}
    $ echo $PS1 
    $ echo $PS1 >> ~/.profile
    
\end{verbatim}




\subsection{\$BASH\_VERSION   : result of last command execution }

\begin{verbatim}
      $ echo "Bash version $BASH_VERSION"
        Bash version 5.1.16(1)-release
\end{verbatim}


\subsection{\$?   : result of last command execution }
\begin{verbatim}

      $ ls -la
        total 772
        drwxrwxr-x 2 amits amits   4096 Feb 28 12:12 .
        drwxrwxr-x 8 amits amits   4096 Feb 28 11:42 ..
        -rw-rw-r-- 1 amits amits 256176 Feb 28 12:11 try.pdf
        -rw-rw-r-- 1 amits amits 148037 Feb 28 11:41 try.tex
      $ echo $?
        0
     https://www.baeldung.com/linux/check-if-command-executed-
     successfully#:~:text=We%20can%20use%20a%20binary%20comparison%20operator%20eq,
     the%20pwd%20command%20displays%20the%20present%20working%20directory.
     
       We can use a binary comparison operator eq to check the
       outcome of the previous command, stored in the $?   
       variable:

         #!/bin/bash
         pwd 
         if [ $? -eq 0 ]; then
            echo "Command Executed Successfully"
         else
            echo "Command Failed"
         fi{verbatim}

\end{verbatim}


\subsection{\$!   : PID of last executed program }
\begin{verbatim}

    Present PID of last program:
    $ geany 1.txt & disown
      [1] 8816
    $ echo $!
      8816


\end{verbatim}


\subsection{\$\$   : current bash Process ID }
\begin{verbatim}

      $ echo $$     # present current BASH pid
        6425

\end{verbatim}


\subsection{\$PROMPT\_DIRTRIM   : PROMPT DIRectory TRIMmer }
\begin{verbatim}

      For people looking for a much simpler solution and don't need           
      the name of the first directory in the path, 
      Bash has built-in support for this using the PROMPT_DIRTRIM 
      variable. For example:

      $ PROMPT_DIRTRIM=2

      If set to a number greater than zero, the value is used as the   
      number of trailing directory components to #retain when 
      expanding the \w and \W prompt string escapes (see Printing a 
      Prompt). Characters removed are replaced with an ellipsis.

\end{verbatim}

\subsection{which     : Locate a program in the user's PATH.}
\begin{verbatim}
       $ type -a which
         which is /usr/bin/which
         which is /bin/which
      
       which is not BASH built in
       which locates a program in the user's $PATH.
       Search the $PATH env. variable and display the location of 
       any matching executables.
       
       - Search the $PATH environment variable and display the location 
         of any matching executables:
           $ which {{executable}}
       - If there are multiple executables which match, display all:
           $ which -a {{executable}}
        Example : 
           $ which -a time
              /usr/bin/time
              /bin/time
\end{verbatim}

\subsection{export   : set or present exported variables}
\begin{verbatim}
      Export command is used to add system variable. 
      It is session specific i.e., when Bash window is closed, 
      it became not valid. If you want you variable to be set 
      always in your shell, add it to .bashrc file:
      
          1. Make sure your .bashrc exists
		      ls -A ~/.bashrc
          2. Add your variable at the bottom
		      export VARIABLE=value
          3. Save, then open a new terminal and verify the variable is set
		      $ echo $VARIABLE
		      value
      
      $ type export
      export is a shell builtin

      To present all exported variables : 
      $ export
            
      Example of set env variable: 
        $ NEW_VAR="Testing export"
        $ export NEW_VAR      
      
      We can set environmental variables in a single step like this:
        $ export NEW_VAR="Testing export" 
 
      If no names are supplied, or if the -p option is given, a
      list of names of all exported variables is displayed:
      $ export | head
        declare -x COLORTERM="truecolor"
        declare -x DESKTOP_SESSION="ubuntu"
        declare -x DISPLAY=":0"
        declare -x GDMSESSION="ubuntu"
      $ export | grep TEST_VAR
        declare -x TEST_VAR="current BASH var"
       
      also see declare that can set exported variable and 
      non exported variable.
\end{verbatim}


\subsection{.bashrc   : (BASH Run Comm.) runs every time new bash opened}
\begin{verbatim}
            .bashrc file location: File bash.bashrc 
            is located in /etc/  folder.
            
            For the current user it must be located at home dir: 
                ~/.bashrc 
                /home/amits/.bashrc
            
            
            
            The .bashrc file is a script file that’s executed 
            when a user logs in.
            
            The file itself contains a series of configurations for the 
            terminal session. This includes setting up or enabling: coloring, 
            completion, shell history, command aliases, exported variables
            and more.
      
            Export command when entered in bash terminal is used 
            to add system variable. It is session specific i.e.,
            when Bash window is closed, it became not valid. 
            If you want you variable to be set 
            always in your shell, (as opposed to the superuser's), 
            add it to .bashrc file:
      
                1. Make sure your .bashrc exists
		            ls -A ~/.bashrc
                2. Add your variable at the bottom
		            export VARIABLE=value
                3. Save, then open a new terminal and verify the 
                   variable is set
		             $ echo $VARIABLE
		               value      
            .bashrc contains also alias commands.
            see https://linuxhandbook.com/linux-alias-command/
            for example add to ~/.bashrc file following text:
               alias lilearn='/home/amits/Documents/RTC/linux_administration/               
                         linux_learn_apps/linux_learn_apps_in_C/
                         tex_file_data_extract'  
           
\end{verbatim}

\subsection{.profile vs. .bash\_profile vs. .bashrc}
\begin{verbatim}
        https://www.baeldung.com/linux/bashrc-vs-bash-profile-vs-profile
         
         The order after login:
         ----------------------
         1.  /etc/profile 
         2. .bash_profile 
         3  .bash_login and .profile (add: echo "Hello .profile")
         
         On every run of bash terminal:
         ------------------------------ 
         1 .bashrc  # runs every time we open new bash
                    (add echo "Hello .bashrc")
         
         Startup files contain commands that are to be executed on 
         shell startup. As a result, the shell executes commands 
         present in these files automatically to set up the shell. 
         This happens even before displaying the command prompt.
         
         1. /etc/profile file
         --------------------
         In an interactive login shell, Bash first looks for the 
         /etc/profile file. If found, Bash reads and executes it
         in the current shell. As a result, /etc/profile sets up
         the environment configuration for ALL users.
              
         2. .bash_profile file
         ---------------------
         Bash then checks if .bash_profile exists in the home 
         directory. If it does, then Bash executes .bash_profile in 
         the current shell. Bash then stops looking for other files 
         such as .bash_login and .profile.
         
         3. .bash_login and .profile
         ---------------------------
         If Bash doesn’t find .bash_profile, then it looks for 
         .bash_login and .profile, in that order, and executes the 
         first readable file only.
         
         4. .bashrc
         ----------
         runs at every terminal run.
         
\end{verbatim}




\subsection{alias   : Creates aliases }
\begin{verbatim}
      alias command.
      ll is alias to ls -alF , see .bashrc file	

      alias is defined at terminal level.
      to allow run at sessin level, add alias to .bashrc file.
      
      Note: alias does not work in bash script !!!


      To see the list of aliases:
      $ alias
        alias alert='notify-send --urgency=low -i "$([ $? = 0 ] && echo        
        terminal || echo error)" "$(history|tail -n1|sed -e '\''s/^\s*[0-9]\+
        \s*//;s/[;&|]\s*alert$//'\'')"'
        alias egrep='egrep --color=auto'
        alias fgrep='fgrep --color=auto'
        alias grep='grep --color=auto'
        alias l='ls -CF'
        alias la='ls -A'
        alias liln='/home/amits/Documents/RTC/linux_administration/       
        linux_learn_apps/linux_learn_apps_in_C/tex_file_data_extract'
        alias ll='ls -alF'
        alias ls='ls --color=auto'


      To define alias: 
      $ alias x=‘echo ‘
      $ x "hello world"
        hello world

      To remove alias:
      $ unalias x
      $ x hello world
       x: command not found

\end{verbatim}

\subsection{unalias   : Remove aliases }
\begin{verbatim}

      Define alias: 
      $ alias x=‘echo ‘
      $ x "hello world"
        hello world

      Remove alias:
      $ unalias x
      $ x hello world
       x: command not found

\end{verbatim}


\subsection{printenv   : Print values of all/specific Env. variables. }
\begin{verbatim}
     Print values of all or specific environment variables
     - Display key-value pairs of all environment variables:
         $ printenv | less

     - Display the value of a specific variable:
         $ printenv USER
           amits

     - Display the value of a variable and end with NUL 
       instead of newline:
        $ printenv --null USER
        amits $
\end{verbatim}


\subsection{env   : Show and set the env. / run a prog. in a modified environment }
\begin{verbatim}

      https://www.geeksforgeeks.org/env-command-in-linux-with-examples/

      https://stackoverflow.com/questions/43793040/how-does-usr-bin-env-work-in-a-linux-shebang-line

      1. env is used to either print environment variables. 
      2. It is also used to run a utility or command in a custom environment.
      3. Improtant : In practice, env has another common use. 
         It is often used by shell 
         scripts to launch the correct interpreter. 
         In this usage, the environment is typically not changed (i.e. used current 
         envirinment).
      
      https://en.wikipedia.org/wiki/Env
      env may also be used in the hashbang line of a script to allow
      the interpreter to be looked up via the $PATH.
      For example, here is the code of a Python script:

            #!/usr/bin/env python3
            print("Hello, World!")
            
      In this example, /usr/bin/env is the full path of the env command.
      The environment is not altered, i.e. the original $PATH will be used.
      Note that it is possible to specify the interpreter without using env, by giving             
      the full path of the python interpreter. A problem with that approach is that on 
      different computer systems, the exact path may be different. 
      By instead using env as in the example, the interpreter is searched for
      and located at the time the script is run


      a. Without any argument : print out a list of all environment variables. 
          $ env
        
      
      b. With -i or –ignore-environment or only – : Ignores the inherited environment
         and starts with a clean slate :
          $ env -i your_command          
         This completely clear out the environment, but it does not prevent
         your_command setting new variables.  
      
      c. With -u or –unset: Unsets a variable.Remove variable from the environment
          $ env -u variable_name    

      
      d. With -0 or –null: End each output line with NULL, not newline
          $ env -0
      
      
      e. With -v:	Prints verbose information for each step (For Debug).
           $ env -v echo hello world
             executing: echo
             arg[0]= ‘echo’
             arg[1]= ‘hello’
             arg[2]= ‘world’
             hello world

      
      env [OPTION]... [-] [NAME=VALUE]... [COMMAND [ARG]...]
      

      
      Script Debugging
      -x  Print commands and their arguments as they are executed.
      The set command is especially handy when you are trying to debug your             
      scripts. Use the -x option with the set command to see which command in 
      your script is being executed after the result, allowing you to 
      determine each command's result.

      Script Exporting
      -a  Mark variables which are modified or created for export.
      Automatically export any variable or function created using the -a 
      option. Exporting variables or functions allows other subshells and 
      scripts to use them.

      Exit When a Command Fails
      -e  Exit immediately if a command exits with a non-zero status.
      Use the -e option to instruct a script to exit if it encounters an error       
      during execution. Stopping a partially functional script helps prevent 
      issues or incorrect results.

      Prevent Data Loss
      -C  If set, disallow existing regular files to be overwritten
          by redirection of output.
      The default Bash setting is to overwrite existing files. However, using 
      the -C option configures Bash not to overwrite an existing file when 
      output redirection using >, >&, or <> is redirected to that file.

      Report Non-Existent Variables
      -u  Treat unset variables as an error when substituting.
      The default setting in Bash is to ignore non-existent variables and work 
      with existing ones. Using the -u option prevents Bash from ignoring 
      variables that don't exist and makes it report the issue.




      Using + rather than - causes these flags to be turned off.


      $ type env
        env is hashed (/usr/bin/env)


      - Show the environment variables:
         $ env

       - Run a program. 
       Often used in scripts after the shebang (#!) for looking up the
       path to the program:
         $ env {{program}}

       - Clear the environment and run a program:
         $ env -i {{program}}

       - Remove variable from the environment and run a program:
         $ env -u {{variable}} {{program}}

       - Set a variable and run a program:
         $ env {{variable}}={{value}} {{program}}

       - Set multiple variables and run a program:
         $ env {{variable1}}={{value}} {{variable2}}={{value}} 
           {{variable3}}={{value}} {{program}}

\end{verbatim}


\subsection{declare   : Declare variables and give them attributes }
\begin{verbatim}
      Declare different types of vars 
      see: https://tldp.org/LDP/abs/html/declareref.html      
       
      allows to cancell "export" variable  
      when using -x to set export , or +x to cancell export
    
       - Declare a string variable with the specified value:
           $ declare {{variable}}="{{value}}"

       - Declare an integer variable with the specified value:
           $ declare -i {{variable}}="{{value}}"

        declare -i number
        # The script will treat subsequent occurrences of "number"
          as an integer.		

          $ number=3
          $ echo "Number = $number"     # Number = 3

          $ number=three
          $ echo "Number = $number"     # Number = 0
          # Tries to evaluate the string "three" as an integer.
          Certain arithmetic operations are permitted for 
          declared integer variables without the need for expr or let.
          $ n=6/3
          $ echo "n = $n"       # n = 6/3
          $ declare -i n
          $ n=6/3
          $ echo "n = $n"       # n = 2

       - Declare an array variable with the specified value:
           $ declare -a {{variable}}=({{item_a item_b item_c}})

       -x export declares a variable as available for exporting outside 
          the environment of the script itself.
            $ declare -x var3
            $ declare -x var3=373

        The declare command permits assigning a value to a variable in 
        the same statement as setting its properties.

        $ type declare
          declare is a shell builtin
        
        To set exported variable :
        $ TEST_VAR="current BASH var"
        $ declare -x TEST_VAR   # set exported - env var     
          or  
        $ declare -x TEST_VAR="current BASH var"
        
        To unset exported variable :  
        $ declare + x TEST_VAR   # set non exported - shell var
        $ export | grep TEST_VAR
        $        

        $ echo Enter some text
        $ read myvar
        $   Hello world
        $ echo '$myvar' now equals $myvar 
            $myva now equals Hello world      

\end{verbatim}


\subsection{"."vc'.'    : Quoting}
\begin{verbatim}

          Normally BASH interprets spaces as argument separators:
              $ echo Hello                 World
                Hello World
                # Argument 0 is : echo
                # Argument 1 is : Hello
                # Argument 2 is : World

          " " used to prevent the shell from interpreting spaces 
                as argument separators. The whole expression
                within quotes ".." is considered as one argument.
          All characters between " " are the same argument.
          
          Also " " disables globing, but does not disable ! and $ :  
               $ echo "You are logged as $USER"
                  You are logged as amits
               $ echo *
                  latex-document-prototype.tex try.docx try.html try.pdf try.tex
               $ echo "*"
                 *
               $ echo *.pdf  
                 try.pdf
               $ echo "*.pdf"  
                 *.pdf
          Single quotes preserve all special characters, while double quotes
          allow some special characters to be interpreted:
               $ echo "You are logged as $USER"
                  You are logged as amits
               $ echo 'You are logged as $USER'   
                  You are logged as $USER
          see : https://saturncloud.io/blog/difference-between-single-
                and-double-quotes-in-bash/         
                 
\end{verbatim}

\subsection{"expansion"      : made in echo command }
\begin{verbatim}
      
      echo     $USER
        |        | 
        v        v
      echo     shmuel
        |        |  
        V        v   
   argument0   argument1
   command !    
  1. check ALIAS
  2. check PATH
  3. check BUILT IN   
   
      
      $ a=Hello
      $ b=World
      $ echo "$a $b"
       Hello World   
      "echo" is not path - search for command named "echo".
      Firsly expansion is made.
      After it Hello is substituted as first argument 
      World is substituted as second argument.  
      
      $ echo $a$b
       ABCXYZ
      $ echo "$a""$b"
       ABCXYZ
      $ echo "$a"123"$b"
       ABC123XYZ


      $ a=ec
      $ b=ho
      $ c=hel
      $ d=lo
      $ $a$b $c$d
        hello

      $ some_path="aaa/bbb/ccc"
      $ extended_path="$some_path"/ddd
      $ echo $extended_path
        aaa/bbb/ccc/ddd
      $ extended_path1=ddd/"$some_path"
      $ echo $extended_path1
        ddd/aaa/bbb/ccc
      
      $ # print undefined variable (default - NOT ERROR, empty value):
      $ echo [$yyyy]
        []      
      $ a=
      $ echo [$a]  
        []
      $ a=""
      $ echo [$a]
        []
           
      Problematic example:
      $ a=123
      $ a=$aABC      # problematic: $a here is not separate veriable !
      $ echo [$a]    # BASH thinks that $aABC is not defined
      $ []     
      Fixed example : 
      $ a=123      
      $ a=${a}ABC    # now it is fixed !
      $ echo $a    # BASH thinks that $aABC is not defined
        123ABC 
        
      $ echo "hello $USER"
        hello amits
      $ echo 'hello $USER'
        hello $USER
      
      Problematic example:  
      $ a="Hello        World"
      $ echo $a
        Hello World   # echo gets Hello as first argument 
                      # and Word as second argument
      Fixed example :       
      $ a="Hello        World"
      $ echo "$a"     # now the result of expansion is
                      # considered as one alone string
      $ Hello        World

          
\end{verbatim}


\subsection{`cmd`(backtick) equals \$(cmd): run cmd and substitute its text}
\begin{verbatim}
      Back tick (key beyond ~) `command` or $(command) means:
      firstly, run command that is within `` or within $() and put its       
      output text here. 

      The expression $(command) is a modern synonym for `command` which       
      stands for command substitution;
      
      $ echo whoami
        whoami
      
      $ echo `whoami`
        amits
      
      $ x=`whoami`
      $ echo $x
        amits
      
      $ x=$(whoami)
      $ echo $x
        amits
               
      $ echo `whoami`
        amits

      $ x=$(tldr ls)
      $ echo $x
            ls List directory contents.More information: https://
            www.gnu.org/software/coreutils/ls. - List files one
            per line:  ls -1  - List all files, including hidden
            files:  ls -a  - List all files, with trailing / added
            to directory names:  ls -F  - Long format list 
            (permissions, ownership, size, and modification date) 
            of all files:  ls -la  - Long format list with size       
            displayed using human-readable units (KiB, MiB, GiB):  ls -
       
      In the example below output of command 2 will not be presented on
      the screen, it will act as input of command 1:
         $ # store in x the 
         $ # PID of ping command that is running : 
         $ pgrep ping
           46034      
         $ x=`pgrep ping`
         $ echo $x
           46034          
         $ x=$(pgrep ping)
           46034          
         
         Use of `` vs $() :
         $ # COMMAND1 $(COMMAND2 $(COMMAND3))
         
         Print hello to the user:
         $ echo hello to $(whoami)
           hello to amit
         $ echo hello to `whoami`
           hello to amit         
         $ echo "You are at `pwd`"
           You are at /home/amits/Desktop
           
         Example with problem:
         --------------------- 
         $ my_var=Hi there
         $ echo `$my_var` # after expansion it is `Hi there`
                          # and "Hi" is the first argument
            Hi: command not found
         
         Example without problem:
	   $ var=123
	   $ echo `echo $var`
	    123  
         $ echo $(echo $(echo $var))  
          123 
\end{verbatim}

\subsection{math..\$[...]...\$((...))...((...)) : command to do simple arithmetic}
\begin{verbatim}
       
       $((  )) is the same as $[    ] : performs integer math operation
       
       
       $ echo $[ 1 + 2 ]  # it is not a command
                          # but part of bash features !!!
         3
       $ echo $[1+2]
         3
       $ echo $((1+2))   # it is not a command,
                         # only part of bash features !!!         3
       $ echo $((2*4))
         8
       $ echo $((5/2))
         2
       $ echo $((5>2))
         1  # NOTE : in this case true is 1,
            # it is different than test command.                  
       $ echo $((5<2))
         0
  
            Operator	    Operation
            +, -, *, /	    addition, subtraction, multiply, divide
            %	          Modulus (Return the remainder after division)         
            >,<,>=,<=       Note : the result logic is this case is similar
                                   to C language.
             
       $ x=4
       $ y=5                            
       $ ((z=x+y))     
       $ echo $z
         9
       $ ((x+=y))
       $ echo x
         9
       $ ((x++))
       $ echo x
         10  
            
       Result of ((   )) returns status code:
       $ ((5>17))        # this is COMMAND !!!
       $ echo $?
         1       # regular "bash" logic
       $ ((5<17))
       $ echo $?
         0
        
        
         
       
        ((5>17)) logic of status code is       :  0 -  true,  1 - false
                -----------
       $((5>17)) logic of output expression is :  1 -  true,  0 -  flase   
                -----------------  
       note: if command in bash uses status code !!!
       Two following expressions are equivalent : 
       
       if [ "$x" -gt "$y" ];then echo YES; else echo NO; fi
       if ((x>y))          ;then echo YES; else echo NO; fi
       
       Additional explanation:
       ((0)) --> fault   ---> status = 1
       ((1)) --> sucess  ---> status = 0     
       
       Recommendation:
       ---------------
          use  ((  )) only for boolean conditions :  if ((  )) ...
          use $((  )) only for math :  echo $((3+4)) ; x=$((5+8))  
             
          
          use $(( )) only for math operations
          
       -------------------   
       $ x=6;y=10
       $ ((y=++x))
       $ echo "x=$x, y=$y"
       
       $ echo "x=$x, y=$y"
         x=7, y=7

       $ x=6;y=10
       $ ((y=x++))
       $ echo "x=$x, y=$y"
       
       $ echo "x=$x, y=$y"
         x=7, y=6    # y gets x value before increment of x      
       
       -------------------    
       ((x--))
       ((--x)) 
         
\end{verbatim}



\subsection{let     : math - command to do simple arithmetic}
\begin{verbatim}
      let is a builtin function of Bash that allows us to do simple arithmetic. 
      It follows the basic format:
        $ let <arithmetic expression>
      Here is a table with some of the basic expressions you may perform. 
      There are others but these are the most commonly used.
      
            Operator	    Operation
            +, -, \*, /	    addition, subtraction, multiply, divide
            var++	          Increase the variable var by 1
            var--	          Decrease the variable var by 1
            %	          Modulus (Return the remainder after division)

               $ let a=5+4
               $ echo $a # 9
               $ let "a = 5 + 4"
               $ echo $a # 9
               $ let a++
               $ echo $a # 10
               $ let "a = 4 * 5"
               $ echo $a # 20
               $ let "a = $1 + 30"  
               $ echo $a # 30 + first command line argument   
\end{verbatim}






\subsection{true     : Returns a successful exit status code of 0: \$? = 0}
\begin{verbatim}
      - Return a successful exit code:
         $ true
         
       Example : 
       $ ! true
       $ echo $?
        1
         
\end{verbatim}

\subsection{false     : Returns a non-zero exit code: \$? = 1}
\begin{verbatim}
      - Return a non-zero exit code:
         $ false
         
       Example : 
       $ ! false
       $ echo $?
        0   
\end{verbatim}


\subsection{tab     : autocomplete}
\begin{verbatim}
     Tab 		: automatic completion
     Tab Tab 	: all options for complete
\end{verbatim}










\subsection{type     : Display the type of command the shell will execute.}
\begin{verbatim}

    The ‘type’ is a built-in shell command used to identify the kind of
    command. For example it checks if it is shell built in or external 
    command. 
     
     type without any key : 
                  if NAME is an alias,
    		        shell reserved word, 
    		        shell function, 
    		        shell builtin, 
    		        disk file
    		        or not found, respectively
     	
     -t	output a single word which is one of 
                  'alias',
                  'keyword',
                  'function', 
                  'builtin',
                  'file'

     -p	Returns the disk file that would be executed
     -a	Returns all locations containing an executable.

      For example:
      $ type echo
          echo is a shell builtin
      $ type cat
          cat is /usr/bin/cat
      $ type code
          code is /usr/bin/code     
      $ type -a time
          time is a shell keyword
          time is /usr/bin/time
          time is /bin/time

\end{verbatim}


\subsection{Comments in BASH}
\begin{verbatim}
      comments in BASH : 
          #comment
           
\end{verbatim}



\subsection{source/"." :  Execute commands from a file in the current shell}
\begin{verbatim}
      
      Execute commands from a file in the current shell
      
      - Evaluate contents of a given file:
         $ source {{path/to/file}}

      - Evaluate contents of a given file 
         (alternatively replacing source with .):
         $ . {{path/to/file}}

           
\end{verbatim}


\subsection{read :  retrieving data from stdin}
\begin{verbatim}
      Read a line from the standard input and split it into fields.
      $ read a b c d
        aaa bbb ccc ddd
      $ echo $a
        aaa
      $ echo $b
        bbb
      $ echo $c
        ccc
      $ echo $d
        ddd
      $ bash  # if parameters are not exported , 
              # the child process does not recognize them
      $ echo "x=$x, a=$a, b=$b, c=$c"
        x=456, a=, b=, c=
      $ exit
        exit
      $ echo "x=$x, a=$a, b=$b, c=$c"
        x=456, a=aaa, b=bbb, c=ccc
      $ 

      Shell builtin for retrieving data from stdin
      $ type read
        read is a shell builtin
      $ help read   # man read is something else
        -a array assign the words read to sequential indices of the array
        -p prompt output the string PROMPT without a trailing newline before
    		attempting to read
    	  -s silent (without output to terminal, to enter passwords)	  
    	Example: 	
      $ read -p "What's your name? " name
      $ echo "Hello $name"	
    	
    	If we insert the input at the same line, the prompt does not appear:
      ex_04.sh:
            ---------------------------    	
    	      $ read -p "What's your name? " name 
    	      $ echo "Hello $name" 
    	      ---------------------------
    	$ echo shmuel | bash ex04.sh 
         Hello shmuel	  # without prompt           
\end{verbatim}

\subsection{gnome-terminal :  open new linux terminal emulator  }
\begin{verbatim}
      gnome-terminal is a terminal emulator application for accessing a UNIX shell
      environment which can be used to run programs available on your system.
      
      $ gnome-terminal . &  # may be . is redundant
\end{verbatim}





\subsection{arguments..of..process  }
\begin{verbatim}

      $# - number of arguments
      S0 - name of script file
      $1 - the first argument
      $2 - the second argument
      $* is full list of arguments
      $@ is full list of arguments
      ------------------------------------------
      ex05.sh 
      ------------------------------------------
      echo "this program was run with $0"  # S0 is name of script file
      echo "received $# arguments."        # $# is number of arguments
      echo "first arguemnt is [$1]"        # $1 is the first argument
      echo "second argument is [$2]"       # S2 is the second argument
      echo "third argument is [$3]"

      echo "full list: $*"                 # $* is full list of arguments
      echo "full list: $@"                 # $@ is full list of arguments
      ------------------------------------------
      $ bash ex05.sh aaaa bbbb ccccc
            this program was run with ex05.sh
            received 3 arguments.
            first arguemnt is [aaaa]
            second argument is [bbbb]
            third argument is [ccccc]
            full list: aaaa bbbb ccccc
            
      Note : $# does not take into account the argument 0 (the command itself) !
             it counts arguments without zero argument !!!      
       
       
       $ bash ex05.sh "aaaa bbbb ccccc" # now it treated as single argument !!!
          this program was run with ex05.sh
            received 1 arguments.
            first arguemnt is [aaaa bbbb ccccc]
            second argument is [] 
            third argument is []
            full list: aaaa bbbb ccccc
            full list: aaaa bbbb ccccc

       
       
            
\end{verbatim}


\subsection{clear :Equivalent to CTRL+L : clears the screen of the terminal. }
\begin{verbatim}
      Clears the screen of the terminal.
       - Clear the screen (equivalent to pressing Control-L in Bash shell):
       $  clear
\end{verbatim}


\subsection{process..output..to..stderr }
\begin{verbatim}
           Trick to print ouput to stderr from process:
           ---------------
           ex07.sh:
             echo "normal output"
             echo "error output" >&2            
           --------------- 
             $ bash ex07*  >/dev/null
             $ bash ex07* 2>/dev/null 

\end{verbatim}


\subsection{exit   : exit the process (script) , exit the shell}
\begin{verbatim}

      Example with exit the process:
      ---------------------
      echo "this program will fail"
      exit 1  # exit with status code other than 0 is: error
      # Note:
      #     exit 
      # is the same as:
      #     exit 0     
      ---------------------
      $ bash ex06.sh
        this program will fail
      $ echo $?
        1
   Note : exit status is limited to 255
   
   Example with exit the bash:
      
      $ whoami
        amits
      $ su david
        password:
      $ whoami
        david 
      $ exit    # exit from david
      $ whoami
        amits   
\end{verbatim}

\subsection{if     : if command}
\begin{verbatim}
   
   
    if expression1; then  # if expression1 return status of sucess (exit status 0)
                          # then statement1 will be executed.
         statement1
    elif expression2  # elif (!), not elsif
         then statement2
    else
         statement3
    fi
    
    Note : elif can appear several times
           if and else can appear only once.
    
    
    Examples (also see test command): 
    ---------------------------------
    if ! grep $USER /etc/passwd
           then echo "your user account is not managed locally";
    fi
    
    grep $USER /etc/passwd
    if [ $? -ne 0 ] ; then
       echo "not a local account" ;
    fi
    ---------------------------------
    leap year example:     
    
    #! /usr/bin/bash
    
    if [ -z "$1" ]; then
       echo "Error: no arguments" >&2
       exit 2
    fi

    if (( $# != 1 )); then
      echo "Error: must have exactly 1 argument for year" >&2
      exit 2
    fi

    if (( ($1%4)==0));then
       if (( ($1%100)==0 ));then
          if (( ($1%400)==0 ));then
            echo "$1 % 4 == 0, $1 % 100 == 0 and $1 % 400 == 0 "  
            echo "leap year"
            exit 0
          else
            echo "$1%4==0 and $1%100==0, but not dividable by 400"; 
            echo "common year"
            exit 1
          fi   
       else
          echo "$1%4==0 && $1%100 != 0"; 
          echo "it is leap year"
          exit 0
       fi 
    else
       echo "$1 % 4 != 0"; 
       echo "it is common year"
       exit 1
    fi 
    -------------------------------
    $ bash leap_year_check_upd.sh 2000
    $ echo $?
        0
    $ bash leap_year_check_upd.sh 2002
    $ echo $?
        1
    $ bash leap_year_check_upd.sh 2002 2000
        Error: must have exactly 1 argument for year

    $ if bash leap_year_check.sh 2020;then echo YES;else echo NO;fi
      2020%4==0 && 2020%100 != 0
      it is leap year
      YES

\end{verbatim}


\subsection{test..[...]   : command that can be used as test in if command}
\begin{verbatim}
      
    Evaluate conditional expression.
    
    Exits with a status of 0 (true) or 1 (false) depending on
    the evaluation of EXPR.  Expressions may be unary or binary.  
    Unary expressions are often used to examine the status of a file.  
    There are string operators and numeric comparison operators as well.
    The behavior of test depends on the number of arguments.  
    
               File operators:
               ---------------
      -a FILE        True if file exists.
      -b FILE        True if file is block special.
      -c FILE        True if file is character special.
      
      -d FILE        True if file exists and is a directory.
      -e FILE        True if file exists.
      -f FILE        True if file exists and is a regular file.
              # Note: regular file: not dir and not soft link
      
      -g FILE        True if file is set-group-id.
      -h FILE        True if file is a symbolic link.
      -L FILE        True if file is a symbolic link.
      -k FILE        True if file has its `sticky' bit set.
      -p FILE        True if file is a named pipe.
      -r FILE        True if file is readable by you.
      -s FILE        True if file exists and is not empty.
      -S FILE        True if file is a socket.
      -t FD          True if FD is opened on a terminal.
      -u FILE        True if the file is set-user-id.
      -w FILE        True if the file is writable by you.
      -x FILE        True if the file is executable by you.
      -O FILE        True if the file is effectively owned by you.
      -G FILE        True if the file is effectively owned by your group.
      -N FILE        True if the file has been modified since it was last read.
      
      FILE1 -nt FILE2  True if file1 is newer than file2 (according to              
                       modification date).
      FILE1 -ot FILE2  True if file1 is older than file2.
      FILE1 -ef FILE2  True if file1 is a hard link to file2.
                       (same device and inode numbers)
      
                String operators:
                -----------------
      -z STRING            True if string is empty.
      -n STRING            True if string is not empty.
      STRING1 = STRING2    True if the strings are equal.
      STRING1 != STRING2   True if the strings are not equal.
      STRING1 < STRING2    True if STRING1 sorts before STRING2  
                           lexicographically.
      STRING1 > STRING2    True if STRING1 sorts after STRING2 
                           lexicographically.
      
                Other operators:
                ----------------
      -o OPTION      True if the shell option OPTION is enabled.
      -v VAR         True if the shell variable VAR is set.
      -R VAR         True if the shell variable VAR is set 
                     and is a name reference.
      ! EXPR         True if expr is false.
      EXPR1 -a EXPR2 True if both expr1 AND expr2 are true.
      EXPR1 -o EXPR2 True if either expr1 OR expr2 is true.
    
      arg1 OP arg2   Arithmetic tests.  OP is one of -eq, -ne,
                     -lt, -le, -gt, or -ge.
      -----------------------------------------------------------------
      $ which test
        /usr/bin/test
      
      Options to perform test:
          test EXPRESSION
          test             # always true
          [ EXPRESSION ]   # ] is argument
          [ ]              # always true
            
      
      
      Strings comparison : 
         NOTE : put"" on compared operands when using test command !!!
      
            "="
            "<"
            ">"
            "<="
            ">="
            "!="
            
            -o : or
            -a : and
             ! : not
            
            $ [ "$password_guess" = 1234 -a "$username_guess" = grisha ]
            $ ["$password_guess" = 1234 ] && ["$username_guess" = grisha ]  
                     # check that both commands succeded !!!
            
            Same options:
            ------------- 
            if [! "$password_guess" != 1234 ] && ["$username_guess" = grisha ]; then  
            if ! ["$password_guess" != 1234 ] || ! ["$username_guess" = grisha ]; then 
            if ! ["$password_guess" != 1234 ] && ["$username_guess" = grisha ]; then   
            if ["$password_guess" = 1234 ] && ["$username_guess" = grisha ]; then 
	         echo "Access granted"
            else 
	         echo "Access denied"
	         exit 1
            fi
            ---------------
            
            
            $ [ 1 = 1 -o asd = 1234 ] # note backspaces
            $ echo $?
              0
            
            
            
             
            $ test "AAA" = "AAA" # note spaces around "="
                     $1  $2  $3
   
   
            $ test "AAA" = "AAA" # Note MUST enter spaces between arguments !!!
            $ echo $?
              0
            $ test "AAA" = "BBB"
            $ echo $?
              1
            $ test "AAA" != "AAA"
            $ echo $?
              1
            $ test "AAA" "<" "BBB"
            $ echo $?
              0
            $ test "AAA" ">" "BBB"
            $ echo $?
              1

      Numbers comparison (integer comparison): 
                -lt (less than)
                -eq (equal)
                -ne (not equal)
                -le (less or equal)
                -gt (greater than)
                -ge (greater or equal) 
      
            $ test 100 -lt 50  # less than 
            $ echo $?
              1
            $ test 100 -eq 50  # equal
            $ echo $?
              1
            $ [ 100 -lt 200 ] && echo TRUE || echo FALSE
              TRUE
            $ [ 200 -lt 100 ] && echo TRUE || echo FALSE
              FALSE
            
            
              [            200  -lt  100 ]
              $0(command)  $1    $2  $3  $4 
            
            
            --------------------------------------------  
            Note : put " " around arguments within test/[ ]
            
            Example1 (problematic, enter "1 = 1 -o asd" to ): 
              
              # read passsword from the user
              # the password is 1234
              # if the password is correct print
              # "Acess granted"
              # else print
              # "Acess denied" and exit with failure
              read -p 'enter password:' password_guess
              if [ $password_guess = 1234 ]; then 
	          echo "Access granted"
              else 
	          echo "Access denied"
	          exit 1
              fi
            
            
            Example2 (now it can not be hacked): 
            
              # read passsword from the user
              # the password is 1234
              # if the password is correct print
              # "Acess granted"
              # else print
              # "Acess denied" and exit with failure
              read -p 'enter password:' password_guess
              if [ "$password_guess" = 1234 ]; then   # now password guess 
                                                      # is one single word
                                                      # because of " " 
	          echo "Access granted"
              else 
	          echo "Access denied"
	          exit 1
              fi            
            
          Example 3 : 
            # read username and passsword from the user
            # the password is 1234
            # the username is grisha
            # if the credentials are correct print
            # 	"Access granted"
            # else print
            # 	"Access denied" and exit with failure
            read -p 'enter username:' username_guess
            read -p 'enter password:' password_guess

            if [ "$password_guess" = 1234 \
                 -a "$username_guess" = grisha ]; then 
	         echo "Access granted"
            else 
	         echo "Access denied"
	         exit 1
            fi  
            --------------------------------------
            $ [ ]        # without any text always 1 !
            $ echo $?
              1
            $ [ anyText ] # with any text always 0
            $ echo $?
              0

\end{verbatim}

\subsection{for   : loop command}
\begin{verbatim}
         
          https://www.cyberciti.biz/faq/bash-for-loop/
          
          For does not work as classic prog. languages.
          but using ((   )) we can do the same
          as in C:
          
          Three-expression bash for loops syntax
          --------------------------------------
          This type of for loop share a common heritage with the C       
          programming language. It is characterized by a three-parameter loop 
          control expression; consisting of an initializer (EXP1), a loop-
          test or condition (EXP2), and a counting expression/step (EXP3).
          for (( EXP1; EXP2; EXP3 ))
                ## The C-style Bash for loop ##
                for (( initializer; condition; step ))
                do
                  shell_COMMANDS
                done
          A representative three-expression example in bash as follows:
                  #!/bin/bash
                  # set counter 'c' to 1 and condition 
                  # c is less than or equal to 5
                  for (( c=1; c<=5; c++ ))
                  do  
                     echo "Welcome $c times"
                  done
          
          How do I use for as infinite loops? ↑
          Infinite for loop can be created with empty expressions, 
          such as:
                  #!/bin/bash
                  for (( ; ; ))
                  do
                     echo "infinite loops [ hit CTRL+C to stop]"
                  done
          
          Conditional exit with break 
          ---------------------------
          You can do early exit with break statement inside the
          for loop. You can exit from within a FOR, WHILE or UNTIL
          loop using break. General break statement inside the for loop:
                  for I in 1 2 3 4 5
                  do
                    statements1      #Executed for all values of 
                                     #''I'', up to a disaster-condition
                                     # if any.
                    statements2
                    if (disaster-condition)
                          then
	                        break      #Abandon the loop.
                    fi
                    statements3    #While good and, no disaster-condition.
                   done
          Early continuation with continue statement
          ------------------------------------------ 
          To resume the next iteration of the enclosing FOR, 
          WHILE or UNTIL loop use continue statement.

            for I in 1 2 3 4 5
            do
              statements1      #Executed for all values of
                               # ''I'', up to a disaster-condition if any.
              statements2
              if (condition)
                 then
	              continue   #Go to next iteration of I 
	                         #in the loop and skip statements3
              fi
             statements3
            done
            
            ---------------------------------------
            *           for loop syntax           *
            ---------------------------------------
             Numeric ranges for syntax is as follows:

             for VARIABLE in 1 2 3 4 5 .. N
             do
	        command1
	        command2
	        commandN
             done

             OR

             for VARIABLE in file1 file2 file3
             do
	        command1 on $VARIABLE
	        command2
	        commandN
             done

             OR

             for OUTPUT in $(Linux-Or-Unix-Command-Here)
             do
	        command1 on $OUTPUT
	        command2 on $OUTPUT
	        commandN
             done
            ------------------------------------
               #!/bin/bash
               for i in 1 2 3 4 5
               do
                 echo "Welcome $i times"
               done
            ------------------------------------
               #!/bin/bash
               for i in {1..5}
               do
                 echo "Welcome $i times"
               done
            ------------------------------------ 
             #!/bin/bash
             echo "Bash version ${BASH_VERSION}..."
             for i in {0..10..2}
             do 
               echo "Welcome $i times"
             done
            ------------------------------------
            Loop with strings (helps to divide string to separate words)
            
            #!/bin/bash
            #Say we have a variable, and we need to loop through
            #a list of strings to install those packages:
            TRY_VAR="Mor Rotem Amit"
            for p in $TRY_VAR
            do
               echo "$p is ready"
            done
            
            Result:
            $ bash try_for.sh
               Mor is ready
               Rotem is ready
               Amit is ready
            ----------------------------------
            #!/bin/bash
            ## $@ expands to the positional parameters, starting from one.  
            for i in $@
            do
               echo "Script arg is $i"
            done
            
            Result:
            $ bash try_for_for_args.sh well bad good
               Script arg is well
               Script arg is bad
               Script arg is good
           ---------------------------------- 
           for var in $(command)
           do
             print "$var"
           done
           
           ## example ##
           ## Show the name of every PDF file in the /nas directory: 
           for f in $(ls ./try_dir)
           do
             printf "File is %s\n" "$f"
           done
           
           Result:
           $ bash for_with_ls_command.sh
             File is 1.txt
             File is 2.txt
             File is my_hostname.txt
             File is my_name.txt
\end{verbatim}


\subsection{case   : case command}
\begin{verbatim}

     1. Example of case : 
     
            #!/bin/bash
            if [ $# -lt 2 ];then
                echo "Usage : $0 Signalnumber PID"
                exit
            fi

            case "$1" in
               1) echo "Sending SIGHUP signal"
                  kill -SIGHUP $2
                  ;;
               2) echo "Sending SIGINT signal"
                  kill -SIGINT $2
                  ;;
               9) echo "Sending SIGKILL signal"
                  kill -SIGKILL $2
                  ;;
               *) echo "Signal number $1 is not processed"
                  ;;
            esac     

     2. Example of case :  
 
        #! env bash
        while true; do
	      echo 'Drinks Menu:
		   1) Water
		   2) Coffee
		   3) Cola
		   4) Tea
		   0) Exit
		   '
	       read -p 'enter drink number:' drink_number
  
       	 case "$drink_number" in
	         1|water)     
	              echo "you got water" 
	              ;;
	         2|coffee|java) 
	              echo "you got coffee" 
	              ;;
         	   3|cola) 
         	        echo "you got cola" 
         	        ;;
	         4|tea) 
	              echo "you got tea" 
	              ;;
	         # 0 or q or exit or quit or anything 
	         # that starts with "x"
	         0|q|"exit"|quit|x*)
		        echo "goodbye"
		        break
		        ;;
	           *)
		     echo "ERROR: invalid drink number"
		     ;;
	           esac
        done
        # if--fi case-esac  while-done for-done
 
\end{verbatim}

\subsection{while   : while command}
\begin{verbatim}


     #! /usr/bin/bash
     read -p "Enter base value: " base_value
     read -p "Enter power value: " power_value
     num_of_iterations=1
     result=1;

     while [ $power_value -ge $num_of_iterations ]; do
         echo "num_of_iterations = $num_of_iterations"
         num_of_iterations=$((num_of_iterations+1))
         result=$((result*base_value))
     done

     echo "result = $result"
     exit 

\end{verbatim}

\section{Processes/Jobs Control }

\subsection{process..input..output  }
\begin{verbatim}
        How process gets information:
        -----------------------------
        1. arguments
        2. read file (side effect)
        3. signals (KILL)
        4. env. variables
        5. stdin (0)
        
        How process outputs information:
        --------------------------------
        1. stdout(1)
        2. stderr(2)
        3. write file (side effect)
        4. exit code (0 - well, others - error)

\end{verbatim}

\subsection{arguments..of..process  }
\begin{verbatim}

      $# - number of arguments
      S0 - name of script file
      $1 - the first argument
      $2 - the second argument
      $* is full list of arguments
      $@ is full list of arguments
      ------------------------------------------
      ex05.sh 
      ------------------------------------------
      echo "this program was run with $0"  # S0 is name of script file
      echo "received $# arguments."        # $# is number of arguments
      echo "first arguemnt is [$1]"        # $1 is the first argument
      echo "second argument is [$2]"       # S2 is the second argument
      echo "third argument is [$3]"

      echo "full list: $*"                 # $* is full list of arguments
      echo "full list: $@"                 # $@ is full list of arguments
      ------------------------------------------
      $ bash ex05.sh aaaa bbbb ccccc
            this program was run with ex05.sh
            received 3 arguments.
            first arguemnt is [aaaa]
            second argument is [bbbb]
            third argument is [ccccc]
            full list: aaaa bbbb ccccc
            
      Note : $# does not take into account the argument 0 (the command itself) !
             it counts arguments without zero argument !!!      
       
       
       $ bash ex05.sh "aaaa bbbb ccccc" # now it treated as single argument !!!
          this program was run with ex05.sh
            received 1 arguments.
            first arguemnt is [aaaa bbbb ccccc]
            second argument is [] 
            third argument is []
            full list: aaaa bbbb ccccc
            full list: aaaa bbbb ccccc
 
            
\end{verbatim}


\subsection{Ctrl+Z     : SIGTSTP - an inTeractive SToP signal  }
\begin{verbatim}
       https://medium.com/@prateeksrivastav598/ 
       understanding-ctrl-z-and-ctrl-c-in-the-
       linux-terminal-c28bd1e00dbe
       
       https://www.gnu.org/software/libc/manual/html_node/
       Job-Control-Signals.html
       
       Unlike SIGSTOP, this signal can be handled and ignored.
       Your program should handle this signal if you have
       a special need to leave files or system tables in a
       secure state when a process is stopped.
       For example, programs that turn off echoing should handle   
       SIGTSTP so they can turn echoing back on before stopping.
       
           Ctrl + Z: Suspending a Process
       Ctrl + Z is used to suspend a running process in the terminal.
       This means that you temporarily : 
        1. pause the execution of a program and
        2. move it to the background, allowing you to continue using
           the terminal. 
       It’s useful when you want to halt a process without terminating it.
       Here’s how to use it:
       
           Example 1: Pausing a Long-Running Process
           -----------------------------------------
       Let’s say you’re compressing a large file using the gzip
       command, and it's taking longer than expected.
       You can press Ctrl + Z to pause the compression process and return
       to the terminal prompt.
       This doesn't terminate the gzip command; it simply stops it
       temporarily.
         $ gzip -9 largefile.txt
         [Press Ctrl + Z]
         [1]+  Stopped       gzip -9 largefile.txt
         We should note that even though the shell says the job is Stopped,
         it is, in fact, paused – not terminated – and we’ll be able 
         to resume it. To resume the operation, we can use the fg command.
         So, the gzip process is paused and can be resumed later using
         the fg (foreground) command:
         $ fg
         
           Example 2: Running a Process in the Background.
           -----------------------------------------------
       You can also use Ctrl + Z to start a command in the background.
       For instance, you can run a command like this:
         $ my_long_running_command &
       However, if you forget to use the ampersand (&) to run it in the 
       background, you can press Ctrl + Z to suspend it and then type bg 
       (background) to resume its execution in the background.
       
\end{verbatim}


\subsection{Ctrl+C     : SIGINT <SIgnal INTerrupt> interrupting a process  }
\begin{verbatim}
       https://medium.com/@prateeksrivastav598/ 
       understanding-ctrl-z-and-ctrl-c-in-the-linux-terminal-c28bd1e00dbe
       
                   SIGINT(^C) - The gentle nudge
                   -----------------------------
       SIGINT(^C), or <SIgnal INTerrupt>, is perhaps the most commonly
       encountered signal for many users. This signal is typically associated with the    
       CTRL+C command you often use in your terminal to stop a running process. 
       The primary purpose of SIGINT is to notify a process that the user 
       has requested an interruption.
       Example of SIGINT:
          $ kill -INT <process_id> # SIGINT (interrupt) signal
            or       
          $ kill -2 <process_id>
       
       Example :
          run on first terminal window:
           $ sleep 100  
          run on second terminal window:       
           $ ps aux | grep sleep
     root   9416  0.0  0.0  3212  1792 ?     S  Feb26 0:00 sleep 3600
     amits  0816  0.0  0.0  8376  1920 pts/0 S+ 00:16 0:00 sleep 100
     amits 10822  0.0  0.0  9216  2432 pts/1 S+ 00:16 0:00 grep --color=auto sleep
           $ kill -2 10816  # this will kill the process 
                            # 10816 on first terminal window.
       
       
       Ctrl + C: Interrupting a Process
       Ctrl + C is used to forcefully terminate a running process. 
       When you press Ctrl + C, the terminal sends an interrupt signal (SIGINT)
       to the process, causing it to stop immediately. 
       It’s a quick way to exit a process or command.

       Example: Stopping an Unresponsive Program
          Imagine you’re running a program that becomes unresponsive
          or is stuck in an infinite loop.
          You can use Ctrl + C to interrupt the program and 
          return to the terminal prompt.

             $ ./my_program
             [Press Ctrl + C]
             ^C
             The ^C represents the interrupt signal. 
\end{verbatim}

\subsection{Ctrl+D     : End of File / exit shell  }
\begin{verbatim}
       Note1: ^D is not signal, it is simply End Of File.
              ^D closes Bash
       Sometimes when we try to exit the shell by pressing Ctrl+D
       or by using logout, we may receive an error message:
       $ logout
         There are stopped jobs. # message when there are paused jobs
       
       When bash shows this message, it also prevents logout.
       This is because bash doesn’t allow us to exit the shell while
       there are paused jobs.
       
       The current shell’s process manages its jobs. 
       When we pause a job, it’s left in an incomplete state. If we 
       exit the shell with jobs paused, we might lose some critical data.
       So, we need to take care of these paused jobs before we can exit the shell.

\end{verbatim}


\subsection{ps     : "Process Status" - snapshot of the current processes }
\begin{verbatim}
        
        see https://sites.ualberta.ca/dept/chemeng/AIX-43/share/man/info/C
        /a_doc_lib/cmds/aixcmds4/ps.htm
        
        The ps program (short for "process status") displays info on the 
        currently-running processes. A related Unix utility named "top"
        provides a real-time view of the running processes.   
        
        This version of ps accepts 3 kinds of options:
          1   UNIX options, which may be grouped and must be preceded by a dash.
          2   BSD options, which may be grouped and must not be used with a dash.
          3   GNU long options, which are preceded by two dashes.

        Options of different types may be freely mixed, but conflicts can
       appear.  There are some synonymous options, which are functionally
       identical, due to the many standards and ps implementations that this
       ps is compatible with.
       ---------------------------------------------------------------------------------------
       To see only the processes owned by a specific user on Linux run: 
          $ ps -u {USERNAME}

          $ ps -u  | grep bash # present all processes of CURRENT USER !
            amits       6425  0.0  0.0  12056  5760 pts/0    Ss   11:31   0:00 bash
            amits       7263  0.0  0.0  11408  5120 pts/1    Ss+  12:27   0:00 /bin/bash
            amits      13283  0.0  0.0  10108  3840 pts/0    T    14:49   0:00 bash ex04_1.sh
            amits      15526  0.0  0.0  10108  3840 pts/0    T    15:06   0:00 bash ex04_1.sh
       ---------------------------------------------------------------------------------------
        -a Select all processes except:
             1. session leaders (see getsid(2))
             2. processes not associated with a terminal.        
        -u userlist Select by effective user ID (EUID) or name.
             This selects the processes whose effective user 
             name or ID is in userlist.
             
        -x is for tty : 
            Today in Linux, the tty is a legacy name used to refer
            to the user interface for text-based input and output,
            otherwise known as a "terminal".
            In Linux systems, there can be multiple tty device
            "consoles", to support potentially dozens of serial
            ports or more.
            tty0 is the current one in use, but Linux allows you 
            to switch to another session by changing to a different tty,
            e.g., tty1. Linux (e.g., Ubuntu) supports up to 6 ttys by default,
            but this number is configurable     
           
        EXAMPLES
       To see every process on the system using standard syntax:
          $ ps -e    #  Select all processes.
          $ ps -ef   #  Do full-format listing.
          $ ps -eF   #  Extra full format.
          $ ps -ely  #  -l : BSD Long format.
                        -y : Do not show flags; show rss in place of addr.
                           This option can only be used with -l.
       By running $ ps l, you can see which processes are running, 
       and which are sleeping. If a process is sleeping, the WCHAN column  
       ("wait channel", the name of the wait queue) will tell you what
       kernel event the process is waiting for. 
          $ ps  l # display BSD long format.
          $ ps -l # Long format.  The -y option is often useful with this   
   
       To see every process on the system using BSD syntax:
          $ ps ax
          $ ps aux

       Example : 
          $ ps aux | grep bash
             amits       5464  0.0  0.0  11408  5120 pts/0    Ss+  13:03   0:00 bash
             amits       6966  0.0  0.0  11408  5248 pts/1    Ss   14:32   0:00 bash
             amits       6974  0.0  0.0   9216  2432 pts/1    S+   14:32   0:00 grep 
                                                                   --color=auto bash
       To print a process tree:
          ps -ejH
          ps axjf
       To get info about threads:
          ps -eLf
          ps axms   

    Column         |     Contents
    -------------------------------------------------------------------
    %CPU           |     How much of the CPU the process is using
    %MEM           |     How much memory the process is using
    ADDR           |     Memory address of the process
    C or CP        |     CPU usage and scheduling information
    COMMAND*       |     Name of the process, including arguments, if any
    NI             |     Nice value
    F              |     Flags
    PID            |     Process ID number
    PPID           |     ID number of the process's Parent Process
    PRI            |     Priority of the process
    RSS            |     Resident set size (process memory occupied in RAM)
    S or STAT      |     Process status code
    START or STIME |	    Time when the process started
    VSZ            |     Virtual memory usage
    TIME           |     The amount of CPU time used by the process
    TT or TTY      |     Terminal associated with the process
    UID or USER    |     Username of the process's owner
    WCHAN          |     Memory address of the event the process is waiting for    
    -------------------------------------------------------------------
    Note 1 :   In computing, resident set size (RSS) is the portion of memory 
    (measured in kilobytes) occupied by a process that is held in main memory 
    (RAM). The rest of the occupied memory exists in the swap space or file 
    system, either because some parts of the occupied memory were paged out, or 
    because some parts of the executable were never loaded.    
    Note 2 : nice is a program found on Unix and Unix-like operating systems
    such as Linux. It directly maps to a kernel call of the same name. 
    nice is used to invoke a utility or shell script with a particular CPU
    priority, thus giving the process more or less CPU time than other processes.  
    A niceness of -20 is the lowest niceness, or highest priority. 
    Note 3 :  The SZ field displays the size of the process in memory.
    The value of the field is the number of pages the process is occupying. 
    On Linux machines, a page is 4,096 bytes.
    
       PROCESS STATE CODES 
       https://askubuntu.com/questions/1144036/process-status-of-s-s-s-sl
       https://www.linusakesson.net/programming/tty/index.php
       Here are the different values that the s, stat and state output
       specifiers (header "STAT" or "S") will display to describe the state of
       a process:

               D    uninterruptible sleep (waiting for some event)
               R    running or runnable (on run queue)
               S    interruptible sleep (waiting for some event or signal)
               T    stopped by job control signal
               t    stopped by debugger during the tracing
               Z    defunct "zombie" process, terminated but not yet
                    reaped by 
                    its parent
       A Linux process can be in one of the following states:       
       
            R------------->Z
       D<-->R     S------->Z
            R<--->S--->T-->Z
            R<-------->T
       
       For BSD formats and when the stat keyword is used, ADDITIONAL
       CHARACTERS may be displayed:

               <    high-priority (not nice to other users)
               N    low-priority (nice to other users)
               L    has pages locked into memory (for real-time and custom IO)
               s    This process is a session leader
               l    This process is multi-threaded (using CLONE_THREAD,
                    like NPTL pthreads do)
               +    This process is in the foreground process group
               
        Job control is what happens when you press ^Z to suspend a program,
        or when you start a program in the background using &. 
        A job is the same as a process group. 
        Internal shell commands like jobs, fg and bg can be used to manipulate
        the existing jobs within a session. 
        Each session is managed by a session leader, the shell, which is 
        cooperating tightly with the kernel using a complex protocol of signals
        and system calls.          
  
    --------------------------------------------------------------    
    The default niceness for processes is inherited from its parent process
    and is usually 0.
    
        Users can pipeline ps with other commands, 
        such as less to view the process status output one 
        page at a time:
        
        $ ps -A | less
        
        Users can also utilize the ps command in conjunction with    
        the grep command (see the pgrep and pkill commands) to find 
        information about a single process, such as its id:
        
        $ # Trying to find the PID of `firefox-bin` which is 2701
        $ ps -A | grep firefox-bin
          2701 ?        22:16:04 firefox-bin
    Custom columns choose : 
        $ ps a -o user -o comm -o pid -o nice    
           USER     COMMAND             PID  NI
           amits    gdm-wayland-ses    2342   0
           amits    gnome-session-b    2345   0
           root     agetty             3891   0
           root     agetty             3931   0
           amits    tex_file_data_e   14438   0
           amits    htop              30232   0
           root     sudo              30469   0
           root     sudo              30471   0
           root     geany             30472 -20
           root     bash              30520 -20
           root     dbus-launch       30529 -20

         AIX FORMAT DESCRIPTORS
            This ps supports AIX format descriptors, which work
            somewhat like the            
            formatting codes of printf(1) and printf(3).  
            For example, the normal default output can be 
            produced with this: ps -eo "%p %y %x %c".  The
            NORMAL codes are described in the next section.

            CODE   NORMAL   HEADER
            %C     pcpu     %CPU
            %G     group    GROUP
            %P     ppid     PPID
            %U     user     USER
            %a     args     COMMAND
            %c     comm     COMMAND
            %g     rgroup   RGROUP
            %n     nice     NI
            %p     pid      PID
            %r     pgid     PGID
            %t     etime    ELAPSED
            %u     ruser    RUSER
            %x     time     TIME
            %y     tty      TTY
            %z     vsz      VSZ
       
       Example: 
         $ ps a -o user -o comm -o pid -o nice
       it is the same as :
         $ ps a -o "%U %c %p %n"
       
       
        The use of pgrep simplifies the syntax and avoids potential race 
        conditions:

\end{verbatim}

\subsection{pstree    : display a tree of processes}
\begin{verbatim}
       
       https://www.geeksforgeeks.org/pstree-command-in-
       linux-with-examples/
       
       pstree [options] [pid or username]
       
       A convenient tool to show running processes as a tree
       The pstree command is a handy little utility that shows the 
       processes currently running on your system and it show them in 
       a tree diagram. 
       To display process tree:
          $ pstree
       
       To include command line arguments in output
          $ pstree -a

       To display PIDs
          $ pstree -p

       To force pstree to expand identical subtrees in output.
          $ pstree -c

       To sort processes
          $ pstree -n
       
       To see who is the owner/user of a process
          $ pstree -u
          
       To highlight the current process or any other process
          $ pstree -h  
          
       To show process group IDs in output
          $ pstree -g   
          
       To make pstree display process tree specific to a user.
          $ pstree amits   
           
          
\end{verbatim}



\subsection{pgrep     : ps+grep }
\begin{verbatim}
      Find or signal processes by name

      Search for a Linux process by name run: 
        $ pgrep -u {USERNAME} {processName}
        Example :
        $ pgrep -u amits sleep # only for "amits" user
          19267
        $ pgrep sleep # all users
          18588
          19267


      Return only PID number , pgrep = ps + grep : 
         $ pgrep ping = ps aux | grep -w [p]ing:

\end{verbatim}


\subsection{top     : "Table Of Processes" }
\begin{verbatim}
        
        To show all processes for a specific user using top/htop
            $ top -U amits
                         
        is a task manager or system monitor program,
        found in many Unix-like operating systems, 
        that displays information about CPU and memory utilization
        
        
        The program produces an ordered list of running processes 
        selected by user-specified criteria, 
        and updates it periodically. 
        Default ordering is by CPU usage, 
        and only the top CPU consumers are shown. 
        top shows how much processing power and memory are being used,
        as well as other information about the running processes.
        Some versions of top allow extensive customization of the display,
        such as choice of columns or sorting method.
        top is useful for system administrators, as it shows which
        users and processes are consuming the most system
        resources at any given time.
        
        top view can be listed with help of arrows.
        
           To get help : 
           $ top
           ^Z         # suspend and move the process to bg
           $ man top
           $ fg       # move the process to fg

\end{verbatim}


\subsection{htop     : similar as top, but more user friendly }
\begin{verbatim}
       More user friendly than top, 
       The title is always on the top of the table.
       
       Press F4 to insert filter word,
        for example, "amits" word in filter will cause to present user "amits"
        "bash" word in filter will present processes of bash.       
            
       Press F5 to see process tree !!!
       Press F5 again to present process list.  
        
       NI:  niceness in htop window. 
          The higher the NI value, the lower the priority-
          the "nicer" the process is to other processes.

\end{verbatim}


\subsection{nice    : run a program with a custom priority (niceness) }
\begin{verbatim}
    Niceness values range from -20 (highest priority - most favorable
    to the process) to 19 (lowest priority - least favorable
    to the process).

    ◦ The nice value determines the priority of a process, 
      with lower values indicating higher priority.
    
     NOTE:  your shell may have its own version of nice, 
     which usually supersedes the version described here.
     
     Please refer to your shell's documentation  for
     details about the options it supports.
     
     $ type nice
        nice is /usr/bin/nice # not built in
     
     
     nice [OPTION] [COMMAND [ARG]...]

     - Launch a program with altered priority:
         $ nice -n {{niceness_value}} {{command}}
       Example: 
         $ sudo nice -n -20 geany
           [sudo] password for amits: **** 
 
    https://www.scaler.com/topics/linux-nice/
    ◦ The Nice command in linux allows users to prioritize
      process execution.
    ◦ The nice value determines the priority of a process, 
      with lower values indicating higher priority.
    ◦ To check the nice value of a running process, 
      use the ps -l <process_id> command.
    ◦ execute a process with given nice value : 
        $ nice -n <nice_value> <command> syntax.  

        $ sudo nice -n -20 geany  # this caused to 
            [sudo] password for amits: 
        $ ps a -o user -o comm -o pid -o nice    
           USER     COMMAND             PID  NI
           amits    gdm-wayland-ses    2342   0
           amits    gnome-session-b    2345   0
           root     agetty             3891   0
           root     agetty             3931   0
           root     agetty             3937   0
           root     login              4195   0
           amits    bash               4294   0
           amits    bash               5964   0
           amits    bash               6101   0
           amits    tex_file_data_e   14438   0
           amits    htop              30232   0
           root     sudo              30469   0
           root     sudo              30471   0
           root     geany             30472 -20
           root     bash              30520 -20
           root     dbus-launch       30529 -20
\end{verbatim}


\subsection{renice   : alters priority of already running process }
\begin{verbatim}
      The Renice command can be used to alter a process's priority while it is     
      already running. Open a terminal and execute the following command:
         $ renice <nice_value> -p <process_id>
      Replace <nice_value> with the new priority level you want to assign    
      and <process_id> with the ID of the running process you wish to modify.
      For example, to increase the priority of a process with ID 1234 to 5, you 
      would run:
         $ renice 5 -p 1234

     Niceness values range from -20 (most favorable to the  
     process) to 19 (least favorable to the process)

     - Change priority of a running process:
         $ renice -n {{niceness_value}} -p {{pid}}

     - Change priority of all processes owned by a user:
         $ renice -n {{niceness_value}} -u {{user}}

     - Change priority of all processes that belong to 
       a process  group:
         $ renice -n {{niceness_value}} --pgrp {{process_group}}
\end{verbatim}

\subsection{eval : combines arguments into a single string and runs }
\begin{verbatim}
      https://www.geeksforgeeks.org/eval-command-in-
      linux-with-examples/
      
      https://www.educative.io/answers/what-is-the-
      eval-command-in-linux
      
      eval runs string as a bash command
      
      eval is a built-in Linux command which is used to execute 
      arguments as a shell command. 
      It combines arguments into a single string and uses it as an 
      input to the shell and execute the commands. 
      
      When you have Linux command saved in a variable and wish to run 
      that command, the eval command is useful. 
      
      Example:
      $ COMM="ls"
      $ eval $COMM
      
      Example:
      $ eval echo hello
      
\end{verbatim}

\subsection{jobs   : Execute arguments as a single command in shell }
\begin{verbatim}
     Display status of jobs in background in the current session     
     
     The jobs command marks the last paused job with the + sign and the immediate previous stopped job   
     with the – sign. If we use fg and other job commands without a job number, the last paused job is  
     implied.
      - Show status of all jobs:
        jobs
      
      - Show status of a particular job:
        jobs %{{job_id}}

      - Show status and process IDs of all jobs:
        jobs -l

      - Show process IDs only of all jobs:
        jobs -p
        
      Example : 
      $ man ed
        ^Z  # ctrl+Z by user
        [1]+  Stopped                 man ed
      $ man cp
        ^Z  # ctrl+Z by user
        [2]+  Stopped                 man cp
      $ nano 2.txt
        ^T^Z # ctrl T + ctrl+Z
        [3]+  Stopped                 nano 2.txt   
      $ man pwd
        ^Z
        [4]+  Stopped                 man pwd
      $ jobs
        [1]   Stopped                 man ed      # "stopped" means paused !!!
        [2]   Stopped                 man cp
        [3]-  Stopped                 nano 2.txt  # job before the last 
        [4]+  Stopped                 man pwd     # last job
      $ jobs -l    # including job number
        [1]   6065 Stopped                 man ed
        [2]   6084 Stopped                 man cp
        [3]-  6120 Stopped (signal)        nano 2.txt
        [4]+  6123 Stopped                 man pwd
      $ fg %1
        man ed
\end{verbatim}


\subsection{kill   : send a signal to a process  }
\begin{verbatim}
       NAME
           kill - send a signal to a process

       SYNOPSIS
            kill [options] <pid> [...]

       DESCRIPTION
           The  default  signal  for kill is TERM.  
           Use -l or -L to list available signals.
           Particularly useful signals include HUP, INT, KILL, STOP,
           CONT, and 0.   
           
           Alternate signals may be specified in three ways: -9, -SIGKILL or -KILL.  
           
           Negative PID values may be  used to choose whole process groups; 
           see the PGID column in ps command output.  
           
           A PID of -1 is special; 
           it indicates all processes except the kill process itself and init.

       kill sends a signal to a process, usually related
       to stopping the process.
       All signals except for SIGKILL/KILL and SIGSTOP/STOP can
       be intercepted by the process to perform a clean exit.

       To present signals list (to be used without the SIG prefix, 
       first 15 signals are mandatory): 
           $ kill -l
    
    1) SIGHUP        2) SIGINT (1,^C)  3) SIGQUIT(^\)  4) SIGILL       5) SIGTRAP
                     ----------------  
    6) SIGABRT       7) SIGBUS         8) SIGFPE       9) SIGKILL(3)  10) SIGUSR1
                                                       -------------    
    11) SIGSEGV      12) SIGUSR2       13) SIGPIP      14) SIGALRM     15) SIGTERM (2)
                                                                       ---------------
    16) SIGSTKFLT    17) SIGCHLD       18) SIGCONT     19) SIGSTOP     20) SIGTSTP(^Z)

    
    Note1: ^D is not signal, it is simply End Of File.
           ^D closes Bash
    Note2: To respond on kill signal the process need to be run, 
           if the process receive the ^C signal while it is suspended with ^Z ,
           it will finish its work only after resume from suspend with fg command.    
    
                   SIGINT(^C) - The gentle nudge
                   -----------------------------
    SIGINT(^C), or <SIgnal INTerrupt>, is perhaps the most commonly
    encountered signal for many users. This signal is typically associated with the    
    CTRL+C command you often use in your terminal to stop a running process. 
    The primary purpose of SIGINT is to notify a process that the user 
    has requested an interruption.
    Example of SIGINT:
       $ kill -INT <process_id> # SIGINT (interrupt) signal
         or       
       $ kill -2 <process_id>
       
       Example :
          run on first terminal window:
           $ sleep 100  
          run on second terminal window:       
           $ ps aux | grep sleep
     root   9416  0.0  0.0  3212  1792 ?     S  Feb26 0:00 sleep 3600
     amits  0816  0.0  0.0  8376  1920 pts/0 S+ 00:16 0:00 sleep 100
     amits 10822  0.0  0.0  9216  2432 pts/1 S+ 00:16 0:00 grep --color=auto sleep
           $ kill -2 10816  # this will kill the process 
                            # 10816 on first terminal window.
                   
                   SIGTERM - The polite request
                   ----------------------------
      see: https://linuxhandbook.com/sigterm-vs-sigkill/               
      The SIGTERM (SIGnal and TERMinate) signal is used for terminating a program.
      SIGTERM is more forceful than SIGINT but still gives a process the opportunity
      to perform clean-up tasks before it ends. It allows the process to catch the signal
      and manage its termination elegantly – saving data or finishing essential tasks.
      The SIGTERM can also be referred as a soft kill because the process 
      that receives the SIGTERM signal may choose to ignore it. 
      SIGTERM doesn’t kill the child processes, With SIGTERM, a process gets
      the time to send the information to its parent and child processes. 
      Its child processes are handled by init.
      To send SIGTERM:
         $ kill <process_id>   # default SIGTERM (terminate)
           or   
         $ kill -s TERM <process id>   

                  SIGKILL - The last resort
                   ------------------------
      see: https://linuxhandbook.com/sigterm-vs-sigkill/               
      
      Now, what if a process does not respond to the SIGTERM signal, 
      or is stuck in an endless loop and not releasing the resources? 
      This is where SIGKILL comes in. SIGKILL, as the name suggests, 
      kills the process immediately. The system does not give the process 
      any chance to clean up or free up resources
      
      SIGKILL (the last option) is used for immediate termination of a process. 
      This signal cannot be ignored or blocked ! 
      
      The process will be terminated along with its threads (if any). 
      It is a brutal way of killing a process, and it should only be
      used as a last resort.
      Suppose you have an unresponsive process that you want to close. 
      The SIGKILL can be used in such a case. 
      SIGKILL kills the child processes as well. 
      Use of SIGKILL may lead to the creation of a zombie process because
      the killed process doesn’t get the chance to tell its parent process
      that it has received a kill signal.
        $ kill -9 <process id>
           or
        $ kill -KILL <process id>   

     Terminate a background job:
        $ kill %{{job_id}}
    
     see https://www.gnu.org/software/libc/manual/html_node/Standard-Signals.html
         https://www.gnu.org/software/libc/manual/html_node/Termination-Signals.html
         https://www.gnu.org/software/libc/manual/html_node/Job-Control-Signals.html
     
     INT  is for ^C
     QUIT is for ^/ , it is similar to ^C but performes the core damp.
     TERM is a generic signal used to program termination. Unlike 
          KILL, this signal can be blocked, handled, and ignored.     
     KILL is for immediate process termination. It cannot be handled or ignored,
          and is therefore always fatal. It is also not possible to block this signal.
     TSTP is ^Z signal The SIGTSTP signal is an interactive stop signal.
          Unlike STOP, this signal can be handled and ignored
     STOP is The STOP signal stops the process. It cannot be handled, 
          ignored, or blocked.
      
\end{verbatim}

\subsection{pkill     : Signal process by name }
\begin{verbatim}
      Signal processes by process's name
       
      see https://www.gnu.org/software/libc/manual/html_node/Standard-Signals.html
      https://www.gnu.org/software/libc/manual/html_node/Termination-Signals.html
      https://www.gnu.org/software/libc/manual/html_node/Job-Control-Signals.html
     
      INT  is for ^C
      QUIT is for ^/ , it is similar to ^C but performes the core damp.
      TERM is a generic signal used to program termination. Unlike 
           KILL, this signal can be blocked, handled, and ignored.     
      KILL is for immediate process termination. It cannot be 
           handled or ignored, and is therefore always fatal.
           It is also not possible to block this signal.
      TSTP is ^Z signal The SIGTSTP signal is an interactive stop signal.
           Unlike STOP, this signal can be handled and ignored
      STOP is The STOP signal stops the process. It cannot be handled, 
           ignored, or blocked. 
       
      First option:
        $ pkill -QUIT ping
        $ pkill -INT ping   # ^C
      Second option:
        $ pkill -TERM ping
      The Last option:
        $ pkill -KILL ping
     To stop process:
        ^Z: $ pkill -TSTP xed 

      - Force kill matching processes (can't be blocked):
           pkill -9 "{{process_name}}"
      - Send KILL signal to processes which match:
           pkill -KILL "{{process_name}}"
      Example :
      # on first terminal : 
       $ sleep 500
      # on second terminal
       $ ps aux | grep sleep
   root       14276  0.0  0.0   3212  1664 ?        S    01:41   0:00 sleep 3600
   amits      18341  0.0  0.0   8376  1920 pts/0    S+   02:36   0:00 sleep 500      
       $ pkill -KILL sleep     
       $ pkill -KILL sleep
   pkill: killing pid 14276 failed: Operation not permitted # because of root !

\end{verbatim}

\subsection{xkill     : Kill a window interactively in a graphical session }
\begin{verbatim}
         - Display a cursor to kill a window when pressing 
           the left mouse button (press any other mouse button to cancel):
            $ xkill

\end{verbatim}


\subsection{bg   : resume suspended process to run it in the background.  }
\begin{verbatim}
        https://www.scaler.com/topics/how-to-run-process-in-background-linux/

      - Resume the most recently suspended job and run it in the background:
        $ bg

      - Resume a specific job (use jobs -l to get its ID) and run it in the background:
        $ bg %{{job_id}}

        How To Run a Process in Background in Linux?
        Running processes in the background is a common requirement in Linux systems. 
        Running a process in the background allows you : 
              1. to continue using the terminal 
              2. execute other commands while the process runs independently. 
        This can be particularly useful for : 
              1.  long-running tasks or 
              2.  when you want to execute multiple commands simultaneously.

        What is a Background Processes
        Let's first understand what a background process is.
        In Linux, a background process refers to a process that runs independently
        of the terminal. 
        When you execute a command, it typically runs in the foreground, 
        meaning it occupies the terminal until it completes. 
        On the other hand, running a process in the background allows you 
        to execute other commands while the process continues to run silently. 

                       How to Use the bg Command in Linux
                       ----------------------------------
        One way to run a process in the background is by using the bg command.
        The bg command is built-in to most Linux distributions and allows you to send
        a running process to the background. Following steps are involved to run the 
        process in background:
           1. Start a process in the foreground by executing a command. 
              For example, let's say we want to compress a large file using the gzip command:
                $ gzip largefile.txt
           2. While the process is running, press Ctrl+Z to pause the process 
              and bring it to the background.
           3. Use the bg command to resume the process in the background:
                $ bg
           4. The process will now continue running in the background, 
              and you will see its job ID displayed on the terminal. 
              You can execute other commands while it progresses silently. 
              However, keep in mind that the process may still produce output, 
              which can appear on the terminal.

\end{verbatim}

\subsection{\& symbol   : start process in the background immediately  }
\begin{verbatim}
     & Symbol
     Another method to run a process in the background is by appending the
     ampersand symbol & to the command you execute. 
     This symbol tells the shell to start the process
     in the background immediately. 
     To just start the command in the background at the beginning: 
     all you need to do is put an ampersand at the end of any Linux command.
     This will allow to use terminal along with running the command:

     Example:

       $ gzip largefile.txt &

     In this example, the gzip command starts running in the background 
     as soon as you execute it. 
     The process ID (PID) will be displayed on the terminal, 
     allowing you to continue using the shell while the process completes silently.

     It's important to note that when using the & symbol, 
     the process may still produce output, which can interfere with
     the prompt on the terminal. 
     
     Note: 
     -----
     To avoid this, it's recommended to redirect the output
     to a file or to /dev/null, which discards the output. 
     You can do this by appending > filename or > /dev/null
     to the command, respectively.
\end{verbatim}

\subsection{disown     : Signal process by process's name }
\begin{verbatim}
     Allow sub-processes to live beyond the shell that they are attached to.
     - Disown the current job:
        $ disown

     - Disown a specific job:
        $ disown %{{job_number}}

     - Disown all jobs:
        $ disown -a
        
     If you close current terminal, the jobs that it run will be also closed,
     because the jobs are terminal’s children.
     Terminal itself is the child of graphical system (that presents windows).
     If you need to close your terminal, and don’t want these commands to stop,
     then you need to “disown” the job(s). 

     Disown command will cause the job to be child of graphical system,
     in place of child of terminal. 

     If you only have one job in the background, the following command will work:
        $ disown
     If you have multiple, then you’ll need to specify the job ID.
        $ disown %1
     You’ll no longer see the job in your jobs table when you execute the jobs command.
     Now it’s safe to close the terminal and your command will continue running.
        $ jobs -l
     You can still keep an eye on your running command by using the ps command.
        $ ps aux | grep sleep
   linuxco+    1650  0.0  0.0   8084   524 pts/0    S    12:27   0:00 sleep 10000
     And if you want to stop the command from running, you can use the kill
     command and specify the process ID.
        $ kill 1650        
     
     Examples : 
      $ geany | disown     
     
      $ geany & disown  # send to bg and disown 
     
\end{verbatim}

\subsection{time   : Run the command and print the time measurements  }
\begin{verbatim}

       Users of the bash shell need to use an explicit path in 
       order to run the  external time command "/bin/time" and not the 
       shell builtin variant "time" !!! 
          
          $ which -a time
             /usr/bin/time
             /bin/time
          $ time sleep 2
            real	0m2.003s
            user	0m0.002s
            sys	0m0.000s
          $ /bin/time --format="%E real, %U user, %S sys" sleep 2                                     
            0:02.00 real, 0.00 user, 0.00 sys     

       On system where
       time is installed in /usr/bin, the first example would become
       /usr/bin/time wc /etc/hosts        
       
       "time" runs the program "COMMAND" with any given arguments "ARG"....
       When COMMAND finishes, time displays information about resources
       used by COMMAND (on the standard error output, by default).
       If COMMAND exits with non-zero status, time
       displays a warning message and the exit status.
   
       Run the command and print the time measurements to stdout:
            $ time <command> <args>
   
\end{verbatim}


\subsection{fg   : resume a specific suspended job by its job number.  }
\begin{verbatim}
     The fg command allows us to resume a specific suspended job by
     its job number. The job’s name is output when it starts:
     Resume last job:
       $ fg 
       Resume job number 1:
       $ fg %1
\end{verbatim}

\subsection{\$?   : result of last command execution }
\begin{verbatim}

      $ ls -la
        total 772
        drwxrwxr-x 2 amits amits   4096 Feb 28 12:12 .
        drwxrwxr-x 8 amits amits   4096 Feb 28 11:42 ..
        -rw-rw-r-- 1 amits amits 256176 Feb 28 12:11 try.pdf
        -rw-rw-r-- 1 amits amits 148037 Feb 28 11:41 try.tex
      $ echo $?
        0
     https://www.baeldung.com/linux/check-if-command-executed-
     successfully#:~:text=We%20can%20use%20a%20binary%20comparison%20operator%20eq,
     the%20pwd%20command%20displays%20the%20present%20working%20directory.
     
       We can use a binary comparison operator eq to check the
       outcome of the previous command, stored in the $?   
       variable:

         #!/bin/bash
         pwd 
         if [ $? -eq 0 ]; then
            echo "Command Executed Successfully"
         else
            echo "Command Failed"
         fi{verbatim}

\end{verbatim}


\section{Commands envistigation}



\subsection{type     : Display the type of command the shell will execute.}
\begin{verbatim}

    The ‘type’ is a built-in shell command used to identify the kind of
    command. For example it checks if it is shell built in or external 
    command. 
     
     type without any key : 
                  if NAME is an alias,
    		        shell reserved word, 
    		        shell function, 
    		        shell builtin, 
    		        disk file
    		        or not found, respectively
     	
     -t	output a single word which is one of 
                  'alias',
                  'keyword',
                  'function', 
                  'builtin',
                  'file'

     -p	Returns the disk file that would be executed
     -a	Returns all locations containing an executable.

      For example:
      $ type echo
          echo is a shell builtin
      $ type cat
          cat is /usr/bin/cat
      $ type code
          code is /usr/bin/code     
      $ type -a time
          time is a shell keyword
          time is /usr/bin/time
          time is /bin/time

\end{verbatim}

\subsection{which     : Locate a program in the user's path.}
\begin{verbatim}
       Locate a program in the user's path.
       - Search the PATH environment variable and display the location 
         of any matching executables:
           $ which {{executable}}
       - If there are multiple executables which match, display all:
           $ which -a {{executable}}
        Example : 
           $ which -a time
              /usr/bin/time
              /bin/time
\end{verbatim}


\subsection{whereis  : locate the source, and manual page for a command.}
\begin{verbatim}
      $ type whereis
         whereis is /usr/bin/whereis


\end{verbatim}

\subsection{command  : locate the source, and manual page for a command.}
\begin{verbatim}



\end{verbatim}


\subsection{history  : command-line history}
\begin{verbatim}
      history   : Display the history list with line numbers. 
      history -c: delete all of the entries  
      history N : Only list the last N commands.
      history 10: list last 10 entries 
      ------------------------------------------------------- 
      !!        : to execute the last command you typed.
      !n        : to execute command number n.
      !-n       : execute the command n times before
      !!==!-1   :equivalent commands
      !"string" : execute the last command starting with
                  that “string” (last command with given name)
      -------------------------------------------------------            
      CTRL-R    : to perform a “reverse-i-search”. For 
                  example, if you wanted to use the command 
                  you used the last time you used snort,
                  you would type: CTRL-R then type “snort”.
			
      An argument of N lists only the last N entries.
		
\end{verbatim}



\subsection{man     : display manual pages}
\begin{verbatim}
     The table below shows the section numbers of the manual
     followed by the types of pages they contain.

        1	General commands
        2	System calls
        3	Library functions, covering in particular 
             the C standard library
        4	Special files (usually devices, those found in /dev) and drivers
        5	File formats and conventions
        6	Games and screensavers
        7	Miscellaneous
        8	System administration commands and daemons       
     
      When you invoke the man command, you can optionally
      specify the section of the manual to read :  
         $ man 3 printf
     
      To see all the man pages available in Ubuntu use 
         $ man -k .
     
      -f : to find page that contains the command :
         $ man -f signal
            signal (7)           - overview of signals
            signal (2)           - ANSI C signal handling
         $ man -f print
            print (1)            - execute programs via entries
                                   in the mailcap file     

      -k : to find command and its page that contains the word:
         $ man -k pgrep
            pgrep (1)            - look up, signal, or wait for processes   
                                   based on name a...
         $ zipgrep (1)           - search files in a ZIP archive for lines 
                                   matching a pat...

      man is an interface to the system reference manuals
        $ man -k : the same as apropos
        $ man -f : the same at whatis
        $ -h  - help on intercative commands
     
      find words in man manual:
       to find "name" option in man manual choose : 
           / "name"
       to find next search result press : 
           press "/" and press “enter”
\end{verbatim}


\subsection{tldr     : "Too Long; Didn't Read" - simplified man pages}
\begin{verbatim}
       simplified man pages
     
       To update : 
         $ tldr –update

       To install:
         $ sudo apt install tldr
         
         
       Update in Ubuntu (there is a problem to update tldr,
       since tldr -u did not work ): 
         
          $ tldr --version
            0.6.4 # this is the old version
          $ apt remove tldr
          $ pip install tldr
          $ tldr --version
            tldr 3.2.0 (Client Specification 1.5)

             
          see: https://forums.linuxmint.com/viewtopic.php?t=407665
          
          
            
\end{verbatim}


\subsection{whatis   : breaf command description}
\begin{verbatim}
      whatis is a command for obtaining the brief description
      of a specific command whose exact name is already known
\end{verbatim}


\subsection{apropos  : command to search the man page files}
\begin{verbatim}
       Often a wrapper for the man -k command.  
       Used to search the "name" sections of all manual pages
       for the specified string or strings (called keywords).
       The output is a list of all manual pages containing
       the search term (case insensitive) in their name or description.
       This is often useful if one knows the action that is desired,
       but does not remember the exact command or page name.
\end{verbatim}


\section{Text and text files manupulations}


\subsection{pandoc     : Convert documents between various formats}
\begin{verbatim}
       Convert documents between various formats:
       - List all supported input formats:
          $ pandoc --list-input-formats
       - List all supported output formats:
          $ pandoc --list-output-formats
       - Convert file to PDF (the output format is determined by file extension):
          $ pandoc {{input.md}} -o {{output.pdf}}
         
       Examples : 
         $ pandoc latex-document-prototype.tex -o try.html
         $ pandoc latex-document-prototype.tex -o try.pdf
         $ pandoc latex-document-prototype.tex -o try.docx      
\end{verbatim}


\subsection{tab     : autocomplete}
\begin{verbatim}
     Tab 		: automatic completion
     Tab Tab 	: all options for complete
\end{verbatim}


\subsection{ |     : pipe - redirect stdout of one cmd to stdin of another cmd.}
\begin{verbatim}


    • “Pipe” redirects standard output of one command to standard 
       input of another command
    • t1, t2, tL are “temporary files” in this explanation.
    • “|” acts as temporary file.
    • Pipe redirects only REGULAR input/output, not error


    keyboard --> command1 ---> | ---> command2 ---> |--> command3---> display
    command1 > t1
    command2 <t1 >t2
    command3 > tL
    
    • Example 1: 
      Redirect output of (print of all files *.log)
      to (grep -i (insensitive) that search lines with “error”)
      and sort all this lines:   
        $ cat *log|grep -i error|sort  
        


    • Example 2:
      Redirect (result of recursive insensitive search ,
      starting with current directory) to (search of lines that not    
      include word “ignored”), the (sort the results with “unic”)
      and send the result to file serious_errors.log:  
        $ grep -ri error .|grep -v “ignored”|sort u> serious_errors.log
        
    • Example 3 (YES command usage):          
      $ cat dir1_rm_inputs.txt
        y
        y
        n
        n
        y
        y
        y
      $ cat dir1_rm_inputs.txt | rm ./dir1/*.* -i

      $ rm -i dir2/*
         rm: remove regular empty file &apos;dir2/6.txt? n
         rm: remove regular empty file &apos;dir2/7.txt? n
         rm: remove regular file &apos;dir2/errors2.log? n
         rm: remove regular file &apos;dir2/errors3.log? n
         rm: remove regular file &apos;dir2/errors.log? n
      $ yes | rm -i dir2/*
         rm: remove regular empty file ‘dir2/6.txt’? rm: remove regular                             
         empty file ‘dir2/7.txt’? rm: remove regular file ‘dir2/
         errors2.log’? rm: remove regular file ‘dir2/errors3.log’? rm: 
         remove regular file ‘dir2/errors.log’?
         
    • Example: automation to enter “yes”  when install: 

      $ yes | sudo apt install some_programs

    • apt has -flag to allow install without interrupts.          
         
         
\end{verbatim}



\subsection{Redirections of I/O }
\begin{verbatim}
    

    Process:    
    
    We can control input and output direction of process:    
      0 ---> stdin
      1 ---> stdout
      2 ---> err    
    • Standard output can be redirected to file using > 
      (remove previous file and creates from the start of file).
    • Standard output can be appended to file using >> 
      (appends to the end of existing file).
    ------------------------------
    Redirect input to file:  
        $ cat operations.txt
          2+3
          4*5
          x=7
          x*100

        $ bc < operations.txt
          5
          20
          700    
    ----------------------------------
    Redirect output to file : 
      
    Redirect result of ls command to file ls_comm_result.txt:
      $ ls -l 1.txt 7.txt > ls_comm_result.txt
         ls: cannot access '7.txt': No such file or directory
    
    Redirect regular output of ls command and also ERR output 
    to the same file ls_comm_result.txt as regular output: 
      $ ls -l 1.txt 7.txt > ls_comm_result.txt 2>&1
         ls: cannot access '7.txt': No such file or directory      

        # or:
      $ ls -l 1.txt 7.txt 1>ls_comm_result.txt 2>&1
         ls: cannot access '7.txt': No such file or directory      

    Redirect ERR output to ls_comm_result.txt and 
    regular output to the same file as error output:  
      $ ls -l 1.txt 7.txt 2>ls_comm_result.txt 1>&2
        # or:
      $ ls -l 1.txt 7.txt 2>ls_comm_result.txt  >&2

    In This case regular output is redirected to errors output: 
      $ echo "error" >&2 
    
    Send grep command output to file :
      $ grep error errors.log > grep_result.txt
    
    Append grep command output to the end of the file: 
      $ grep error errors.log >> grep_result.txt
     
    Append cat command output to file: 
      $ cat 1.txt
        this is very interesting.
        what ?
        this.
      $ cat 2.txt
        also this is very nice and interesting
        exectly
      $ cat 2.txt >> 1.txt
      $ cat 1.txt
        this is very interesting.
        what ?
        this.
        also this is very nice and interesting
        exectly     
    
    Print cat command output to file: 
      $ cat 1.txt > 2.txt
      $ cat 1.txt
        this is very interesting.
        what ?
        this.
        also this is very nice and interesting
        exectly
      $ cat 2.txt
        this is very interesting.
        what ?
        this.
        also this is very nice and interesting
        exectly 
     
    Print ls command output to file:  
      $ ls
        1.txt  bc    dir2        fruits.txt       numbers.txt  sizes.txt
        2.txt  dir1  errors.log  grep_result.txt  out1
      $ ls > out
      $ cat out
        1.txt
        2.txt
        bc
        dir1
        dir2
        errors.log
        fruits.txt
        grep_result.txt
        numbers.txt
        out
        out1
        sizes.txt
 
    Use echo command output to prepare file in place of nano: 
      $ echo 'this
      > is a file
      > that try to use echo command
      > in place of nano command'> my_file.txt
      $ cat my_file.txt 
        this
        is a file
        that try to use echo command
        in place of nano command


    Append line to file with echo command output: 
      $ echo "ups fogot this line" >> my_file.txt
      $ cat my_file.txt 
        this
        is a file
        that try to use echo command
        in place of cat command
        ups fogot this line
     
    Redirect errors output to additional file errors.txt:    
      $ ls 5.txt z.txt
	    ls: cannot access ‘z.txt’ No such file or directory # err output 
	    5.txt                                               # regular output
      $ ls 5.txt z.txt >out.txt 2>error.txt # same as 1>out.txt
      $ cat out.txt
	    5.txt
      $ cat error.txt
	    ls: cannot access ‘z.txt’: No such file or directory  
        
        
\end{verbatim}



\subsection{echo     : display line of text}
\begin{verbatim}
     
     Display line of text/string that are passed as an argument.
     Mostly used in shell scripts and batch files 
     to output status text to the screen or a file
     Helps to see a coomand after expansion
 
       $ echo hello
       hello

       $ echo hello   world   # 2 different arguments : Hello and World
       hello world

       $ echo "hello   world" # 1 argument "Hello world"
       hello   world

       $ echo ~
       /home/amit

       $ echo "~"
       ~

       echo ls *.txt # to see the command after expansion

       ? replaces exactly one character:
       $ echo ?.txt # ? replace exactly one character
          1.txt 3.txt 5.txt 7.txt               

       The list of files with extention of 4 characters :
       $ echo *.????
          2.html 5.html

       $ echo ??.txt
          11.txt 12.txt
 
       $ echo ???.txt
          111.txt

       $ echo 1?.txt
          11.txt 12.txt

       $ echo 1??.txt
          111.txt

       $ echo [0-9][0-9][0-9].txt
          111.txt
       
       ^ is for "not"  : 
       $ echo [^0-9].txt  # not(numbers 0 to 9)
          a.txt b.txt g.txt i.txt

       $ echo [a-z].txt
          a.txt b.txt g.txt i.txt

       $ echo [1a2x3456bc].txt
          1.txt 3.txt 5.txt a.txt b.txt

       $ echo {1..9}
          1 2 3 4 5 6 7 8 9

       $ echo {1..100..5}  # from 1 to 100 with step 5
          1 6 11 16 21 26 31 36 41 46 51 56 61 66 71 76 81 86 91 96

       $ echo a{,x,y,z}4  # {,x,y,x} means : nothing,x,y,z
          a4 ax4 ay4 az4

       $ echo a{“”,x,y,z}4
          a4 ax4 ay4 az4

       $ echo {,.}*.txt # with or without "."
          *.txt .hidden.txt .secret.txt
       
       $ echo *.txt .*.txt
          *.txt .hidden.txt .secret.txt       
       
       $ echo */*.txt
            docs/111.txt docs/11.txt docs/12.txt docs/1.txt
            docs/3.txt docs/5.txt docs/7.txt docs/a.txt
            docs/b.txt docs/g.txt docs/i.txt
       $ echo "a{b,c,d}e"
            a{b,c,d}e
       $ echo a{b,c,d}e
            abe ace ade     

         to enable ** recursive glob :
            $ shopt -s globstar  // set recursive glob on
            $ shopt globstar
               globstar       	on
            $ echo **/*.txt   # recursive glob: 
                              # find in the current dir (!)
                              # and all sub dirs under it 
              1.txt 2.txt 3.txt 4.txt 5.txt   # files from current dir also
              temp1/a.txt temp1/b.txt temp1/c.txt temp1/d.txt
              temp2/10.txt temp2/1.txt temp2/2.txt temp2/3.txt
              temp2/4.txt temp2/5.txt temp2/6.txt temp2/7.txt
              temp2/8.txt temp2/9.txt
            $ shopt -u globstar
            $ shopt globstar
               globstar       	off
            $ echo **/*.*    # now it will search only in sub dirs
              temp1/a.txt temp1/b.txt temp1/c.txt temp1/d.txt 
              temp2/10.txt temp2/1.txt temp2/2.txt temp2/3.txt 
              temp2/4.txt temp2/5.txt temp2/6.txt temp2/7.txt 
              temp2/8.txt temp2/9.txt
       
           search in subdirs only : 
           	$ echo */*.txt
              temp1/a.txt temp1/b.txt temp1/c.txt temp1/d.txt 
              temp2/10.txt temp2/1.txt temp2/2.txt temp2/3.txt 
              temp2/4.txt temp2/5.txt temp2/6.txt temp2/7.txt 
              temp2/8.txt temp2/9.txt  
           
           split text in several strings:   
             $ echo 'this
             > is afile
             > that is vary long'
             this
             is afile
             that is vary long

           Trick to print ouput to stderr from process:
           ---------------
           ex07.sh:
             echo "normal output"
             echo "error output" >&2            
           --------------- 
             $ bash ex07*  >/dev/null
             $ bash ex07* 2>/dev/null 
               
                      
\end{verbatim}

\subsection{printf     : display line of text}
\begin{verbatim}
     Format and print text.

    Similar to echo command, but does not output 
    new line, if not defined particularly.
      $ printf “\rA”   – return to start of the same line and 
                         print A.
      $ printf “A\n” – print A and add enter.


      $ printf "Hello World\n"
        Hello World
      $ printf "Hello World"
        Hello World $
       
      This will output A, wait 3 sec., clear line and output B:
      $ printf "A"; sleep 3; printf "\rB\n"
      
           
\end{verbatim}


\subsection{grep     : "Global Regular Expression search and Print" }
\begin{verbatim}
     https://en.wikipedia.org/wiki/Grep
     
     print lines that match patterns
     grep allows to enter regular expression for the search.
     
     returns status  $? = 0 if string is found,
     returns status  $? = 1 if string is not found
     returns status  $? = 2 if error occured
     
     
     grep is a command-line utility for searching 
     plain-text data sets for lines that match a regular expression
     i.e. to search files for certain patterns
     Filter only lines with desired contents

     - Search for a pattern within a file:
          grep "{{search_pattern}}" {{path/to/file}}
          
          $ grep "is" *.txt
				1.txt:this is an adittional try
				try1.txt:This is very convinient editor, 
				try1.txt:Its name is tilde.    

          $ grep is *.txt
				1.txt:this is an adittional try
				try1.txt:This is very convinient editor, 
				try1.txt:Its name is tilde. 
     ----------------------------------------------------
     Options:
         -L : list of FILES that don’t include the word
         -l : list of FILES that include the result
         -i : Insesitive to upper case     
         -w : Select  only  lines with whole word.
         -o : present ONLY WORD we search
         -r : recursive search
         -v : (inVerterd) list of rows that does not include the word
         -c : Count number of lines with word
         -q : quiet, exit immediately with 0 status if any match is found.
     ----------------------------------------------------
     
     - Search for an exact string (disables regular expressions):
          grep --fixed-strings "{{exact_string}}" {{path/to/file}}

     - Search for a pattern in all files recursively in a directory, showing   
       line numbers of matches, ignoring binary files:
          grep --recursive --line-number --binary-files={{without-match}}   
          "{{search_pattern}}" {{path/to/directory}}
         
         Global Regular Expression "error" search:
             $ grep error errors.log
                error: aeasasadf
                dfdfdfdf error sfsadgasdg
                ldksldjks no errors

         -i : Insesitive to upper case 
            $ grep error errors.log -i
                error: aeasasadf
                dfdfdfdf error sfsadgasdg
                ERROR dssdsdsdsd
                ldksldjks no errors

         -w : exact word
            $ grep errors errors.log -i -w
               ldksldjks no errors

         -o : present only word we search, without text 
            $ grep error errors.log -o -i
               error
               error
               ERROR
               error
         
         -r : recursive search in all subdirectories : 
            $ grep error -r
                dir2/errors2.log:error: aeasasadf
                dir2/errors2.log:dfdfdfdf error sfsadgasdg
                errors.log:dfdfdfdf error sfsadgasdg
                errors.log:ldksldjks no errors

         -l : list of files that include the result : 
            $ grep error -r -l
                 dir2/errors2.log
                 dir2/errors.log
                 dir2/errors3.log
                 errors.log

         -L : list of files that don’t include the word : 
            $ grep error -r -L
                .secret.txt
                .hidden.txt
                dir1/5.txt

         -v : (inVerterd) list of rows that does not include the word "error": 
            $ grep error errors.log -v
                aaaaa
                bbbb

                dfdf

                ERROR dssdsdsdsd
                sddsd
                dsdsds
             
         -v -i : (inVerterd) list of rows that does not include the word "error"
                  insensitive : 
            $ grep error errors.log -v -i
              aaaaa
              bbbb

              dfdf

              sddsd
              dsdsds 

         List all files, that contain at least one row without word error:                   
            $ grep error -v -r -l
              dir1/b.txt
              dir1/dudu.txt
              dir2/errors2.log
      
         
         List all files that does not contain any word “error”: 
            $ grep error -r -L
              .secret.txt
              .hidden.txt
              dir1/5.txt
              dir1/b.txt
              
         Count number of lines with word "error":      
             $ grep error errors.log -c
                 3
             $ grep error errors.log -i -c
                 4   
                
                
                       
\end{verbatim}


\subsection{cat      : "ConcATenite" : concatenite and print.}
\begin{verbatim}
     The cat utility serves a dual purpose: concatenating and printing.
     
     With a single argument, it is often used to print a file to the 
     user's terminal emulator.
     
     With more than one argument, it concatenates several files

     - Print the contents of a file to stdout:
           cat {{path/to/file}}

     - Concatenate several files into an output file:
           cat {{path/to/file1 path/to/file2 ...}} > {{path/to/output_file}}

     - Append several files to an output file:
           cat {{path/to/file1 path/to/file2 ...}} >> {{path/to/output_file}}

    

\end{verbatim}

\subsection{more     : present text file in parts. }
\begin{verbatim}

       Same as cat , but presents the contents
       in several parts with more option
       
       $ cat /usr/include/stdio.h 
       $ more /usr/include/stdio.h 

\end{verbatim}

\subsection{less     : present text file in parts and allows search }
\begin{verbatim}
       
       Same as cat , presents the contents
       in option to move with arrows
       allows search
       
       $ less /usr/include/stdio.h

       allows to find with
       / prin....
       /+”enter”

\end{verbatim}

\subsection{head     : Display n rows at the head of file }
\begin{verbatim}
       Present n rows at the head of fil
       
       
       Present 3 rows at the head:
             $ head /usr/include/stdio.h -n 3

\end{verbatim}

\subsection{tail     : Display n rows at the end of file }
\begin{verbatim}
         Present n rows at the end of file
        
        - Show last 'count' lines in file:
                tail --lines {{count}} {{path/to/file}}

             $ tail --lines 3 /usr/include/stdio.h

\end{verbatim}

\subsection{sort     : Display lines in sorted order}
\begin{verbatim}
           prints the lines of its input or concatenation 
           of all files listed in its argument list in sorted order
        regular sort:
           $ sort fruits.txt
             apple
             apple
             banana
             banana
             pineapple
             strawberry
             tomato
             watermelon

        -u: unic sort remove the same words: 
          $ sort fruits.txt -u
            apple
            banana
            pineapple
            strawberry
            tomato
            watermelon


        -r : reverse sort :
          $ sort fruits.txt -u -r
            watermelon
            tomato
            strawberry
            pineapple
            banana
            apple

        -h: sort file sizes in human order :
          $ sort files sizes.txt -h
            10
            30K
            200K
            500K
            5M
            1G
            2G
    • sort numbers: 
         $ sort numbers.txt
            1
            1
            1212
            19898990
            2
            33
            45
\end{verbatim}

\subsection{wc     : "Word Count" - standard input statistics}
\begin{verbatim}
     Count lines, words, and bytes.     
     The program reads either standard input or a list of computer
     files and generates one or more of the following statistics: 
          1. newline count,
          2. word count,
          3. byte count. 
          4. charachters count
     If a list of files is provided, both individual file and total
     statistics follow.
       - Count all lines in a file: 
           wc --lines {{path/to/file}}
      - Count all words in a file:
           wc --words {{path/to/file}}
      - Count all bytes in a file:
           wc --bytes {{path/to/file}}
      - Count all characters in a file (taking multi-byte characters into account):
           wc --chars {{path/to/file}}
      - Count all lines, words and bytes from stdin: 
          {{find .}} | wc
      - Count the length of the longest line in number of characters:
          wc --max-line-length {{path/to/file}}
          
       Example: 
       $ wc ex1.c
          8  20 138 ex1.c

          #8   lines
          #20  words,
          #138 characters   
          
          
\end{verbatim}


\section{Compressing and archiving}

\subsection{gzip/gunzip  : shrink huge files and save space}
\begin{verbatim}

        GNU zip compression utility. 
        Creates .gz files (original file is deleted)
        Ordinary performance (similar to Zip).
        Very popular in Linux, almost every Linux has gzip.
        Relative fast and well decompression  
        Has no password 
        
        NOTE: Does not work with directories !
         
        $ gzip <file>    # to compress  
        $ gunzip <file>  # to decompress
        
        
        Example : 
        Suppose in current directory there are 3 files:
        1.txt , 2.txt , 3.txt
        
        $ gzip 1.txt 2.txt 3.txt
        Normally this will replace the original file by compressed file : 
          1.txt.gz in place of 1.txt
          2.txt.gz in place of 2.txt
          3.txt.gz in place of 3.txt         
        
        Note: if the file has a hard link, the compression will not be 
        done by gzip: 
         $ gzip 1.txt 2.txt 3.txt
           zip: 1.txt has 1 other link  -- unchanged
         $ file 2.txt
         $ file 2.txt.gz
            2.txt.gz: gzip compressed data, was "2.txt", 
            last modified: Thu Feb 15 11:40:21 2024, from Unix, truncated
         $ gunzip 2.txt # produce 2.txt
         $ gunzip 3.txt # produce 3.txt
   
         -k : keep input files during compression (don't replace)
         $ gzip -k 2.txt # now there are 2 files : 2.txt and 2.txt.gz
         
         
\end{verbatim}

\subsection{bzip2/bunzip2  : better compression than gzip}
\begin{verbatim}
        Similar usage as gzip.         
        Compression level is higher than in in gzip
        Decompression sometimes problematic (may be).
        Does not work with directories
        
        
        $ bzip2 <file>    # to compress
        $ bunzip2 <file>  # to decompress
        
        Example: 
        $ bzip2 3.txt        # replace 3.txt with 3.txt.bz2
        
        $ bunzip2 3.txt.bz2  # replace 3.txt.bz2 with 3.txt
        
        $ bzip2 -k  3.txt    # generates 3.txt.bz2 
                             # and dont change 3.txt 
        
\end{verbatim}


\subsection{tar  : Tape ARchive - back up set of files}
\begin{verbatim}
      
      Archiving utility (without compressing !!!).
      Often combined with a compression method, such as gzip or bzip2
      
      So for extraction of information within *.tar.gz we perform 
      following steps:       
      *.tar.gz ---> gzip decompress ---> *.tar ----> tar extract

      Note : jar in Java has the same flags as tar in linux      
      
      Useful to backup or release a set of files into a single file
      In the past used to move files from Linux to tape.
      tar stores also structure of directories, inods , etc. and send 
      it to tape archive.
      
      Note : tar stores the file system along with inode data
      so it stores the original permissions of the files !!!
      
      Creating an archive:
      --------------------
        $ tar -c  -f <archive_file_name>  <files or directories>
      Extracting archive:
      -------------------      
              -x  
      Listing archive:
      ----------------      
              -t
          -c: create. The opposite action is extract (-x)
          -x: extract
          -t: test (listing) check that archive is OK
          In the past tasr used >> to sent to tape
          now we use -f to tell tar to use file:
          -f: file. Archive created in/extracted from file 
                   (tape used otherwise if no -f).
           Note, that if we don't add -f flag 
           tar will read from stdin and write the result to stdout         
          --------------------------------
          Also:
          -v: verbose. Useful to follow archiving progress.
          --------------------------------
          Creating an archive:
             $ tar -cvf <archive> <files or directories>
          Extracting an archive:
             $ tar -xvf <archive> <files or directories>       
          Viewing the contents of an archive or integrity check:
             $ tar -tvf <archive>
           
          Combination of compressing and archiving on the fly.
          Useful to avoid creating huge intermediate files
          Much simpler than with tar and bzip2\gzip!

          -j --bz2  option: [un]compresses on the fly with bzip2
          -z --gzip option: [un]compresses on the fly with gzip 
          -J --xz   option: [un]compresses on the fly with xz  
          
          create tar.gz file (- czf):
           $ tar -czfv files.tar.gz 1.txt 2.txt dir1 dir2 # create z file  
          unzip  tar.gz file:          
           $ tar -xzf files.tar.gz # eXtract z file
          list tar.gz file : 
           $ tar -tzf files.tar.gz # list z file        
          
          --------------------------------
          tar -czfv files.tar.gz 1.txt
          tar -cfv 1.txt | gzip > files.tar.gz  # the same operation in
                                                # another form
          --------------------------------                                   
          Files or directories are given with paths relative to 
          the archive’s directory.
          
          Use -C DIR or --directory=DIR, to specify the 
          extraction directory location.
          
          Suppose we have following directories tree:
          Desktop
          |          
          |--- linux_2312
                 |
                 |---- playground
                          |
                          |--- 1.txt
                          |--- 2.txt
                 
          And we want like to compress playground directory,
          including name playground:
            amits:Desktop$ tar -czf play.tar.gz -C linux_2312 -v playground/
               playground/
               playground/2.txt
               playground/1.txt
            amits:Desktop$ ls
               Documents     LinuxAdminLearn.sh     play.tar.gz
            amits:Desktop$ tar -tzf play.tar.gz
              playground/
              playground/2.txt
              playground/1.txt
            
            amits:Desktop$ mkdir try_dir 
            amits:Desktop$ mv play.tar.gz ./try_dir
            amits:Desktop$ cd try_dir
            amits:Desktop$ tar -xzf play.tar.gz
            amits:Desktop$ ls
              playground  play.tar.gz 
            amits:try_dir$ cd playground/
            amits:playground$ ls
              1.txt  2.txt
          --------------------------------
          Example : 
          create archive "text files.tar": 
            $ tar -c -f text_files.tar 1.txt 2.txt ./dir1 ./dir2
            
          to check what is inside the archive text_files.tar : 
            $ tar -t -f text_files.tar 
            or
            $ tar -t <text_files.tar            
                      
           Now if we want to extract only 1.txt from tar frile 
           text_files.tar : 
            $ tar -x -f text_files.tar 1.txt
            ------------------------------------------------  
              tar   -x    -f text_files.tar        1.txt
            command flag  flag with its argument   argument        
            ------------------------------------------------
            Now we run gzip on file text_files.tar:
            $ gzip text_files.tar # in place of text_files.tar
                                  # it place text_files.tar.gz
            $ gunzip text_files.tar
            ------------------------------------------------- 
            $ mkdir try_dir
            $ mv text_files.tar try_dir
            $ gzip text_files.tar  # run gzip on tar
            $ ls
              text_files.tar.gz    # this file stores inode and compressed
            $ gunzip text_files.tar.gz
            $ ls
              text_files.tar
            $ tar -xf text_files.tar
            $ ls
              1.txt  2.txt  dir1  dir2  text_files.tar
            ------------------------------------------  
            Example with xz : 
            $ tar --xz -cf data.xz 1.txt  
            $ tar --xz -xf data.xz 
            ------------------------------------------ 

\end{verbatim}


\subsection{xz/  : relative new, better compression}
\begin{verbatim}
        New application with huigh compression level.
   
\end{verbatim}


\subsection{7zip/  : better compression than bzip2}
\begin{verbatim}
    7zip supports strong AES256 encryption.
    No need to encrypt in a separate pass.
    At last a solution available for Unix and Windows!
    The tool supports most other compression formats:
    zip, cab, arj, gzip, bzip2, tar, cpio, rpm , deb.
     
    Option 1 to install : 
        
    https://www.linuxcapable.com/how-to-install-7-zip-on-ubuntu-linux/
    Download 7-Zip using the wget command below:

     $ wget https://www.7-zip.org/a/7z2301-linux-x64.tar.xz  
     $ mkdir 7zip
     $ tar -xJf 7z2301-linux-x64.tar.xz -C 7zip
     $ cd 7zip
     $ ./7zz
     
    Option 2 to install (preferred): 
        https://www.geeksforgeeks.org/how-to-install-7zip-in-ubuntu/
        
        $ sudo apt-get update
        $ sudo apt-get install p7zip-full
        

        Much better compression ratio than bzip2 (up to 10 to 20%).
        Actually its compression is the best of all options.
        Free of charge, but was developed by company
         $ man 7z      # 7zip manual 
        
        Add files to 7z : 
        
          $ 7z a data.7z 1.txt # add files to archive
           
           7-Zip [64] 16.02 : Copyright (c) 1999-2016 Igor Pavlov : 2016-05-21
           p7zip Version 16.02 (locale=en_IL,Utf16=on,HugeFiles=on,64 bits,12      
           CPUs 12th Gen Intel(R) Core(TM) i7-1255U (906A4),ASM,AES-NI)
           Scanning the drive:
           1 file, 474 bytes (1 KiB)
           Creating archive: data.7z
           Items to compress: 1
           Files read from disk: 1
           Archive size: 409 bytes (1 KiB)
           Everything is Ok
         
         List data in archive:
          $ 7z x py2403_day03.zip  # to see contents of zip file
          $ 7z l data.7z

                 Date      Time    Attr         Size   Compressed  Name
     ------------------- ----- ------------ ------------  -----------------
     2024-02-26 15:18:03 ....A          474          295  1.txt
     ------------------- ----- ------------ ------------  -----------------
     2024-02-26 15:18:03                474          295  1 files

         Decompress files: 
          
          $ 7z x data.7z 
        
         Unzip zip file:  
          $ 7z x py2403_day03.zip
          
\end{verbatim}

\section{File comparison and file integrity}

\subsection{file..integrity..check}
\begin{verbatim}

     We can look at file data as single lage number.
     We can define function that translates this large number
     to another number called "checksum"
     
     Can another file generate the same checksum ?
     crc32 number present 4.294.967.296 different numbers 
     there are more than 4e9 different files
     
     According to shovah ha yonim there are two different files with the 
     same crc number.

\end{verbatim}

\subsection{crc32  : calculate crc32}
\begin{verbatim}
   install: 
   $ sudo apt install libarchive-zip-perl
    
   crc32 - compute CRC-32 checksums for the given files
   crc32 filename [ filename ... ]  
   
   For example:
   $ cat 1.txt
     this is David's file
   $ crc32 1.txt
     7b6accd4
  
     
     
\end{verbatim}


\section{Find files and check directory size}

\subsection{du       : "Disk Usage" - file and dir usage space}
\begin{verbatim}
      
      https://phoenixnap.com/kb/show-linux-directory-size
      
      Display the current directory size by running du without any arguments:
      The system lists the current directory content and subdirectories with a    
      number representing the object size in kilobytes. 
      The current directory size is the last line.
       $ du
      Print the Current Directory Size in a Human-Friendly Format
      Add the -h option to the du command to make the output easily readable:
       $ du -h
      
      Find the Size of a Specific Directory
       $ du -h ./example1
      
      Show Only the Total Disk Usage of a Directory
      To print an extra total line at the end of the output and summarize the disk     
      space usage for all files and directories, use the -c argument:
       $ du -c ./example1
      
      Scan up to a Chosen Subdirectory Level
      Limit the scan to a certain subdirectory level with the --max-depth option.     
      For example, to scan only the top-level directory (example1), 
      use --max-depth=0:
       $ du -h --max-depth=0 ./example1
      To list the top directory (example1) and the first layer of subdirectories    
      (example2 and example3), use --max-depth=1:
       du -h --max-depth=1 ./example1
      --------------------------------------------------------------------------
      
      Disk usage: estimate and summarize file and directory space usage
       - List the sizes of a directory and any subdirectories, in human-   
         readable form (i.e. auto-selecting the appropriate unit for each
         size):
            $ du -h {{path/to/directory}}

       - List the human-readable sizes of a directory and of all the files and 
         directories within it:
            $ du -ah {{path/to/directory}}


        Summary of disk usage for two directories:
           $ du -h
             8.0K	./dir2/2024
             12K	./dir2
             8.0K	./dir1/2024
             20K	./dir1
             76K	.
           $ du -h | sort -h
             8.0K	./dir1/2024
             8.0K	./dir2/2024
             12K	./dir2
             20K	./dir1
             76K	.
           $ du -h -s ./dir1 ./dir2  # -s is for summary of size
             20K	./dir1
             12K	./dir2
             
         Apparent size is the actual amount of data that can be read from the 
         file. Block-oriented devices can only store in terms of blocks, not 
         bytes. As a result, the disk usage is always rounded up to the next 
         highest block     
           $ du -h -s --apparent-size ./dir1 ./dir2  
            
            
             
\end{verbatim}




\subsection{locate   :  command}
\begin{verbatim}


    Uses database fo find files.
    Note: the database is recomputed periodically (usually weekly or daily).
    Manual : the database is located in /var/lib/mlocate/mlocate.db
    If we add file to disk or delete file from disk we need to update the data 
     base. To update database:
        $ sudo updatedb
    - locate a file by its exact filename
       $ locate stdio.h | less
    - locate all text files in current directory:       
       $ locate "$PWD/*.txt"
    - locate file *.txt starting with root directory:
       $ locate "/*1.txt" | head



\end{verbatim}

\subsection{find   :  command}
\begin{verbatim}
       Find filenames quickly. 

       Find files or directories under the given directory 
       tree, recursively


        Suppose we need to find file stdio.h
        $ cd /
        $ find . -name "stdio.h" 2>/dev/null
        
        Suppose we need to find all *.txt files starting from ~
        $ cd ~
        $ find . -name "*.txt"
        $ cd ~/Desktop
        
        There are many flags in find command , 
        the flags are divided in categories :
	     Actions (-print by default)
	     Operators (-o logical OR, -a logical AND)
	     Tests (+n(>n) -n(<n)  n (=n)  ) check size, time , etc.
                For example: 
                   -mtime +5 : if file modified in 5 last day

         List all files that modified in 14 last days : 
         $ find -mtime -14

        
         -mtime n
              File's data was last modified less than, more  than  or  exactly
              n*24  hours  ago.
         
         List all files that size of them is larger that 2G :
         $ find -size +2G 2>/dev/null 
         
         List all files from ~ that are at least 100Mb :
         $ find ~ -size +100M -name "*cl*" 2>/dev/null
        
\end{verbatim}

\subsection{tree    : show curr. dir as tree.}
\begin{verbatim}
       Show current directory as tree , 
       with sizes in human style and permissions : 
       
       
       Combine tree with du : 
       ----------------------        
       $ tree --du -h -p       
       
       -L     Level
       -s     Print the size of each file in bytes along with the name.
       -h     Print the size of each file but in a more  human  readable  way,
              e.g.  appending  a size letter for kilobytes (K), megabytes (M),
              gigabytes (G), terabytes (T), petabytes (P) and exabytes (E).             
       
       --du   For each directory report its size as the accumulation of  sizes
              of  all  its  files and sub-directories
              

       --inodes
              Prints the inode number of the file or directory
       
       -p     Print  the  file  type  and permissions for each file (as per ls
              -l).
     
       
       -u     Print the username, or UID # if no username is available, of the
              file.
       -g     Print the group name, or GID # if no group name is available, of
              the file.
            
      - Print files and directories up to 'num' levels of 
        depth (where 1 means the current directory):
           $ tree -L {{num}}

      - Print directories only:
           $ tree -d

      - Print hidden files too with colorization on:
           $ tree -a -C
       
       
\end{verbatim}


\section{Dirs and files manupulation}

\subsection{directory.permissions.in.Linux}
\begin{verbatim}
  see https://unix.stackexchange.com/questions/21251/
  execute-vs-read-bit-how-do-directory-permissions-in-linux-work

        When applying permissions to directories on Linux, the permission bits 
        have different meanings than on regular files.


        The read bit (r) allows the affected user to list the files 
        within the directory

        The write bit (w) allows the affected user to create, rename, or 
        delete files within the directory, and modify the directory's 
        attributes

        The execute bit (x) allows the affected user to enter the directory,  
        and access files and directories inside
        
        The sticky bit (T, or t if the execute bit is set for others) states 
        that files and directories within that directory may only be deleted 
        or renamed by their owner (or root)


\end{verbatim}




\subsection{basename  : store the name of dir based on full path}
\begin{verbatim}
     Remove leading directory portions from a path.
      - Show only the file name from a path:
         $ basename {{path/to/file}}

      - Show only the rightmost directory name from a path:
         $ basename {{path/to/directory/}}

\end{verbatim}


\subsection{nautilus  : file explorer in Ubuntu}
\begin{verbatim}
     nautilus : native file explorer in Ubuntu
\end{verbatim}


\subsection{chmod  : CHange file MODe (permissions and spec. modes)}
\begin{verbatim}

    • At the begin of -rwxrwxrwx is the type of file 
    (soft link (l), dir(d), regular file(-)) 

      r - read rights (ex: present the contents with cat)
      w – write rights (ex: write with nano)
      x – execute rights (ex: run executable)

      u – user (file’s owner)
      g – group
      o – others

      Process knows who is its owner and what are groups that run it.
      When process wants TO USE FILE it can check the rights to use it. 
      inode of file contains info of file’s owner and what is the file’s
      group. Sometimes operational system prevents from use the files, and
      sometimes process checks the match between process’s owner and group 
      again file’s inode information. 
      If there is no match – the process stops its execution. 
      In case of files: all the above info relates only to the CONTENTS 
      of the file, not INODE and not FILE's NAME. 
      
      In case of file of type "directory", the content of file contains the  
      list of files in the folder that can be read with "ls" command.
      Read the file of type "directory" is to run ls command. 
      File of type "directory" contains the list of files and directories  
      within the directory.
	  -------------------------------------------------------------------	
	  FILE PROTECTION must be done at two (!) levels : 
	  1. Permission of the file itself.
	  2. Permissions of the folder that contains the file (to protect the 
	     info about the file).
	  Otherwise the PROTECTED file can be DELETED from the directory's list 
	  of files (by delete the directory) !!!
	  Because "rm" command only deletes the registration of the file within 
	  the directory. 
	  
	  Delete of the write protected file related to directory permissions. 
	  If directory of the file has write permissions and at the same time the 
	  file has no write permissions – it will be deleted !!!:    

         $ ls -l 1.txt
	        -rw-rw-r-- 1 amit amit 146 Sep 22 12:54 1.txt
         $ chmod u-w 1.txt
         $ ls -l 1.txt
	        -r--rw-r-- 1 amit amit 146 Sep 22 12:54 1.txt
         $ rm 1.txt
	          rm: remove write-protected regular file '1.txt'? Y
         $ ls -l 1.txt
	       ls: cannot access ‘1.txt’: No such file or directory

      Only write protection to file’s directory really protects the files in    
      this directory from being deleted. But it does not protect the files 
      from being edited !!! So both protection is needed : the file protect 
      and the directory protect. 
      So, protection of the file from being deleted must be done in two
      levels simulataneously: 
             1. Protection of file’s folder (write option - off). 
             2. Protection of the file by setting the file’s own
                protection option (write option off).
	  -------------------------------------------------------------------	
      CMOD command:
        $ chmod [symbolic permissions/octal premissions] filename
          
          Symbolic permissions : 
          Who:         Action:        Permission type:
          u - user     + (turn on )   r(read)
          g - group    - (turn off)   w(write)
          o - others   = (set)        x(execute)
          a - all

         Take combination of u,g,o,a ; one of +/-/= ; combination of r,w,x
          • u+r : turn on “r” of user
          • additional permissions separated by “,”:
               u+r,g-wx : set on r to user and set off wx to group
          • ug+x : set on x for u and g
          • unset all r : ugo-r ; a-r
          • u=  #u equals nothing
          • unset rxw for u : u-rxw
          • use of = : u=r
	      
          Octal permissions:
             $ chmod 777 1.txt
                U    G    O 
               rwx  rwx  rwx
                7    7    7
           	 
           	 $ chmod 746 1.txt          
                U    G    O 
               rwx  rw-  r--
                7    6    4
           ----------------------------------------------------------------------
                           -R : Apply chmod recursively.
           Note : uppercase X will be set only for folders or executable files.
                  always if you want to permit execution to folders and 
                  files - use X not x !!!   
           -R : change files and directories recursively
                $ chmod -R o-rwx ./  # Unset rwx permissions for o(thers) recursevely
                $ chmod -R o+rwX ./  # Set rwx permissons recuresevely 
                                     # Uppercase X : x will be set only for folders or  
                                     # executable files ! 
                $ tree -p            # present permissions also                     
           
           
           To make directory dir1 and everything in it available to everyone: 
           $ chmod -R a+rX dir1
           $ ls -liR dir1 
           
           ----------------------------------------------------------------------
           
           See https://linuxhandbook.com/suid-sgid-sticky-bit/
           In addition of 9 options {u(r,w,x), g(r,w,x), o(r,w,x)} there are 3 
           additional options
                • Set uid (SUID)
                • Set gid (SGID)
                • Sticky bit
           ----------------------------------------------------------------------
           SUID : -rwsr-xr-x # s here is SUID permission
           When the SUID bit is set on an executable file, this means 
           that the file will be executed with the same permissions as the 
           <owner of this file> (for example - root).
           Practical example: If you look at the binary executable file
           of the passwd command, it has the SUID bit set:
              $ ls -lai /usr/bin/passwd 
                2885961 -rwsr-xr-x 1 root root 59976 Feb  6 14:54 /usr/bin/passwd
           This means that any user running the passwd command will
           be running it with the same permission as root. 
           The passwd command needs to edit files like /etc/passwd,
           /etc/shadow to change the password. These files are owned 
           by root and can only be modified by root.
           But thanks to the setuid flag (SUID bit), a regular user will also be
           able to modify these files (that are owned by root) and change his/her password.   
           This is the reason why you can use the passwd command to 
           change your own password despite of the fact that the files
           this command modifies are owned by root.
           $ stat /usr/bin/passwd
                File: /usr/bin/passwd
                Size: 59976     	Blocks: 120        IO Block: 4096   regular file
                Device: 10303h/66307d	Inode: 2885961     Links: 1
                Access: (4755/-rwsr-xr-x)  Uid: (    0/    root)   Gid: (    0/    root)
          How to set SUID bit?
          I find the symbolic way easier while setting SUID bit.
          You can use the chmod command in this way:

           $ chmod u+s file_name
           
          Difference between small s and capital S as SUID bit
          The S as SUID flag means there is an error that you should look into.
          You want the file to be executed with the same permission as the owner
          but there is no executable permission on the file.
          Which means that not even the owner is allowed to execute 
          the file and if file cannot be executed, you won’t get
          the permission as the owner.  
          -----------------------------------------------------------------------
           SGID # example is "s" within : drwx rws r-x
           is similar to SUID. With the SGID bit is set, 
           any user executing the file will have same permissions 
           as the group owner of the file.
           It’s benefit is in handling the directory. 
           When SGID permission is applied to a directory, all sub directories
           and files created inside this directory will get the same group 
           ownership as main directory (not the group ownership of the user that
           created the files and directories).
           How to set SGID?
           You can set the SGID bit in symbolic mode like this:

           $ chmod g+s directory_name
          -----------------------------------------------------------------------
           The Sticky Bit  # Its example is "t" within drwx rwx rwt  
           Steaky bit: 
           • for the files: in the past used for file mark 
           • for the directory: the sticky bit sets shared directory:
             every user can create in this directory file, but the other user 
             cannot delete this files that are not of the user. The user 
             can update the files of other users but cannot delete them (because 
             the deletion is related to directory permissions).             
           
           Sticky Bit works mostly on the directory. 
          
           With sticky bit set on a directory, all the files
           in the directory can only be DELTED or RENAMED by the
           file owners ONLY or the root.
           
           This is typically used in the /tmp directory that works 
           as the trash can of temporary files. This means that a user (except root)
           cannot delete the temporary files created by other users in the /tmp directory.
           
             $ ls -ldi /tmp
                5242881 drwxrwxrwt 25 root root 20480 Feb 20 12:45 /tmp
             $ stat /tmp
                File: /tmp
                Size: 20480     	Blocks: 48         IO Block: 4096   directory
                Device: 10303h/66307d	Inode: 5242881     Links: 24
                Access: (1777/drwxrwxrwt)  Uid: (    0/    root)   Gid: (    0/    root)
           How to set the sticky bit?
           As always, you can use both symbolic and numeric mode to set the sticky 
           bit in Linux.
              $ chmod +t my_dir
 
\end{verbatim}

\subsection{ln     : Creates links to files and directories }
\begin{verbatim}
    
    Creates links to files and directories.
    Two types of links in Linux:

    • Symbolic links
    • Hard links

    - Create a symbolic link to a file or directory:
         ln -s {{/path/to/file_or_directory}} {{path/to/symlink}}

    - Create a hard link to a file:
         ln {{/path/to/file}} {{path/to/hardlink}}
         
    - Overwrite an existing symbolic link to point to a different file:
         ln -sf {{/path/to/new_file}} {{path/to/symlink}}     
         

    To create symbolic link (-s option):
       $ ln -s 1.txt 1_soft.txt

    To create hard link: 
       $ ln 1.txt 1_hard.txt

    Linux file system EXT and hard link
    How files are stored in disk: 
    Linux file system name is EXT (Extended). 
    All files are stored with 3 linked parts: 
         file name-->inode(information node)-->data of the file
    --------------------------------------------------------------
    • Hard link functions as original file , 
      Hard link is simply new <file name> that linked 
      to the inode of the original file. 
      Hard link is only additional name to the same file (inode + data) !!!
      Hard link is difficult to recognize. It can be recognized only 
      by inode number. If inode number of the file is similar to inode
      number of another file – it is hard link.
      Sector on disk is analogues to address.
      So hard link must be located on the same partition, 
      Hard link can’t link to another partition or the network address of     
      folders!!!  For example, we cannot create hard link on USB disk with 
      link to file on the hard disk.
      User cannot create hard link for folders, because recursive command on 
      folders tree with hard link would cause to infinite loop.
      But he can create soft link to 
      folders. 
      Operational system Linux creates hard link to directories:
      "." is a hard link to current directory.
      ".." is the hard link with same inode in subdirs as "." in top dir.
      Every folder contain at list one hard link named “.”. 
      Every subfolder contains hard link “..”.
      
      Hard link counter of folder less 1 shows how many subfolders has the folder.       
      $ mkdir dir31
      $ mkdir dir31/dir32
      $ mkdir dir31/dir32/dir33
      
      $ ls -li

      1502894 -rw-rw-r-- 1 amit amit   86 Sep 22 05:42 2_txt
      1502893 drwxrwxr-x 3 amit amit 4096 Sep 22 07:53 dir31  # 2 subfolders
      
      rm of hard link does not delete the original file !
      rm of orogonal file does not delete the hard link and file contents !!! 
      In the inode there is inode counter. If we delete the hard link 
      (additional name of the same file), inode counter is diminished by 1. 
      IF inode counter reaches zero – the operational system knows that the 
      space in the disk in empty and it can be used for other purposes.
      
      $ ls -li  # here inode counter is 2
       7753157 -rw-rw-r-- 2 amits amits       45 Feb 19 19:26  1.txt
       7753157 -rw-rw-r-- 2 amits amits       45 Feb 19 19:26  hard.1.txt

      
        $ ln 1.txt 1.hard
    
         File name     inode        data
         1.txt    ---> 123     ---> How do you do ?
                       Shemuel    
                         ^
         1.hard   _______|       

        Example : 
        $ ln 1.txt hard.1.txt
        $ stat hard.1.txt
            File: hard.1.txt
            Size: 45        	Blocks: 8          IO Block: 4096   regular file
            Device: 10303h/66307d	Inode: 7764243     Links: 2
            Access: (0664/-rw-rw-r--)  Uid: ( 1001/   amits)   Gid: ( 1001/   amits)
            Access: 2024-02-19 15:46:06.573765125 +0200
            Modify: 2024-02-15 16:31:55.646124081 +0200
            Change: 2024-02-19 15:45:28.276807762 +0200
            Birth: 2024-02-15 16:31:55.646124081 +0200
        $ stat 1.txt
            File: 1.txt
            Size: 45        	Blocks: 8          IO Block: 4096   regular file
            Device: 10303h/66307d	Inode: 7764243     Links: 2
            Access: (0664/-rw-rw-r--)  Uid: ( 1001/   amits)   Gid: ( 1001/   amits)
            Access: 2024-02-19 15:46:06.573765125 +0200
            Modify: 2024-02-15 16:31:55.646124081 +0200
            Change: 2024-02-19 15:45:28.276807762 +0200
            Birth: 2024-02-15 16:31:55.646124081 +0200
            
        $ ls -i # returns inode numbers of files !!!
           $ ls -i
               7764243  1.txt
               7764243  hard.1.txt

    -------------------------------------------------------------
    • Soft link has <its own inode> and the data that includes
      path to linked file. 
      Soft link is separate file with data that contain address of 
      linked file.
      Soft link is signed with "l" at the head of list of permissions: 
      Link counter of soft link is always 1.
       $ ls -la         
         lrwxrwxrwx 1 amits amits    5 Feb 19 21:15 soft1.txt -> a.txt
    
         File name     inode        data
         1.txt    ---> 123     ---> How do you do ?
                       Shemuel    
         
         1.soft   ---> 456     ---> ./1.txt
                       L      
    ------------------------------------------------------------
    Sector on disk is analogues to address.
	So hard link must be located on the same partition, 
	Hard link can’t link to another partition or the network address of 	
	folders!!!  For example, we cannot create hard link on USB disk with 
	link to file on the hard disk.
    ------------------------------------------------------------    
           
\end{verbatim}



\subsection{find  : Search for a file}
\begin{verbatim}
          The find command in UNIX is a command line utility for
          walking a file hierarchy. 
          It can be used to find files and directories and perform   
          subsequent operations on them. 
          It supports searching by file, folder, name, creation date,   
          modification date, owner and permissions. 
          By using the ‘-exec’ 
          other UNIX commands can be executed on files or folders 
          that found. 

          Syntax: 
              $ find [where to start searching from]
                     [expression determines what to find] 
                     [-options] [what to find]

          1. Search for a file with a specific name.
                $ find ./GFG -name sample.txt

          2. Search for a file with pattern.
                $ find ./GFG -name *.txt

          3. Find and delete a file with confirmation.
                $ find ./GFG -name sample.txt -exec rm -i {} \; 

          4. Find empty files and directories.
                $ find ./GFG -empty

          5. Search for file with given permissions.
                $ find ./GFG -perm 66

          6. Displays repositories and sub-repositories
                $ find . -type d

          7. Search text 'Geek' within multiple text files.
                $ find ./ -type f -name "*.txt" -exec grep 'Geek'  {} \;
          
          - Find files by extension:
              find {{root_path}} -name '{{*.ext}}'

          - Find files matching multiple path/name patterns:
                 find {{root_path}} -path '{{**/path/**/*.ext}}' 
                 -or -name '{{*pattern*}}'

          - Find directories matching a given name,
            in case-insensitive mode:
                 find {{root_path}} -type d -iname '{{*lib*}}'

          - Find files matching a given pattern, excluding
            specific paths:
                 find {{root_path}} -name '{{*.py}}' -not 
                 -path '{{*/site-packages/*}}'

          - Find files matching a given size range, limiting
            the recursive depth to "1":
                 find {{root_path}} -maxdepth 1 -size
                 {{+500k}} -size {{-10M}}

          - Run a command for each file (use {} within the command to          
            access the filename):
                 find {{root_path}} -name '{{*.ext}}' 
                 -exec {{wc -l {} }}\;

          - Find files modified in the last 7 days:
                find {{root_path}} -daystart -mtime -{{7}}

          - Find empty (0 byte) files and delete them:
                find {{root_path}} -type {{f}} -empty -delete

          Examples : 

          Find all files *.c starting from current directory:
               $ find . -name "*.c"

          Find all files *.jpg starting with root directory: 
               $ find / -name "*.jpg" -print
                           
          Find all files *.jpg starting from home directory:   
               $ find ~ -name "*.jpg" -print

          Insert “” to make globing at find command level, 
          no at bash level,  before expansion for find command:

               $ find -name "*.txt"
                  ./amit.txt
                  ./docs/11.txt
                  ./docs/.secret.txt
                  ./docs/5.txt
                  ./docs/7.txt

\end{verbatim}




\subsection{touch    : Create files or update access times}
\begin{verbatim}
     Update the access and modification times if the file exists
     A FILE argument that does not exist is created empty, unless -c   
     or -h is supplied.      
     
     1. In its default usage, it is the equivalent of creating
        or opening a file and saving it without any change to the file  
        contents. touch avoids opening, saving, and closing the file.  
        Instead it simply updates the dates associated with the file 
        or directory. An updated access or modification date can be 
        important for a variety of other programs such as backup 
        utilities.
     2. The touch command can also be useful for quickly creating 
        files for programs or scripts that require a file with a 
        specific name to exist for successful operation of the 
        program, but do not require the file to have any specific          
        content.
        
		Example:
		     create files 1.txt , 2.txt, ... , 5.txt
		     touch {1..5}.txt
		           
        
\end{verbatim}

\subsection{stat    : display file (inod) and file system info}
\begin{verbatim}
     stat() is a Unix system call that returns file attributes
     about an inode. 
     
     The semantics of stat() vary between operating systems. 
     As an example, Unix command ls uses this system call to retrieve
     information on files that includes:


     atime: time of last access (ls -lu)
     mtime: time of last modification (ls -l)
     ctime: time of last status change (ls -lc)
     
     
     The stat() and lstat() functions take a filename argument.
     If the file is a symbolic link, stat() returns attributes of the 
     eventual target of the link, 
     while lstat() returns attributes of the link itself. 
     
     File system information:   
     $ stat -f try.py 
         File: "try.py"
         ID: 5a2e6f11fdd5900a Namelen: 255     Type: ext2/ext3
         Block size: 4096       Fundamental block size: 4096
         Blocks: Total: 120454536  Free: 91353786   Available: 85216660
         Inodes: Total: 30670848   Free: 29126190
      
\end{verbatim}


\subsection{ls       : "LiSt" files and dirs}
\begin{verbatim}
    
    Lists the contents of files and directories. 
    Without options and arguments displays a list 
    of the names of all files in the current working directory
    in alphabetical order
       -i : return inode numbers of files !!
       
       -l : displays detailed information about files and directories
            in a long format

       -a : includes hidden files and directories in the listing

       -t : sorts files and directories by their last modification time,
            most recently modified ones first

       -S : sorts files and directories by their sizes,
            listing the largest ones first

       -r : reverse order (most recent at the end)

       -R : lists files and directories recursively, including subdirectories

       -h : prints file sizes in a human-readable format (e.g., 1K, 234M, 2G) 

       -d : info for directories only.
   

    ls command examples:
      ls           # files list (alef bet order)
      ls /         # root dir
      ls .         # current dir
      ls           # current dir
      ls ~         # current dir
      ls -l        # display detailed info: long format
      ls -a        # include hidden files
      ls -lh       # size of files in human format
      ls -lah      # detailed info long format, include hidden files
                     size of files in human format
      ls [0-9].txt # 
      ls -R        # include subdirs
      ls -ld docs  # dir info only
      ls -t        # most recent first
      ls -r        # oldest first
      ls -S        # sort by size (largest first)
      ll           # alias of ls -l

    file types:
         - : regular file
         d : directory (file listing a set of files)
         l : soft link file (Files referring to the name of another file
         s : socket file (inter process communication)
         p : pipe file (Used to cascade programs)
         c : char file
         b : block file driver (disk) communication files
         
         search hidden files:
         ls *.txt .*.txt
         
         $ ls [0-9].txt
             1.txt  3.txt  5.txt  7.txt
             
         amits:temp1$ ls -l [a-c].txt
             -rw-rw-r-- 1 amits amits 0 Feb 15 13:55 a.txt
             -rw-rw-r-- 1 amits amits 0 Feb 15 13:55 b.txt
             -rw-rw-r-- 1 amits amits 0 Feb 15 13:55 c.txt
         amits:temp1$ ls -l {a..c}.txt                      # same result
             -rw-rw-r-- 1 amits amits 0 Feb 15 13:55 a.txt
             -rw-rw-r-- 1 amits amits 0 Feb 15 13:55 b.txt
             -rw-rw-r-- 1 amits amits 0 Feb 15 13:55 c.txt
         $ ls -i # returns inode numbers of files !!!
           7764243  1.txt
           7764243  hard.1.txt
\end{verbatim}


\subsection{du       : "Disk Usage" - file and dir usage space}
\begin{verbatim}
      
      https://phoenixnap.com/kb/show-linux-directory-size
      
      Display the current directory size by running du without any arguments:
      The system lists the current directory content and subdirectories with a    
      number representing the object size in kilobytes. 
      The current directory size is the last line.
       $ du
      Print the Current Directory Size in a Human-Friendly Format
      Add the -h option to the du command to make the output easily readable:
       $ du -h
      
      Find the Size of a Specific Directory
       $ du -h ./example1
      
      To print an extra total line at the end of the output and summarize the disk     
      space usage for all files and directories, use the -c argument:
       $ du -c ./example1
      
      
      
      Disk usage: estimate and summarize file and directory space usage
       - List the sizes of a directory and any subdirectories, in human-   
         readable form (i.e. auto-selecting the appropriate unit for each
         size):
            $ du -h {{path/to/directory}}

       - List the human-readable sizes of a directory and of all the files and 
         directories within it:
            $ du -ah {{path/to/directory}}


        Summary of disk usage for two directories:
           $ du -h
             8.0K	./dir2/2024
             12K	./dir2
             8.0K	./dir1/2024
             20K	./dir1
             76K	.
           $ du -h | sort -h
             8.0K	./dir1/2024
             8.0K	./dir2/2024
             12K	./dir2
             20K	./dir1
             76K	.
           $ du -h -s ./dir1 ./dir2  # -s is for summary of size
             20K	./dir1
             12K	./dir2
             
         Apparent size is the actual amount of data that can be read from the 
         file. Block-oriented devices can only store in terms of blocks, not 
         bytes. As a result, the disk usage is always rounded up to the next 
         highest block     
           $ du -h -s --apparent-size ./dir1 ./dir2  
            
            
             
\end{verbatim}


\subsection{ncdu       : "Norton Commander and Disk Usage" - file and dir usage space}
\begin{verbatim}
       - Analyze the current working directory:
               $ ncdu

       - Colorize output:
               $ ncdu --color {{dark|off}}  # dark is color , off is regular

       - Analyze a given directory:
               $ ncdu {{path/to/directory}}

       - Save results to a file:
               $ ncdu -o {{path/to/file}}
\end{verbatim}

\subsection{fdisk       : partitions in hard disk}
\begin{verbatim}

 - List partitions:
    $ sudo fdisk -l

\end{verbatim}

\subsection{lsblk       : Lists information about block(storage) devices}
\begin{verbatim}
      - List all storage devices in a tree-like format:
          $ lsblk

       - Also list empty devices:
          $ lsblk -a

       - Print the SIZE column in bytes rather than in a 
         human-readable format:
          $ lsblk -b

       - Output info about filesystems:
          $ lsblk -f
\end{verbatim}


\subsection{df       : overview of the filesystem disk space usage}
\begin{verbatim}
      
       - Display all filesystems and their disk usage in human-readable form:
         $ df -h

      - Display all filesystems and their disk usage:
         $ df
\end{verbatim}


\subsection{cd       : "Change Dir"}
\begin{verbatim}
      go to parent directory : cd ..
      go to home directory   : cd ~
      go to root directory   : cd /
      go to prev. dir.       : cd -  

      $ cd -
        /home/amit
      $ cd -
        /home/amit/Documents

\end{verbatim}


\subsection{pwd      : "Print Working Directory"}
\begin{verbatim}

\end{verbatim}



\subsection{mkdir    : "MaKe DIRectory"}
\begin{verbatim}
    The mkdir command in Linux allows the user to create directories. 
    This command can create multiple directories at once 
    as well as set the permissions for the directories.

    This command can create multiple directories at once
    as well as set the permissions for the directories.

   - Create specific directories:
        mkdir {{path/to/directory1 path/to/directory2 ...}}

   - Create specific directories and their [p]arents if needed:
        mkdir -p {{path/to/directory1 path/to/directory2 ...}}
   -p to create parent folders : 
        $ mkdir -p aa/bb/cc

   - Create directories with specific permissions:
        mkdir -m {{rwxrw-r--}} {{path/to/directory1 path/to/directory2 ...}}



     $ mkdir aa/bb/cc
        mkdir: cannot create directory ‘aa/bb/cc’: No such file
               or directory
     $ mkdir -p aa/bb/cc  # now it is performed     
     
\end{verbatim}

\subsection{cp       : "CoPy" file to file, file to dir, dir to dir}
\begin{verbatim}
        Copy files and directories
     
     By default, the cp command copies files without preserving 
     their attributes such as permissions, timestamps, and ownership. 
     
     Recursively copy a directory's contents : 
              cp -R {source_directory} {target_directory}
              cp -r {source_directory} {target_directory}
        When the program has two arguments of path names to <files>, the program      
     copies the contents of the first <file> to the second <file>, creating  
     the second <file> if necessary.
     
        When the program has one or more arguments of path names of <files> and   
     following those an argument of a path to a <directory>, then the program   
     copies <each source file> to <the destination directory>, creating any 
     files not already existing (!).

        When the program's arguments are the path names to two <directories>, cp
     copies all files in the source <directory> to the destination 
     <director>y, creating any files or directories needed. This mode of 
     operation requires an additional option flag, typically -r, to indicate 
     the recursive copying of directories. If the destination directory 
     already exists, the source is copied into the destination, while a new 
     directory is created if the destination does not exist.
     
     Example:
        mkdir temp{1..5}
    
    
    Renaming a file by copying and deleting it
    Linux users copy a file by using the “cp” command. 
    When you copy a file, you give the source files and rename the files. 
      $ cp old_file new_file
    As an example, if we were to rename “august.png” to “september.png”:
      $ cp august.png september.png
    This will create a copy of the file with a new name. 
    But now you’ll have two copies. 
    You’ll resolve this by deleting the old file. 
    You delete with “rm.”

      $ rm august.png
    
     
\end{verbatim}



\subsection{mv       : "MoVe" files or dirs from one place to another.}
\begin{verbatim}
         If both filenames are on the same filesystem, this results
     in a simple file rename; otherwise the file content is copied
     to the new location and the old file is removed.
     
         Using mv requires the user to have write permission for 
     the directories the file will move between. This is because mv 
     changes the content of both directories (i.e., the source
     and the target) involved in the move. 
     When using the mv command on files located on the same filesystem, 
     the file's timestamp is not updated.
     
      rm of hard link does not delete the original file !
      rm of original file does not delete the hard link and file contents that is linked by it!!!
      

     Example : move interactive (-i) : 
     $ mv 1.txt b.txt -i
         cp: overwrite ‘b.txt’? y

     $ mv dir1 dir2     # move dir1 under dir2
     $ mv dir2/dir1 ./  # move dir1 under current dir

     You can rename a file by “moving it” to a new file name. 
         $ mv august.png september.png
     This actually does what you would think of as a “rename” command and in a single step. 
     But it should be noted that it’s also the general “move” command. 
     So, you are moving the file at the same time. 
     If you move it to a different directory, it will be moved as well as renamed.
\end{verbatim}


\subsection{rm       : "ReMove" removes references(!) to objects from the filesystem}
\begin{verbatim}
     
     The rm command removes references to objects 
     from the filesystem using the unlink system call,
     where those objects might have had multiple references
     (for example, a file with two different names), 
     and the objects themselves are discarded only when 
     all references have been removed and no programs 
     still have open handles to the objects.
   
     The command generally does not destroy file data,
     since its purpose is really merely to unlink references,
     and the filesystem space freed may still contain 
     leftover data from the removed file.
     This can be a security concern in some cases, 
     and hardened versions sometimes provide for 
     wiping out the data as the last link is being cut, 
     and programs such as <shred> and <srm> are available
     which specifically provide data wiping capability.
     
     rm    : Remove specific files
     
     rm -i : Remove specific files [i]nteractively 
             prompting before each removal.
     rm -r : (r)ecursively delete a directory and all its contents (!) 
     
     Note : normally rm will not delete directories, 
            while rmdir will only delete empty
            directories). But -r allows to remove directories:
            
            $ rm dir3
                 rm: cannot remove “dir3” Is a directory
            $ rm dir3 -r # remove dir3 and all child directories
            $            # Done directory remove
            
            $ rm dir1/{2..4}.txt
            $ rm dir2/{1,4,5}.txt
      
      Remove all files with answer "yes" to all prompts : 
            $ yes | rm -r Documents/      

\end{verbatim}

\subsection{rmdir    : "ReMove DIR" remove empty(!) dir from filesystem.}
\begin{verbatim}
     Note : removes only empty directories !!!
     
     $ rmdir dir1
        rmdir: failed to remove “dir1”: Directory not empty

     rmdir     : remove specific empty directory
     rmdir -p  : remove specific directory and all its ancestors - "AVOT" 
\end{verbatim}




\subsection{tree    : show curr. dir as tree.}
\begin{verbatim}
       Show current directory as tree , 
       with sizes in human style and permissions : 
               
       $ tree --du -h -p       
       
       -L     Level
       -s     Print the size of each file in bytes along with the name.
       -h     Print the size of each file but in a more  human  readable  way,
              e.g.  appending  a size letter for kilobytes (K), megabytes (M),
              gigabytes (G), terabytes (T), petabytes (P) and exabytes (E).             
       
       --du   For each directory report its size as the accumulation of  sizes
              of  all  its  files and sub-directories
              

       --inodes
              Prints the inode number of the file or directory
       
       -p     Print  the  file  type  and permissions for each file (as per ls
              -l).
     
       
       -u     Print the username, or UID # if no username is available, of the
              file.
       -g     Print the group name, or GID # if no group name is available, of
              the file.
            
      - Print files and directories up to 'num' levels of 
        depth (where 1 means the current directory):
           $ tree -L {{num}}

      - Print directories only:
           $ tree -d

      - Print hidden files too with colorization on:
           $ tree -a -C
       
       
\end{verbatim}



\subsection{file    : determine file type.}
\begin{verbatim}
 - Give a description of the type of the 
   specified file. Works fine for files with no
   file extension:  
        file {{path/to/file}}
   -b, –brief : This is used to display just
    file type in brief mode
        
   
        
   Examples:
       $ type -a cat       
       $ file [a-z]*
       $ file  /bin/cat 
       $ file tex_file_data_extract       
        
       $ file $(which test) # $( ) means substitute command
         /usr/bin/test: ELF 64-bit LSB pie executable, x86-64, version 1 (SYSV), 
         dynamically linked, interpreter /lib64/ld-linux-x86-64.so.2, 
         BuildID[sha1]=68ccce12c6a2a96d3a40fb5868f8b8c45e6015a5, for GNU/Linux 3.2.0, 
         stripped
       $ file `which test`  # ` ` means substitute command
         /usr/bin/test: ELF 64-bit LSB pie executable, x86-64, version 1 (SYSV), 
         dynamically linked, interpreter /lib64/ld-linux-x86-64.so.2, 
         BuildID[sha1]=68ccce12c6a2a96d3a40fb5868f8b8c45e6015a5, for GNU/Linux 3.2.0, 
         stripped        
\end{verbatim}


\section{User Groups }

\subsection{CreateNewUser - steps to create new user  }
\begin{verbatim}
     
      1. Create Home    : $ sudo useradd tal -m # new user with home directory
      2. Skeleton files : copy files from /etc/skel to /home/tal 
                          Note : copy without "." and ".."
      3. set shell type : $ sudo usermod -s /bin/bash tal            
      4. set pw         : sudo passwd tal
      5. add user to sudo group
                        : $ sudo usermod -aG sudo david
         Debian : "sudo"         
         RedHat : "wheel"    
            or...
         visudo (edit sudoer file)   
\end{verbatim}

\subsection{visudo - edit sudoer file  }
\begin{verbatim}
        Safely edit the sudoers file.

        - Edit the sudoers file:
          sudo visudo

        - Check the sudoers file for errors:
          sudo visudo -c

 - Edit the sudoers file using a specific editor:
   sudo EDITOR={{editor}} visudo

 - Display version information:
   visudo --version
    
  
       
\end{verbatim}




\subsection{passwd   : Change a user's PASSWorD.}
\begin{verbatim}
     
     Change user password
     
     - Change the password of the current user interactively:
         passwd

     - Change the password of a specific user:
         passwd {{username}}

     - Get the current status of the user:
         passwd -S

     - Make the password of the account blank 
       (it will set the named account passwordless):
         passwd -d
\end{verbatim}


\subsection{useradd   : add user to the system}
\begin{verbatim}
    
    It is not simple, more complicated
    relative to adduser (simple perl script that uses useradd)

    When we use useradd, it creates the account without password : 
      $ sudo useradd tal
        password for amit:       
      $ 

    But we cannot enter the account with su – tal 
    because it has not user: 
      $ su - tal
        Password: 
        scdsds
        ssu: Authentication failure

    To give to user tal password, run following command (with sudo):
      $ passwd tal
        passwd: You may not view or modify password information for tal.
      $ sudo passwd tal
        New password: 
        Retype new password: 
        passwd: password updated successfully
        
     Create user with home directory as his name use $ useradd -m   
      $ sudo useradd tal -m
      $ su - tal
        Password: 
      $ sudo passwd tal
        New password: 
        Retype new password: 
        passwd: password updated successfully
      $ su - tal
        Password: 
      $ 
      $ pwd
        /home/tal
      $ echo $SHELL  # when we create user with useradd , shell will be 
        /bin/sh 
   
        By default the user tal uses shell "dash", to define shell to user 
        we need define option -s /bin/bash when we define new user with 
        useradd command
        But now, when the user already exists, we can change the shell 
        type with usermod command.

      $ sudo usermod -s /bin/bash tal           
      $ su - tal
        Password: 
      $ ll
        total 36
        drwxr-x--- 4 tal  tal  4096 Sep 27 16:02 ./
        drwxr-xr-x 5 root root 4096 Sep 27 15:58 ../
        -rw-r--r-- 1 tal  tal   220 Jan  6  2022 .bash_logout
        -rw-r--r-- 1 tal  tal  3771 Jan  6  2022 .bashrc
        drwx------ 2 tal  tal  4096 Sep 27 16:02 .cache/
        drwxr-xr-x 5 tal  tal  4096 Jul 11 19:21 .config/
        -rw-r--r-- 1 tal  tal    22 Sep  8  2011 .gtkrc-2.0
        -rw-r--r-- 1 tal  tal   516 Dec 17  2013 .gtkrc-xfce
        -rw-r--r-- 1 tal  tal   807 Jan  6  2022 .profile   
        
        Above skeleton files are within /etc/skel directory, 
        and can be copied from it manually. But when we copy the file from
        one user to another user, it will have the user name of previous
        user !!!
        There are two ways to solve this:
          • Allow the user to copy files.
          • Change the user of the files.

        Note: don’t run command cp -r /etc/skel/.* because it will copy          
        also “.” And “..” that is not desirable.

        So it is preferred to copy files by exact list of its names. 
        
        Thinks to concern when we create new user: 
          • User’s home
          • User’s skeleton files
          • User’s shell
          • User’s password
        
        Note: for useradd don’t use -P option !!!
        -P, --prefix PREFIX_DIR  # Gibrish
           Apply changes in the PREFIX_DIR directory and use the 
           configuration files from the PREFIX_DIR directory.
           This option does not chroot and is intended for preparing a
           cross-compilation target. Some limitations: NIS and LDAP 
           users/groups are not verified. PAM authentication is using
           the host files. No SELINUX support.

\end{verbatim}

\subsection{\$SHELL   : env. variable (check current shell type) }
\begin{verbatim}
      https://unix.stackexchange.com/questions/43499/difference-between-
      echo-shell-and-which-bash
      
      $ echo $SHELL
       /bin/zsh
      $ which bash
       /bin/bash
       The first command tells me which shell will be executed by login
       when you log in. In my case, /bin/zsh.
       The second command tells me the first occurrence in my
       $PATH the bash command can be found.

\end{verbatim}



\subsection{adduser   : perl script, that uses useradd}
\begin{verbatim}

        simple perl script (that uses useradd command internally):
          
          $ sudo adduser david
            [sudo] password for amit:       
            Adding user `david’ ...
            Adding new group `david’ (1001) ...
            Adding new user `david’ (1001) with group `david’; ...
            Creating home directory `/home/david’ ...
            Copying files from `/etc/skel’ ...
            New password: 
            Retype new password: 
            passwd: password updated successfully
            Changing the user information for david
            Enter the new value, or press ENTER for the default
	        Full Name []: David
	        Room Number []: 1
	        Work Phone []: 046778723
	        Home Phone []: 0584266865
	        Other []: no
            Is the information correct? [Y/n] Y

          $ cat /etc/passwd | grep bash # check users of bash
             root:x:0:0:root:/root:/bin/bash
             amits:x:1001:1001:AmitS,,,:/home/amits:/bin/bash
             david:x:1000:1000:David,1,046778723,0584266865,no:/home/david:/bin/bash
\end{verbatim}


\subsection{userdel   : delete user from the system}
\begin{verbatim}

    it is not simple, more complicated
    relative to deluser (simple perl script that uses userdel)

    Remove the user:
       $ userdel -r moshe
      
      Files in the user’s home directory will be removed 
      along with the home directory itself and the user’s 
      mail spool. Files located in other file systems
      will have to be searched for and deleted manually.
\end{verbatim}

\subsection{deluser   : perl script, that uses userdel}
\begin{verbatim}

        It is simple Perl script (that uses userdel command)


        Delete a user account and related files.
        The userdel command modifies the system account files, 
        deleting all entries that refer to the user name LOGIN.
        The named user must exist.
        
        Try to delete user tal: 
          $ sudo deluser --remove-all-files tal 
           ….
           /usr/sbin/deluser: Cannot handle special file /proc/3614/fd/1
           /usr/sbin/deluser: Cannot handle special file /proc/3614/fd/2
           /usr/sbin/deluser: Cannot handle special file /etc/systemd/   
           system/sudo.service
           Removing user `tal’ ...
           Warning: group `tal’ has no more members.
           Done.
\end{verbatim}


\subsection{/etc/group - file with list of groups }
\begin{verbatim}
         $ head /etc/group
           root:x:0:
           daemon:x:1:
           bin:x:2:
           sys:x:3:
           adm:x:4:syslog,amits
           tty:x:5:syslog
           disk:x:6:
           lp:x:7:
           mail:x:8:
           news:x:9:

          $ cat /etc/group | grep amits
           adm:x:4:syslog,amits
           cdrom:x:24:amits
           sudo:x:27:amits
           dip:x:30:amits
           plugdev:x:46:amits
           lpadmin:x:120:amits
           sambashare:x:131:amits
           amits:x:1001:


     check who has permission to use sambashare (program to share files):
         $ cat /etc/group | grep sambashare
           sambashare:x:136:amit
\end{verbatim}

\subsection{/etc/passwd   : conteins info what users are in the system}
\begin{verbatim}
       $ cat /etc/passwd | grep amit
           amits:x:1001:1001:AmitS,,,:/home/amits:/bin/bash  
                      # /bin/bash : what shell the user uses
\end{verbatim}


\subsection{groups   : present all groups of the user}
\begin{verbatim}
     groups - print the groups a user is in
       
     $ groups david
       david : david sudo
     $ groups amits
       amits : amits adm cdrom sudo dip plugdev lpadmin sambashare
     $ groups
       amits adm cdrom sudo dip plugdev lpadmin sambashare  

\end{verbatim}

\subsection{chown   : CHange OWNer (user/group) of files and dirs}
\begin{verbatim}
     
     Change user and group ownership of files and directories.
         
     Unprivileged (regular) users who wish to change the group membership
     of a file that they own may use chgrp. 

     - Change the owner user of a file/directory:
        chown {{user}} {{path/to/file_or_directory}}

     - Change the owner user and group of a file/directory:
        chown {{user}}:{{group}} {{path/to/file_or_directory}}

     - Recursively change the owner of a directory and its contents:
        chown -R {{user}} {{path/to/directory}}

     - Change the owner of a symbolic link:
        chown -h {{user}} {{path/to/symlink}}

     - Change the owner of a file/directory to match a reference file:
        chown --reference={{path/to/reference_file}} {{path/to/file_or_directory}}
     
     
     
         
        
     $ mkdir project_dir
     $ sudo chown david:family project_dir  # david : owner , family : group

     
\end{verbatim}




\subsection{members   : present all users of the group}
\begin{verbatim}
     $ groups
        amits adm cdrom sudo dip plugdev lpadmin sambashare
     $ members lpadmin
        cups-pk-helper amits 
     $ members sudo
        amits david
\end{verbatim}

\subsection{groupadd   : Add user groups to the system}
\begin{verbatim}

    Add user groups to the system. See also: groups, groupdel, groupmod 
     - Create a new group:
        $ sudo groupadd {{group_name}}

     - Create a new system group:
        $ sudo groupadd --system {{group_name}}

     - Create a new group with the specific groupid:
        $ sudo groupadd --gid {{id}} {{group_name}}

     $ sudo groupadd family
     $ sudo usermod -aG family amits
     $ groups amits
       amits : amits adm cdrom sudo dip plugdev lpadmin sambashare family
     $ groups
       amits adm cdrom sudo dip plugdev lpadmin sambashare


\end{verbatim}



\subsection{passwd   : change password}
\begin{verbatim}
    Change password for current user : 
       passwd
\end{verbatim}

\subsection{whoami   : currently logged-in user}
\begin{verbatim}
    allows users to see the currently logged-in user.
\end{verbatim}


\subsection{su   : "Substitute User"}
\begin{verbatim}
      The 'su' command, short for 'substitute user', is a simple 
      yet powerful command in Linux. It allows you to switch to 
      another user account on your system, right from the command 
      line. This can be particularly useful for system 
      administrators who need to perform tasks on behalf of other 
      users.
      
      To enter David’s user account including David’s environment
      variables:
      $ su - david

      enter and exit from user david :
      $ whoami
        amits
      $ su david
        password:
      $ whoami
        david 
      $ exit
      $ whoami
        amits 
        
      To connect as root user: 
         $ sudo su – root # or sudo su -.
      It allows to work without sudo for every command. 
      Don’t forget to run exit.  
      
         $ sudo su - root
           [sudo] password for amit:       
         # 
         #
         # whoami
           root
         # 

       Note: "#" signs that the user is root.
       Note : sudo is not user, sudo is group
          
\end{verbatim}

\subsection{exit   : exit the shell}
\begin{verbatim}

      $ whoami
        amits
      $ su david
        password:
      $ whoami
        david 
      $ exit    # exit from david
      $ whoami
        amits   
\end{verbatim}




\subsection{usermod   : MODifies a USER account/add user to group.}
\begin{verbatim}
       
        - Change a username:
            sudo usermod --login {{new_username}} {{username}}

        - Change a user id:
            sudo usermod --uid {{id}} {{username}}

        - Change a user shell:
            sudo usermod --shell {{path/to/shell}} {{username}}

        - Add a user to supplementary groups (mind the lack of whitespace):
            sudo usermod --append --groups {{group1,group2,...}} {{username}}

        -a, --append  # the same as --append --groups
            Add the user to the supplementary group(s). 
            Use only with the -G option (!) : -aG

        -G, --groups GROUP1[,GROUP2,...[,GROUPN]]]
            A list of supplementary groups which the user is also a member of.
            Each group is separated from the next by a comma, with no intervening 
            whitespace. The groups are subject to the same restrictions as
            the group given with the -g option.

        - Change a user home directory:
            sudo usermod --move-home --home {{path/to/new_home}} {{username}}

       $ echo $SHELL # check the shell type
          /bin/sh       
       
       Change the shell type for user "david" with usermod command:
            $ sudo usermod -s /bin/bash david
            
       Add the user "david" to group "sudo":
         amits $ su - david
              Password: 
         david $ sudo echo hello
             [sudo] password for david: 
             david is not in the sudoers file.  This incident will be reported.
         david $ exit 
             logout
         amits $ sudo usermod -aG sudo david
             [sudo] password for amits: 
         amits $ su - david
             Password: 
     To run a command as administrator (user "root"), use "sudo <command>". 
     See "man sudo_root" for details.
         david $ sudo echo hello # to check that david is sudoer
             [sudo] password for david: 
             hello 
             
         $ sudo usermod -aG family david # add user david to group family
         $ groups david # check in what groups is david ?
           david : david sudo family
  
\end{verbatim}



\subsection{usermod : add Tal to sudo group and remove tal from sudo group.}
\begin{verbatim}

to add user to sudo group : 
 amits$ sudo usermod -a -G sudo tal

to remove tal from sudo group:
 amits$ sudo usermod -r -G sudo tal
  



\subsection{sudo   : run command as administrator.}
\begin{verbatim}
       
   To run a command as administrator (user   
   "root"), use "sudo <command>".
   See "man sudo_root" for details.
   To be "sudoer" the user must be added to 
   the group "sudo".
   
   $ su - david
     Password: 
   $ sudo echo hello
     [sudo] password for david: 
     david is not in the sudoers file.  This incident will be reported.
   $ exit 
     logout
   amits:Temp_Amit$ sudo usermod -aG sudo david
     [sudo] password for amits: 
   $ su - david
     Password: 
     To run a command as administrator (user "root"), use "sudo <command>". 
     See "man sudo_root" for details.

david@amits-vostro-3520:~$ sudo echo hello # to check that david is sudoer
        [sudo] password for david: 
        hello                      

\end{verbatim}



\section{Special Device files }


\subsection{Device..files with special behaviour }
\begin{verbatim}
        These files are related to drivers, placed in /dev 
        directory and are of two types:
           • char file
           • block file
        Read or write to char or block file causes to some system action. 
        These both files are related to writing data and reading data from
        one place to other place. 
        But the difference between them is how they read/write data.  
        
        Block file:
        -----------
        A block file is a hardware file which read/write data in blocks
        instead of character by character. This type of files are very
        much useful when we want to write/read data in bulk fashion.
        All our disks such are HDD, USB and CDROMs are block devices.
        This is the reason when we are formatting we consider block size.
        The write of data is done in asynchronous fashion and it is 
        CPU intensive activity. These devices files are used to store data
        on real hardware and can be mounted so that we can access the data
        we written. Have a look at our other post on getting block size of a       
        device.

        Block file examples:/dev/loop0, /dev/ram0, /dev/sda1, /dev/sr0.
        
        Character file: 
        ---------------
        A char file is a hardware file which reads/write data in character by    
        character fashion. Some classic examples are keyboard, mouse, serial 
        printer. If a user use a char file for writing data no other 
        user can use same char file to write data which blocks access
        to other user. 
        Character files uses synchronise technic to write data. 
        Of you observe char files are used for communication purpose
        and they can not be mounted.
        Example character files: 
                               /dev/autofs,
                               /dev/console,
                               /dev/crash,
                               /dev/lp0, 
                               /dev/null, 
                               /dev/ppp, 
                               /dev/random, 
                               /dev/tty ,etc
         ls -l output for a char file:
              $ ls -l /dev/null
                  crw-rw-rw- 1 root root 1, 3 Sep 21 06:54 /dev/null
              $ ls -l /dev/zero
                  crw-rw-rw- 1 root root 1, 5 Sep 21 06:54 /dev/zero
              $ ls -l /dev/random
                  crw-rw-rw- 1 root root 1, 8 Sep 21 06:54 /dev/random
              $ ls -l /dev/urandom
                  crw-rw-rw- 1 root root 1, 9 Sep 21 06:54 /dev/urandom

        Note: stdout, stdin and stderr are also placed in /dev directory: 
        (fd/0,1,2 – file descriptor 0,1,2; stdout,stdin,stderr are soft links)
 
  $ ls -li /dev
  482 lrwxrwxrwx   1 root  root   15 Feb 20 12:41 stderr -> /proc/self/fd/2
  480 lrwxrwxrwx   1 root  root   15 Feb 20 12:41 stdin  -> /proc/self/fd/0
  481 lrwxrwxrwx   1 root  root   15 Feb 20 12:41 stdout -> /proc/self/fd/1
 
       
\end{verbatim}


\subsection{/dev/null and /dev/zero : data sink files}
\begin{verbatim}
     Data written to the /dev/null and /dev/zero special files is discarded
     
     Reads  from  /dev/null file always return end of file ^D, 
     whereas reads from /dev/zero file always return bytes 
     containing zero ('\0' characters with 0x0 ASCII code).     

     Read 20 zeros in decimal format from /dev/zero file : 
     $ od -tu1 -N20 /dev/zero
       0000000   0   0   0   0   0   0   0   0   0   0   0   0   0   0   0   0
       0000020   0   0   0   0
       0000024
     
     $ od -tu1 -An -N20 /dev/zero # without present byte count
       0   0   0   0   0   0   0   0   0   0   0   0   0   0   0   0
       0   0   0   0
 

 
     See: $ man null 
     
     If you want to sent the output to nowhere, send it to /dev/null: 
            $ echo "this will not be output to any place" > /dev/null
             
     If you don’t want to see error messages, you can send errors to /dev/null:
            $ ls -l dir1 dir2
                 ls: cannot access ‘dir2’: No such file or directory
                 dir1:
                 total 4
                 -rw-r--r-- 1 amit amit 13 Sep 23 22:51 a.txt
            $ ls -l dir1 dir2 2>/dev/null # rederect errors to dev/null
                 dir1:
                 total 4
                 -rw-r--r-- 1 amit amit 13 Sep 23 22:51 a.txt

     If you want to see only error messages, you can send regular output to /dev/null:
            $ ls -l dir1 dir2 1>/dev/null
                 ls: cannot access ‘dir2’: No such file or directory

     If we don’t want to see messages of any type (error and regular) 
     we can send both outputs to /dev/null:
            $ ls -l dir1 dir2 1>/dev/null 2>&1
\end{verbatim}

\subsection{/dev/random file - random numbers generator file}
\begin{verbatim}
        Read 5 random numbers from /dev/random file:
        Random is quality generator !
          $ od -tu1 -An -N5 /dev/random
            253  85  48 224 101
          $ od -tu1 -An -N5 /dev/random
            166   2 125 190 224
          $ od -tu1 -An -N5 /dev/random
            215 185 237  29 144

        4 byte numbers :
         $ od -tu4 -An -N4 /dev/random # -tu4 is 4 byte decimal number
           1255275280
         $ od -tu4 -An -N4 /dev/random # another random number
           1409517909
         $ od -tu4 -An -N4 /dev/random
           1099118913

        Randomization of /dev/random is of less quality
        Randomization of /dev/urandom in of high quality
         
\end{verbatim}


\subsection{/dev/full - file for device full simulation }
\begin{verbatim}
       If we try to write to /dev/full the system returns that device is full. 
       $ echo 'hello' > /dev/full
          bash: echo: write error: No space left on device
         
\end{verbatim}



\subsection{/etc/group - file with list of groups }
\begin{verbatim}
         $ head /etc/group
           root:x:0:
           daemon:x:1:
           bin:x:2:
           sys:x:3:
           adm:x:4:syslog,amits
           tty:x:5:syslog
           disk:x:6:
           lp:x:7:
           mail:x:8:
           news:x:9:

          $ cat /etc/group | grep amits
           adm:x:4:syslog,amits
           cdrom:x:24:amits
           sudo:x:27:amits
           dip:x:30:amits
           plugdev:x:46:amits
           lpadmin:x:120:amits
           sambashare:x:131:amits
           amits:x:1001:


     check who has permission to use sambashare (program to share files):
         $ cat /etc/group | grep sambashare
           sambashare:x:136:amit
\end{verbatim}

\subsection{/etc/passwd - info what users are in the system }
\begin{verbatim}
         $ head /etc/group
           root:x:0:
           daemon:x:1:
           bin:x:2:
           sys:x:3:
           adm:x:4:syslog,amits
           tty:x:5:syslog
           disk:x:6:
           lp:x:7:
           mail:x:8:
           news:x:9:

          $ cat /etc/group | grep amits
           adm:x:4:syslog,amits
           cdrom:x:24:amits
           sudo:x:27:amits
           dip:x:30:amits
           plugdev:x:46:amits
           lpadmin:x:120:amits
           sambashare:x:131:amits
           amits:x:1001:


     check who has permission to use sambashare (program to share files):
         $ cat /etc/group | grep sambashare
           sambashare:x:136:amit
     
     Users that are using BASH: 
        $ cat /etc/passwd | grep bash
            root:x:0:0:root:/root:/bin/bash
            amit:x:1000:1000:amit,,,:/home/amit:/bin/bash      
           
           
           
\end{verbatim}

\subsection{/etc/skel - skeleton files for new user}

\begin{verbatim}
        Skeleton files (.bashrc , .profile) for new user are 
        within /etc/skel directory, and can be copied from it manually. 
        But when we copy the file from one user to another user, 
        it will have the user name of previous user !!!
        
        There are two ways to solve this:
          • Allow the user to copy files.
          • Change the user of the files.

        Note: don’t run command cp -r /etc/skel/.* because it will copy          
        also “.” And “..” that is not desirable.

        So it is preferred to copy files by exact list of its names. 
        
        $ ls -lai /etc/skel
          total 28
          9699438 drwxr-xr-x   2 root root  4096 Nov 15 10:14 .
          9699329 drwxr-xr-x 154 root root 12288 Feb 20 23:18 ..
          9701502 -rw-r--r--   1 root root   220 Feb 25  2020 .bash_logout
          9701503 -rw-r--r--   1 root root  3771 Feb 25  2020 .bashrc
          9701504 -rw-r--r--   1 root root   807 Feb 25  2020 .profile

\end{verbatim}


\section{Binary files treatment }


\subsection{binary files compare}
\begin{verbatim}

1. Convert files in hex format and compare as unicode files:

$ hexdump file1.bin > file1.hex
$ hexdump file2.bin > file2.hex
$ diff file1.hex file2.hex

2. A visual binary diff tool that allows you to compare binary files       
   interactively.
   
$ vbindiff file1.bin file2.bin


\end{verbatim}

\subsection{nm      :"Name Mangling" damp symbol table from object files.}
\begin{verbatim}
     
     https://en.wikipedia.org/wiki/Nm_(Unix)
     
     - List global (extern) functions in a file (prefixed with T):
          nm -g {{path/to/file.o}}

     - List only undefined symbols in a file:
          nm -u {{path/to/file.o}}

     - List all symbols, even debugging symbols:
          nm -a {{path/to/file.o}}

     - Demangle C++ symbols (make them readable):
          nm --demangle {{path/to/file.o}}
       
       Example:
       amit@amit-vm-mint21:~/C_programs$ nm ex1
          000000000000038c r __abi_tag
          0000000000004010 B __bss_start
          0000000000004010 b completed.0
                 w __cxa_finalize@GLIBC_2.2.5
          0000000000004000 D __data_start
          ...

      Symbol Type	Description
        A	Global absolute symbol.
        a	Local absolute symbol.
        B	Global bss symbol.
        b	Local bss symbol.
        D	Global data symbol.
        d	Local data symbol.
        f	Source file name symbol.
        L	Global thread-local symbol (TLS).
        l	Static thread-local symbol (TLS).
        R	Global read only symbol.
        r	Local read only symbol.
        T	Global text symbol.
        t	Local text symbol.
        U	Undefined symbol.



\end{verbatim}

\subsection{readelf  : ELF file info.}
\begin{verbatim}
       display information about ELF files ONLY
       including their metadata
       readelf [-a|--all]
               [-h|--file-header]
               [-l|--program-headers|--segments]
               [-S|--section-headers|--sections]
               [-g|--section-groups]
               [-t|--section-details]
               [-e|--headers]
               [-s|--syms|--symbols]

       
\end{verbatim}


\subsection{od  : "Octal Damp" - display binary file.}
\begin{verbatim}
       od is a command on various operating systems 
       for displaying ("dumping") data in various human-readable 
       output formats. The name is an acronym for "octal dump" 
       since it defaults to printing in the octal data format.

       Display file contents in octal, decimal or hexadecimal          
       format. Optionally display the byte offsets and/or  
       printable representation for each line.

       Default settings: Octal format, 8 bytes per line, 
                         ------------
       byte offsets in octal, and duplicate lines replaced with *:
             $ od tex_file_data_extract 
       
         -tu1 - decimal format
         -tx1 - hexadecimal format
         -N   - limit dump to N bytes:
         -v   - do not use * to mark line duplicates
            
       Decimal:
          $ od -tu1 a.txt
       Hexadecimal: 
        - Display file in hexadecimal format (1-byte units), 
          and 4 bytes per line:
          $ od --format={{x1}} --width={{4}} -v {{path/to/file}}
          $ od -tx1 tex_file_data_extract -N 1024          
      
      Read text file with od to see ASCII codes : 
          $ echo "ABCD" > a.txt
          $ cat a.txt
            ABCD
          $ od -tu1 a.txt # decimal ASCII values
            0000000  65  66  67  68  10
            0000005
          $ od -tx1 a.txt # 1 Byte (8 bit) hexadecimal ASCII values
            0000000 41 42 43 44 0a
            0000005 
            Note : 0x0a is for line feed character 
          $ od -tx4 a.txt # 4 BYTE (32 BIT) hexadecimal ASCII values
            0000000 44434241 0000000a
            0000005

\end{verbatim}

\section{Editors, text operations}

\subsection{nano  : run text editor.}
\begin{verbatim}
       Run text editor
       nano ex1.c

       
\end{verbatim}

\subsection{vim      :v improved editor.}
\begin{verbatim}
     vim (v improved) , to run :
         vim 1.txt 
 
     2 modes of work : 
         • Edit (:i to enter insert mode),  
           i.e. type ":" and after it "i"
         • Commands (esc to enter command mode)
            > Shift +:e [file] – Opens a [file] that you want
              to open. Here [file] is the filename that you want to open
            > Esc, :w    - Save the file but do not exit.
            > Esc, :q!   - To quit the file without first
              saving what you were working on.
            > Esc, :wq   - To save the file and exit from Vim.
            > Esc, :help - To get help
 
\end{verbatim} 


\subsection{tilde      : intuitive text editor.}
\begin{verbatim}
     tilde - an intuitive text editor
     for the console/terminal
 
\end{verbatim} 

\section{Install applications}


\subsection{dpkg   : Debian PaKaGe install}
\begin{verbatim}

     Command to install debian package (adds it to apt)

    • -i OR --install: Install a package using the dpkg command.
     The command will extract all control files for the specified
     package, remove any previously installed older 
     instance of the package, and install the new package on our system.
    • -r OR --remove: Remove an installed package 
     from our system. It removes every file belonging to the 
     specific package except the configuration files. 
     This can be seen as the uninstallation option.
    • -P OR --purge: An alternative way to remove an in
     stalled package from our system. It completely removes
     every fie belonging to the specific package, including
     the configuration files. This can be seen as the 
     ‘complete uninstallation’ option.
    • -l OR --list: Lists all installed packages on your system.
    • -L OR --listfiles: List installed package’s file locations.
    • -s OR --status: Shows the status of an installed package.
    • -S OR --search: Search for a pattern in installed packages.

     Example of VS Code install: 
     $ sudo dpkg -i code_1.81.1-1691620686_amd64.deb
     
     DPKG – s command : shows the status of an installed package: 
     $ dpkg -s code     
     
     sudo dpkg -i *.deb

     sudo dpkg –install *.deb 

     see APT to find differences between APT and DPKG

\end{verbatim}




\subsection{apt   : "Advanced Package Tool"}
\begin{verbatim}
        
        https://en.wikipedia.org/wiki/APT_(software)       
       
          APT provides a high-level commandline interface
       for the package management system. 
     
     - Update the list of available packages and versions 
       (it's recommended to run this before other apt commands):
          sudo apt update

     - Search for a given package:
          apt search {{package}}
          Example : apt search code # for Visual Studio
     - Show information for a package:
          apt show {{package}}
          Example : apt show code   # for Visual Studio
     - Install a package, or update it to the latest available version:
          sudo apt install {{package}}

     - Remove a package (using purge instead also removes its configuration files):
          sudo apt remove {{package}}

     - Upgrade all installed packages to their newest available versions:
          sudo apt upgrade

     - List all packages:
          apt list

     - List installed packages:
          apt list --installed



     *********Apt vs. Dpkg*********************************
     https://www.geekbits.io/what-is-the-difference-between-apt-and-dpkg/
     
     Apt (Advanced Package Tool) is a newer package manager
     that offers more features than dpkg. It can 
     automatically handle dependencies, so you don't have to
     worry about manually installing them. 
     Apt also has a wider range of repositories from which
     you can install packages.

     Dpkg (Debian Package Management System) is the older
     package manager. It doesn't have as many features as apt,
     but it's still a powerful tool. 
     Dpkg can't automatically handle dependencies, 
     so you'll need to install them manually. 
     Dpkg also has a limited range of repositories, 
     but you can still find the most common packages.


     Apt is a higher-level tool that provides a more straightforward
     interface to work with than dpkg. 
     While dpkg can only install, remove, or query individual packages,
     apt offers additional functionality, such as automatically 
     handling dependencies,retrieving and installing packages from
     remote locations, and upgrading all installed packages to
     their latest versions.

     Additionally, apt offers several useful commands
     for managing repositories and upgrading the system.
     For example, the apt-get command can install or remove packages,
     while the apt-cache command can search for available packages.

     Another difference is that dpkg is primarily a command-line tool,
     while apt also provides a graphical user interface (GUI)
     in the form of the aptitude tool.

     Which Package manager should you use?
     So, which package manager should you use? It depends on your needs.
     If you want a more powerful tool with automatic dependency management,
     then apt is the way to go. 
     On the other hand, if you want a more straightforward package manager
     or if you need to install packages from a specific repository,
     then dpkg is the better choice.

     Overall, apt (Recommended tool) is a more user-friendly and robust
     option than dpkg, making it the preferred package manager 
     for most users. It provides a simpler and more automated way to
     manage packages on your system. 
     
     However, dpkg can still be helpful in certain situations,
     such as when working with old or custom packages 
     that are unavailable through apt.
\end{verbatim}
            
\section{GCC and binary files operations}


\section{Ethernet }


\subsection{Wireshark install}
\begin{verbatim}

     https://askubuntu.com/questions/700712/how-to-install-wireshark
     
      1.Add the stable official PPA. To do this, go to terminal by pressing Ctrl+Alt+T
      and run:

          sudo add-apt-repository ppa:wireshark-dev/stable

      2.Update the repository:

          sudo apt-get update

      3.Install wireshark 2.0:

          sudo apt-get install wireshark

      4.Run wireshark:

          sudo wireshark
          
      5.If you get an error 
      "couldn't run /usr/bin/dumpcap in child process: Permission Denied", go to 
      the terminal again and run:

          sudo dpkg-reconfigure wireshark-common
          
      Say YES to the message box. This adds a wireshark group. Then add user to the
      group by typing

          sudo adduser $USER wireshark
          
      Then restart your machine and open Wireshark. It works. Good Luck.
     
\end{verbatim}     


\subsection{packet.names.at.different.levels.of.TCP/IP}
\begin{verbatim}
         
         packet names at different levels of TCP/IP
         
         application  |  message
         transport    |  segment
         network      |  datagram
         link         |  frame
\end{verbatim}

\subsection{TCP.vs.UDP}
\begin{verbatim}
        TCP is Transmission Control Protocol
        UDP is User Datagram Protocol
        
        TCP use processes to check that the packets are correctly sent
        to the receiver.
        
        UDP send packets without any checks but with the advantage
        of being quicker than TCP. 

\end{verbatim}



\subsection{sysctl : change kernel Networking params}
\begin{verbatim}
        All TCP/IP tuning parameters are located under /proc/sys/net/
        
        important tuning parameters:
        /proc/sys/net/core/rmem_max - Maximum TCP Receive Window
        /proc/sys/net/core/wmem_max - Maximum TCP Send Window
        /proc/sys/net/ipv4/tcp_rmem - memory reserved for TCP receive buffers
        /proc/sys/net/ipv4/tcp_wmem - memory reserved for TCP send buffers
        
        we dont change the above parameters directly,
        sysctl command is used to modify kernel parameters at runtime.
        The parameters available are those listed under /proc/sys/.
        
        Example for modifying kernel receiver buffer to 16Mbyte:
           $ sysctl -q -w net.core.rmem_max = 16777216
           $ sysctl -q -w net.core.rmem_default=16777216
           
           -q : quiet
           -w : write
           
\end{verbatim}




\subsection{IP..netmask..explanation}
\begin{verbatim}
       
       Four basic elements are combined to make an Ethernet system
          1. The FRAME, which is a standardized set of bits used to carry             
             data over the system.

          2. The MEDIA ACCESS CONTROL protocol, which defines how multiple 
             computers access the shared Ethernet channel in a fair manner.

          3. The SIGNALING COMPONENTS, which consists of standardized 
             electronic devices that send and receive signals over an 
             Ethernet channel.

           4. The PHYSICAL MEDIUM, which consists of the cables and other 
              hardware used to carry the digital Ethernet signals between 
              computers attached to the network.
       
       The Ethernet frame (MAC Frame)
       ------------------------------
       7 bytes  Preamble (for synchronizing)
       1 byte   SFD      (Start Frame Delimeter)
       6 bytes  DESTINATION ADDRESS (MAC address)
       6 bytes  SOURCE ADDRESS
       2 bytes  LENGTH/TYPE (frame type or frame length)
       
       Data field:
       46 TO 1500 or 1504  DSAP and SSAP (Destination and Source
       or 1982 bytes       Service Access Point)+CTRL (1 byte each)
      

        4 bytes FRAME CHECK SEQUENCE (CRC 32)
       
       
       
       The Media Access Control protocol
       ---------------------------------
       MAC protocol defines how the physical medium is shared among
       stations connected.
       
       Defines how multiple computers access the shared Ethernet 
       channel in a fair manner.
       
       Ethernet-equipped computer operates independently of all other 
       stations on the network; there is no central controller
       
       Ethernet uses a broadcast delivery mechanism, in which each frame
       that is transmitted is heard by every station.
       
       Every frame sent has an Destination Address and Source Address
       
       MAC address (true HW address, used for switching):
       ---------------------------------------------       
       MAC refers to a unique hardware address assigned to network devices  
       for identification at the data link layer.
       Every Ethernet Station has a unique 48 bit MAC Address (6 bytes).
       
       IP address (can be changed, used for routing):
       ----------------------------------------------
       Two computers with different MAC addresses can
       have the same IP address . Also for example, in case when we change
       HW adapter in the same computer 
       
       This use of m.a.c. address rise some problems : 
        1. there are network hw adapters that allow to put different m.a.c. 
           address.
        2. In the case when in the network one of network adapters is damaged 
           and replaced by new one (for ex. fix the adaptor problem). 
        That is why defined IP (Internet Protocol), and IP address that was 
        defined by Internet Protocol. There are IP adresses 
        ver. 4 and ver. 6 (4 is more common).
        
        Internet protocol ---> IP ----> IP address  
        
        IP address consists of two parts : (subnet number)&(Host number)       
        
        10.0.0.100 --- 10.0.0.101----|         |----10.0.1.100 ---- 
            |                        |         |                  |
            |                      10.0.0.1--10.0.1.1             |
            |                        |GW     GW|                  |
        10.0.0.102 ----10.0.0.103----|         |----10.0.1.101-----  
          Subnet mask : 255.255.255.0               255.255.255.0              
                  10.0.0.*                           10.0.1.*
         
        GW : Gate Way (Router)        
        
        IP address ver. 4 consists of 4 bytes:
        
        0-255 | 0-255 | 0.255 | 0-255
          192 . 168   . 0     . 1
        
             
        Subnet mask number 
        ------------------- 
        netmask/subnet number helps to computer to understand if the 
        package will be send inside the subnet or directly to GW router
        to pass the message outside of subnet.         

              
        We would like to identify between local networks. For this purpose  
        network mask (netmask) is defind.
        We tell that some bits of the IP number present the network number - 
        netmask number is used for this purpose.
        The bits in netmask number that are equal to 1 define area field of  
        network number.
        The bits in netmask number that are equal to 0 define area  field of 
        computer number within local network. 
        If the computer 10.0.0.100 wants to send message to computer 
        10.0.0.103 - it knows that it must send the message in local network 
        10.0.0 and there is no need to pass the message out of the local 
        network 10.0.0. Enough that 10.0.0.100 will send broadcast to all the 
        computers in its local network 10.0.0...until it will get the answer 
        from 10.0.0.103. 
        But if the computer 10.0.0.100 will try to send the message to       
        10.0.1.101...Because it knows that network 10.0.1 is not its own       
        network, it will send the request to Gate Way computer 10.0.0.1.
        Gate Way computer 10.0.0.1 sends the message to Gate Way computer  
        10.0.1.1 and the Gate Way of 10.0.1 will send broadcast on its local 
        network 10.0.1.
        After number of hopes the message from 1.0.0.100 will arrive to 
        computer 1.0.1.101 At home network we connected by WiFi to router 
        that si the GW of local home newtwork. The router is connected to 
        Internet provider... and so on , until the message arrives its 
        target.
        In the following example the network number can be 24 bit
        and the computer number can be 8 bits.
        
        10  .  0   .  0   . 100                  IP address
        255 . 255  .  255 .  0                   Subnet mask
          Network number  | Computer Number      255.255.255.0 <--> /24
          (subnet mask)   | (Host number)
        ------------------------------------        
                          | Num. of computers that can be connected: 
                          |          2^8 =256
                          | 255 is for brodcast
                          | 1   is for GW

           10  .  0   .  0   . 100                  IP address
           255 . 255  .  0   .  0                   Subnet mask 
       Network number | Computer Number             255.255.0.0 <--> /16
       (subnet mask)  | (Host number)       
       -------------------------------------        
                      | Num. of computers that can be connected:
                      |             2^16 = 65000
                      | 255 is for brodcast
                      | 1   is for GW

               10  .  0   .  0   . 100                IP address
               255 . 240  .  0   .  0                 Subnet mask 
       Network number | Computer Number               255.240.0.0 <--> /12
            ---------------------------------        
                      | Num. of computers: 2^20 = 1048576
                      | 255 is for brodcast
                      | 1   is for GW
        240d=11110000b 


        Loopback channel
        ---------------- 
        Intended for inter process communication - IPC:      
        In some computers 127.0.0.1 is the same computer (for loopback 
        perform), in other computers it is address 0.0.0.0 (for loopback 
        perform).We can take 2 programs in the same computer, without common 
        memory and send messages from one program to another (inter process 
        communication - IPC).  


       IF one of computers has the same IP than another computer in the local 
       network, we can change of IP and netmask of one of computers with 
       ifconfig command with following options : 
         $ ifconfig eth0 192.168.0.102               # only IP number chagne
         $ ifconfig eth0 192.168.0.102 255.255.255.0 # IP + Mask
         $ ifconfig eth0 192.168.0.102 /24 # define how many bits takes 
                                           # Network address  
       An additional example: 
         $ ifconfig wifi0 10.0.0.100 255.255.0.0
         $ ifconfig wifi0 10.0.0.100 /16   # define how many bits takes
                                           # Network address     
\end{verbatim}

\subsection{DHCP   : Dynamic Host Configuration Protocol}
\begin{verbatim}
     
    https://en.wikipedia.org/wiki/Dynamic_Host_Configuration_Protocol     
     
      DHCP server runs within router (Gate Way). 
      DHCP client runs within computer in the same subnet as router.
      
      DHCP server manage internal table with all computers in subnet:
      Leas time helps to manage time given to computer to have some IP. 
      
        MAC addr.           IP addr       Time
        a1-a1-a1-b1-c2-d4   10.0.0.100    ~~~~~~   
        b2-c4-d2-12-28-aa   10.0.0.101    ~~~~~~
     
     We can use DHCP (Dynamic Host Configuration Protocol) to define IP 
     numbers automatically.
     
     DHCP stands for Dynamic Host Configuration Protocol. 
     
     It is a protocol that allows devices on a network (DHCP clients) 
     to get an IP address and other network settings automatically from a 
     DHCP server. This makes it easier to connect multiple devices to a 
     network without manually configuring each one.  
     
     The DHCP protocol also provides a mechanism whereby a client can learn 
     important details about the network to which it is attached, such as 
     the location of a default router (GW), the location of a name server, 
     and so on.  
     
     Some of the benefits of using DHCP are:
    • It reduces the chance of IP address conflicts, where two 
       devices have the same IP address.
    • It simplifies network management, as the DHCP server can
       assign and reclaim IP addresses centrally.
    • It supports mobility, as devices can move from one network
       to another and get a new IP address automatically.  
     
    Some of the drawbacks of using DHCP are:
    • It depends on the availability and reliability of the DHCP server. If 
      the server is down or unreachable, devices cannot get an IP address or 
      renew their lease.
    • It may introduce security risks, as malicious users can spoof DHCP 
      messages or set up rogue DHCP servers to intercept or redirect network 
      traffic. 
    • It may not be suitable for some devices that need a fixed or static IP 
      address, such as servers or printers.
     
      Router usually is used as DHCP server. 
      He is also Gate Way. 
      It's address is usually constant : 192.168.0.1. 
      Operation systems are familiar with such addresses.
      
      Computer that is turned on for the first time in the network, send 
      messages to check which computer is GateWay. 
      After it finds it , the computer ask the G.W. router to get IP address 
      (this operation is calles IP lease).
      The router manages internally the table of mac address, IP , expiration 
      date of IP , etc.
      --------------------------------------------------------
                    Internet                       A  | ... . 100
                       |                           B  | ... . 101
                router(DHCP Server)  192.168.0.1   C  | ... . 102
                       |                           D  | ... . 103
              --------------------------       
             |         |        |      |
             A         B        C      D
      192.168.0.100
      
      Now computer D is disconnected and in place of D connected 
      new computer E. If address of D is not expired, router will set 
      addrerss of E ...104, but if address of D is expired, 
      then the router can set the address of E to be ...103.  
      
                    Internet                       A  | ... . 100
                       |                           B  | ... . 101
                router(DHCP Server) 192.168.0.1    C  | ... . 102
                       |                           D  | ... . 103
              ------------------------------       E  | --->  104
             |         |        |          |  
             |         |        |          |      or if 103 expired :
             A         B        C      D   E       A  | ... . 100
      192.168.0.100                                B  | ... . 101
                                                   C  | ... . 102
                                                   D  | xxx   xxx
                                                   E  | ... . 103
        
        But now the computer D connected again and D asks from 
        DHCP server to reftesh its IP address. Now DHCP server will set its 
        address to ...104:
  
                    Internet                       A  | ... . 100
                       |                           B  | ... . 101
                router(DHCP Server) 192.168.0.1    C  | ... . 102
                       |                           D  | ---> (104)
              ------------------------------       E  | ... . 103
             |         |        |          |  
             A         B        C      D   E       
      192.168.0.100                                 
                                                    
        If the address ...103 in any case is important for D,
        it can ask E computer to release its address, to allow ifconfig, in 
        order not to  make collision with existing IP.
        E is calling the router and ask him to release the IP address.
        In this case D can also make release to get the address.
        Additional option is to define IP address per m.a.c. address
        within router (i.e. edit routing table within the router).                                            
                     
      --------------------------------------------------------
                  DHCP protocol
         
        CLIENT          SERVER

              discover
          --------------> (1)
            
              offer
          <-------------- (2)  
   
              request
          --------------> (3)
    
            acknowledge
          <-------------- (4)   
   
         DHCP operations fall into four phases: 
              server discovery, 
              IP lease offer, 
              IP lease request, 
              IP lease acknowledgement. 
         These stages are often abbreviated as DORA for 
         Discovery, Offer, Request, and Acknowledgement.
   
      Discovery
      The DHCP client broadcasts a DHCPDISCOVER message on the network subnet             
      using the destination address 255.255.255.255 (limited broadcast) or 
      the specific subnet broadcast address (directed broadcast). A DHCP 
      client may also request an IP address in the DHCPDISCOVER, which the 
      server may take into account when selecting an address to offer.

      Offer 
      When a DHCP server receives a DHCPDISCOVER message from a client, which       
      is an IP address lease request, the DHCP server reserves an IP address 
      for the client and makes a lease offer by sending a DHCPOFFER message 
      to the client. This message contains the client's client id 
      (traditionally a MAC address), the IP address that the server is 
      offering, the subnet mask, the lease duration, and the IP address of 
      the DHCP server making the offer.

      Request
      In response to the DHCP offer, the client replies with a DHCPREQUEST       
      message, broadcast to the server, requesting the offered address. A 
      client can receive DHCP offers from multiple servers, but it will 
      accept only one DHCP offer. 

      Acknowledgement
      When the DHCP server receives the DHCPREQUEST message from the client,       
      the configuration process enters its final phase. The acknowledgement 
      phase involves sending a DHCPACK packet to the client. This packet 
      includes the lease duration and any other configuration information 
      that the client might have requested. At this point, the IP 
      configuration process is completed. 
      
\end{verbatim}

\subsection{dhclient : obtain/release the IP address drom DHCP server}
\begin{verbatim}
      DHCP server runs within router (Gate Way). 
      DHCP client runs within computer in the same subnet as router.
      
      - Get an IP address for the eth0 interface:
           $ sudo dhclient {{eth0}}

      - Release an IP address for the eth0 interface:
           $ sudo dhclient -r {{eth0}}

      To obtain fresh IP address from DHCP server (router (Gate Way)):
         $ sudo dhclient enp0s3

      To release IP address within DHCP server:
         $ sudo dhclient -r enp0s3       
      Ones current lease has been released, the client exits.      
      
      
      Example:
      $ ip address
        2: enp0s3: <BROADCAST,MULTICAST,UP,LOWER_UP> mtu 1500 qdisc 
        fq_codel state UP group default qlen 1000
        link/ether 08:00:27:78:7a:16 brd ff:ff:ff:ff:ff:ff
        inet 10.0.2.15/24 brd 10.0.2.255 scope global dynamic 
      $ sudo dhclient -r enp0s3
      $ ip address
        2: enp0s3: <BROADCAST,MULTICAST,UP,LOWER_UP> mtu 1500 qdisc
        fq_codel state UP group default qlen 1000
        link/ether 08:00:27:78:7a:16 brd ff:ff:ff:ff:ff:ff
        inet6 fe80::e013:b6a6:6af3:9b07/64 scope link noprefixroute 
      $ sudo dhclient enp0s3
      $ ip address
        2: enp0s3: <BROADCAST,MULTICAST,UP,LOWER_UP> mtu 1500 qdisc
        fq_codel state UP group default qlen 1000
        link/ether 08:00:27:78:7a:16 brd ff:ff:ff:ff:ff:ff
        inet 10.0.2.15/24 brd 10.0.2.255 scope global dynamic  


\end{verbatim}


\subsection{SSH : OpenSSH remote login client}
\begin{verbatim}
     
     ssh (SSH client) is a program for logging into a remote 
     machine and for executing commands
     on a remote machine.  It is intended to provide secure 
     encrypted communications between two
     untrusted hosts over an insecure network.
     
     
     Used to access a remote computer on most Linux/Unix systems
     (write shell commands just as if you were sitting at the 
     workstations) just like telnet.
     
     SSH is Secure SHell (connection to remote computer), 
     SCP is Secure CoPy (copy to remote computer).
     SSH allows connect another Linux computer in safety manner.
     SCP allows copy files from one computer to another computer 
     in safety manner.
     
     1. check if ssh exists on the computer:
         $ ssh -V     
          OpenSSH_8.9p1 Ubuntu-3ubuntu0.6, OpenSSL 3.0.2 15 Mar 2022
     
     2. If not exists , install SSH server (daemon): 
         $ sudo apt update        
         $ sudo apt install openssh-server
     
     3. run SSH server on both computers: 
         $ sudo systemctl start ssh   #start ssh for this session 
            or
         $ sudo systemctl enable ssh  #start ssh from boot to computer  
     4. connect to computer : 
         $ ssh user_name@computer_name
            or
         $ ssh user_name@computer_ip_address  
         
         Example : 
         $ ssh amit@192.168.1.12
          The authenticity of host '192.168.1.12 (192.168.1.12)' 
          can't be established.
          ED25519 key fingerprint is SHA256:P5+OF1k8R/kSq4BSPgypkbda/
          o4XhXsdgHjVtvhuz5c.
          This key is not known by any other names
          Are you sure you want to continue connecting 
          (yes/no/[fingerprint])? yes
          Warning: Permanently added '192.168.1.12' (ED25519) 
          to the list of known hosts.
          amit@192.168.1.12's password: 
          Last login: Sun Oct 29 06:17:17 2023 from 192.168.1.18
         
          amit@amit-vm-mint21:~$ 
          amit@amit-vm-mint21:~$ 
          amit@amit-vm-mint21:~$ 
          
          Note1 : $USER NAME is amit,
                  $HOSTNAME is amit-vm-mint21
                  
          Note2 : SSH uses computer fingerprint to recognize the remote 
                  computer.     
     
     There are two commands to operate SSH : systemctl and service
     SSH operate commands: 
        $ sudo systemctl status ssh  #check the status
        $ sudo systemctl stop ssh    #stop ssh for this session
        $ sudo systemctl start ssh   #start ssh for this session
        $ sudo systemctl disable ssh #stop ssh from boot to computer
        $ sudo systemctl enable ssh  #start ssh from boot to computer    
     

     $ sudo systemctl status ssh
        [sudo] password for amits: 
        o ssh.service - OpenBSD Secure Shell server
        Loaded: loaded (/lib/systemd/system/ssh.service; disabled;
         vendor preset: enabled)
        Active: inactive (dead)
        Docs: man:sshd(8)
              man:sshd_config(5)

      Virtual box settings to access external computer network : 
      ------------------------------------------------------------      
      To connect virtual machine to external network we must choose
      bridge adapter (it connects virtual machine to home network).
      Choose: 
          Settings, 
          Network, 
          Attached to : Bridged adapter
          Name : MedaTek Wifi 6 M7921 Wireless LAN Card  
      Note: you must be connected to Internet in outer computer !     
     
      Check if network connection exists: 
        $ ifconfig
        enp0s3: flags=4163<UP,BROADCAST,RUNNING,MULTICAST>  mtu 1500
        inet 192.168.1.12  netmask 255.255.255.0  broadcast 192.168.1.255
      ------------------------------------------------------------
      Note: NAT network in virtual box is internal network between
      different machines inside the virtual box. If we define NUT network for
      two machines, we can run ssh and scp commands on both virtual computers.
      May there is needed to run 
          $ dhclient enp0s3 # in order to get IP address  
          or
          
          $ ip address add
     
\end{verbatim}

\subsection{fail2ban :  security intrusion prevention software framework}
\begin{verbatim}

      https://en.wikipedia.org/wiki/Fail2ban

      Fail2ban is an intrusion prevention software framework. 
      Written in the Python programming language, it is designed to prevent   
      brute-force attacks.    
\end{verbatim}




\subsection{who : display who is logged in}
\begin{verbatim}
       - Display the username, line, and time of all currently 
         logged-in sessions:
           $ who
         Example: 
         $ who
           amit     tty7         2024-03-10 07:17 (:0)
           amit     pts/1        2024-03-10 17:37 (192.168.1.18)

       
\end{verbatim}




\subsection{SCP : OpenSSH secure file copy}
\begin{verbatim}
     
     scp - remote file copy program, copies files between hosts on 
     a network
     
     SCP is Secure CoPy (copy to remote computer).
     SCP allows copy files from one computer to another computer 
     in safety manner.
     
     uses ssh for data transfer authentication and provides
     the same security as SSH does.
     
     Copies between two remote hosts are permitted also !
     
     1. check if ssh exists on the computer:
         $ ssh -V     
          OpenSSH_8.9p1 Ubuntu-3ubuntu0.6, OpenSSL 3.0.2 15 Mar 2022
     
     2. If not exists , install SSH server (daemon): 
         $ sudo apt update        
         $ sudo apt install openssh-server
     
     3. run SSH server on both computers: 
         $ sudo systemctl start ssh   #start ssh for this session 
            or
         $ sudo systemctl enable ssh  #start ssh from boot to computer  
     4. Copy a local file to a remote host: 
 scp {{path/to/local_file}} {{remote_host}}:{{path/to/remote_file}}
        Copy a file from a remote host to a local directory:
 scp {{remote_host}}:{{path/to/remote_file}} {{path/to/
                                                  local_directory}}
        Recursively copy the contents of a directory from a remote 
        host to a local directory:
 scp -r {{remote_host}}:{{path/to/remote_directory}} {{path/to/                        
                                                  local_directory}}                                         

     Example:
     
     $ scp p_data.txt amit@192.168.1.12:/home/amit/Documents
       amit@192.168.1.12's password: 
       p_data.txt              100%  648   126.7KB/s   00:00    
         
     Recursive copy of folders:
     $ scp -r /home/amits/Documents/RTC    
                       amit@192.168.1.12:/home/amit/Documents/RTC
     
     Windows :
     $ pscp.exe bennyc@192.168.42.38:/home/minicom.log c:


     
\end{verbatim}

\subsection{pscp : transfer files between Windows and linux}
\begin{verbatim}

       The open source PSCP utility makes it easy to transfer files and folders 
       between Windows and Linux computers

       Using PSCP
          PSCP (PuTTY Secure Copy Protocol) is a command-line tool for transferring     
          files and folders from a Windows computer to a Linux computer.

       Download pscp.exe from its website.
       https://www.chiark.greenend.org.uk/~sgtatham/putty/latest.html

          Move pscp.exe to a folder in your PATH (for example, Desktop\App if you 
          followed the PATH tutorial here on Opensource.com). If you haven't set a PATH 
          variable for yourself, you can alternately move pscp.exe to the folder 
          holding the files you're going to transfer.

          Open Powershell on your Windows computer using the search bar in the Windows 
          taskbar (type 'powershell` into the search bar.)

          Type pscp –version to confirm that your computer can find the command.












       Using PSCP

       PSCP (PuTTY Secure Copy Protocol) is a command-line tool for transferring files  
       and folders from a Windows computer to a Linux computer.

       Download pscp.exe from its website.

       Move pscp.exe to a folder in your PATH (for example, Desktop\App if you followed   
       the PATH tutorial here on Opensource.com). If you haven't set a PATH variable 
       for yourself, you can alternately move pscp.exe to the folder holding the files 
       you're going to transfer.

       Open Powershell on your Windows computer using the search bar in the Windows 
       taskbar (type 'powershell` into the search bar.)

       Type pscp –version to confirm that your computer can find the command.
       

\end{verbatim}

\subsection{ping : check connection}
\begin{verbatim}
       This is convinient way to check that the opposite side
       is connected and supports ping protocol.
       
       It sends packets to given machine and get ackonowlegment 
       packets in return.
       
       If you can ping GW (router) - your network interface is OK
       If you is able to ping an external IP address, your network settings 
       are correct.
       
       ping also evaluates time and hops that takes to 
       the command to arrive the opposite side.
       
       Ask ping to check google dns site 3 times only : 
       $ ping -c 3 8.8.8.8
         PING 8.8.8.8 (8.8.8.8) 56(84) bytes of data.
         64 bytes from 8.8.8.8: icmp_seq=1 ttl=119 time=12.2 ms
         64 bytes from 8.8.8.8: icmp_seq=2 ttl=119 time=13.5 ms
         64 bytes from 8.8.8.8: icmp_seq=3 ttl=119 time=13.8 ms

--- 8.8.8.8 ping statistics ---
3 packets transmitted, 3 received, 0% packet loss, time 2003ms
rtt min/avg/max/mdev = 12.215/13.174/13.791/0.687 ms


          --- 8.8.8.8 ping statistics ---
          8 packets transmitted, 8 received, 0% packet loss, time 7008ms
          rtt min/avg/max/mdev = 12.861/14.049/15.676/0.881 ms

          TTL : how many hops it can made maximum !!!
          
          TTL : time to leave
          
          https://www.cloudflare.com/en-gb/learning/cdn/glossary
          /time-to-live-ttl/  
          Time to live (TTL) refers to the amount of time or “hops” 
          that a packet is SET to exist inside a network BEFORE being
          DISCARDED BY A ROUTER.
          
          Time To Leave:
          How many jumps(hops) the packet maximum will pass 
          until the next router will throw it.
       
          to set TTL:
             $ ping rt-ed.co.il -t 50 # set TTL=50  
           
          
          In following picture : rectangular is computer and circle is 
          router. We want to send message from one computer to another. 
          Firstly the message arrives to 1-st router, it must to calculate   
          the best rout from it to second computer. 
          In the well case the data arrives the destination :
          
          []---> O ----> O ----> O ----> []      
                    
          In case of fault the packet can enter the infinite loop between 
          routers. It can cause traffic jam. In order to prevent such 
          situation Internet uses TTL - time to live, how time the packet
          can hop maximum between different routers.
        
          []---> O <---- O xxx O xxx []      
                 |       ^
                 |       | 
                 v       |
                 O-----> O

          When some packet is in closed loop it can cause to conjesction it
          the net, when other packets arrive. 
          To prevent conjesction , defined TTL.
\end{verbatim}


\subsection{iperf : measure quality and bandwidth of link}
\begin{verbatim}
         iperf is a tool to measure the bandwidth and the quality of a 
         network link.
         The network link is delimited by two hosts running iperf.
         The quality of a link can be tested as follows:
           - Latency (response time or RTT): can be measured with the Ping  
              command.
           - Jitter (latency variation): can be measured with an 
              iperf UDP test.
           - Datagram loss: can be measured with an iperf UDP test.
           
   On the first computer (with IP 192.168.1.18) we run the server : 
      $ iperf -s
   ------------------------------------------------------------
   Server listening on TCP port 5001
   TCP window size:  128 KByte (default)
   ------------------------------------------------------------
   [  1] local 192.168.1.18 port 5001 connected with 192.168.1.12 port 38950
   [ ID] Interval       Transfer     Bandwidth
   [  1] 0.0000-10.2928 sec  55.5 MBytes  45.2 Mbits/sec


   
   On the second computer we run the client , enter the IP address of
   the first computer : 
      $ iperf -c 192.168.1.18
   ------------------------------------------------------------
   Client connected to 192.168.1.18, TCP port 5001
   TCP window size:  85 KByte (default)
   ------------------------------------------------------------
   [  1] local 192.168.1.12 port 38950 connected with 192.168.1.18 port 5001
   [ ID] Interval       Transfer     Bandwidth
   [  1] 0.0000-10.2928 sec  55.5 MBytes  45.2 Mbits/sec        
    
    
   Send ctrl+C signal to stop the iperf server.
   
   To run the server in UDP mode : 
     $ iperf -s -u 
   
   Change Network parameters and run iperf server : 
     $ sysctl -q -w net.core.rmem_max=16777216
     $ sysctl -q -w net.core.rmem_default=16777216
     $ iperf -s –u   
   
   Run client in UDP aith defined size : 
   iperf -c 2.2.2.11 -u -b 2G
   
           
\end{verbatim}


\subsection{TTL : Time To Leave}
\begin{verbatim}
       In case of fault the packet can enter the infinite loop between 
          routers. It can cause traffic jam. In order to prevent such 
          situation Internet uses TTL - time to live, how time the packet
          can hop maximum between different routers.
        
          []---> O <---- O xxx O xxx []      
                 |       ^
                 |       | 
                 v       |
                 O-----> O
          When some packet is in closed loop it can cause to conjesction it
          the net, when other packets arrive. 
          To prevent conjesction , defined TTL.
       
       
       
       
       
       Time To Leave:
       How many jumps(hops) the packet maximum will pass 
       until the next router will throw it.
       
       to set TTL:
         $ ping rt-ed.co.il -t 50 # set TTL=50       
      
      https://www.cloudflare.com/en-gb/learning/cdn/glossary
      /time-to-live-ttl/  
      Time to live (TTL) refers to the amount of time or “hops” 
      that a packet is SET to exist inside a network BEFORE being
      DISCARDED BY A ROUTER. 
      
      check Google dns server : 
      $ ping 8.8.8.8
        PING 8.8.8.8 (8.8.8.8) 56(84) bytes of data.
        64 bytes from 8.8.8.8: icmp_seq=1 ttl=118 time=16.4 ms
        64 bytes from 8.8.8.8: icmp_seq=2 ttl=118 time=15.5 ms
        64 bytes from 8.8.8.8: icmp_seq=3 ttl=118 time=14.4 ms
        64 bytes from 8.8.8.8: icmp_seq=4 ttl=118 time=14.4 ms 
       
\end{verbatim}

\subsection{Gate Way : Time To Leave}
\begin{verbatim}
       GW is routher/computer that is used to connect us out 
       of our local network.     
       
\end{verbatim}



\subsection{traceroute : shows hopes(jumps) between devices}
\begin{verbatim}
       
       Print the route of packets trace to network host

     $ traceroute 192.168.1.12
     
traceroute to 192.168.1.12 (192.168.1.12), 30 hops max, 60 byte packets
 1  192.168.1.12 (192.168.1.12)  11.090 ms  11.044 ms  11.284 ms       
       
       $ traceroute rt-ed.co.il
       
traceroute to rt-ed.co.il (212.199.114.183), 30 hops max, 60 byte packets
 1  InnboxG84 (192.168.1.1)  6.118 ms  6.123 ms  6.462 ms
 2  100.64.32.1 (100.64.32.1)  9.977 ms  11.990 ms  12.110 ms
 3  45.80.88.197 (45.80.88.197)  12.803 ms  12.791 ms  12.779 ms
 4  82.102.142.225 (82.102.142.225)  36.986 ms  37.342 ms  37.331 ms
 5  212.116.177.49.static.012.net.il (212.116.177.49)  25.972 ms  25.961 ms 
    25.950 ms
 6  212.116.177.50.static.012.net.il (212.116.177.50)  20.347 ms  9.816 ms
    9.730 ms
 7  * * *
 8  * * *
 9  * * *
10  * * *

      $ traceroute google.com
      
traceroute to google.com (216.239.38.120), 30 hops max, 60 byte packets
 1  InnboxG84 (192.168.1.1)  3.762 ms  3.687 ms  3.808 ms
 2  100.64.32.1 (100.64.32.1)  10.539 ms  14.998 ms  15.167 ms
 3  45.80.88.197 (45.80.88.197)  18.678 ms  18.653 ms  18.617 ms
 4  72.14.221.92 (72.14.221.92)  18.128 ms  18.109 ms  17.911 ms
 5  * * *
 6  any-in-2678.1e100.net (216.239.38.120)  17.687 ms  11.035 ms  10.995 ms


\end{verbatim}

\subsection{route : set/display the routing table}
\begin{verbatim}

   Use route cmd to set the route table
   This command allows us to define the default Gate Way IP .

    - Display the information of route table:
      $ route -n
      or 
      $ route
    - Add route rule:
      $ sudo route add -net {{ip_addr}} netmask {{netmask_addr}} gw {{gw_addr}}

    - Delete route rule:
      $ sudo route del -net {{ip_addr}} netmask {{netmask_addr}} dev {{gw_addr}}
      
      $ route
Kernel IP routing table
Destination     Gateway         Genmask         Flags Metric Ref    Use Iface
default         InnboxG84       0.0.0.0         UG    600    0        0 wlp0s20f3
link-local      0.0.0.0         255.255.0.0     U     1000   0        0 wlp0s20f3
192.168.1.0     0.0.0.0         255.255.255.0   U     600    0        0 wlp0s20f3
   

      $ route -n
Kernel IP routing table
Destination     Gateway         Genmask         Flags Metric Ref    Use Iface
0.0.0.0         192.168.1.1     0.0.0.0         UG    600    0        0 wlp0s20f3
169.254.0.0     0.0.0.0         255.255.0.0     U     1000   0        0 wlp0s20f3
192.168.1.0     0.0.0.0         255.255.255.0   U     600    0        0 wlp0s20f3


      To define default GW : 

     $ sudo route add default gw 192.168.1.0
         SIOCADDRT: Network is unreachable



\end{verbatim}


\subsection{ip/ip route : ip command}
\begin{verbatim}
        
        ip - show / manipulate routing, network devices, 
             interfaces and tunnels
        
        ip [ OPTIONS ] OBJECT { COMMAND | help }
        
        OBJECT := { link | address | addrlabel | route | rule | neigh |           
                    table | tunnel | tuntap | maddress | mroute
                    | mrule | monitor | xfrm | netns | l2tp | tcp_metrics |  
                    token | macsec | vrf | mptcp | ioam }

        ------------------------------------------------------
         -Similar but newer than ifconfig command.
          display network status : 
        
            $ ifconfig
           or 
            $ ip address show # 
            $ ip addr show    # short form

         - Make an interface up/down:
            $ ip link set enp0s3 up   # enable enp0s3 interface
            $ ip link set enp0s3 down # disable enp0s3 interface     
           or 
            $ ifconfig enp0s3 up   # enable enp0s3 interface
            $ ifconfig enp0s3 down # disable enp0s3 interface           
       ------------------------------------------------------------    
         - Set IP address of eth0:
         
            $ ifconfig eth0 10.0.0.100
           or 
            $ ip address add 10.0.0.100 dev eth0

         - Change IP of eth0 to 10.0.0.100
           and netmask to 255.255.255.0
            $ ifconfig eth0 10.0.0.100 netmask 255.255.255.0
           or
            $ ifconfig eth0 10.0.0.100/24
           or
            $ ip address add 10.0.0.100/24 dev eth0         
       ------------------------------------------------------------    
         - Change default gateway to 10.0.0.1 for eth0 : 
            $ route add default gw 10.0.0.1
            
            $ ip route add default via 10.0.0.1 dev eth0
       ------------------------------------------------------------
        Show/manipulate routing, devices, policy routing and tunnels
        
         - List interfaces with detailed info:
             $ ip address
         
         1: lo: <LOOPBACK,UP,LOWER_UP> mtu 65536 qdisc noqueue state UNKNOWN           
            group default qlen 1000
            link/loopback 00:00:00:00:00:00 brd 00:00:00:00:00:00
            inet 127.0.0.1/8 scope host lo
            valid_lft forever preferred_lft forever
            inet6 ::1/128 scope host 
            valid_lft forever preferred_lft forever
         2: enp1s0: <NO-CARRIER,BROADCAST,MULTICAST,UP> mtu 1500 qdisc 
            fq_codel state DOWN group default qlen 1000
            link/ether 04:bf:1b:72:b9:5d brd ff:ff:ff:ff:ff:ff
         3: wlp0s20f3: <BROADCAST,MULTICAST,UP,LOWER_UP> mtu 1500 qdisc 
            noqueue state UP group default qlen 1000
            link/ether f4:6d:3f:8e:03:dc brd ff:ff:ff:ff:ff:ff
            inet 192.168.1.18/24 brd 192.168.1.255 scope global dynamic 
            noprefixroute wlp0s20f3
            valid_lft 74418sec preferred_lft 74418sec
            inet6 fe80::dad7:1350:f3b3:451d/64 scope link noprefixroute 
             valid_lft forever preferred_lft forever

         - Display the routing table:
             $ ip route
         - Show neighbors (ARP table):
             $ ip neighbour


            --------------------------------------
       
       This command can be used to display or modify the existing IP 
       routing table. We can add, delete, or modify specific static routes
       to specific hosts or networks using ip route command

      - Display the routing table:
          $ ip route {{show|list}}

      - Add a default route using gateway forwarding:
          $ sudo ip route add default via {{gateway_ip}}

      - Add a default route using eth0:
          $ sudo ip route add default dev {{eth0}}

      - Add a static route:
          $ sudo ip route add {{destination_ip}} via {{gateway_ip}} dev {{eth0}}

          Example : 
          
          $ sudo ip route add default via 192.168.7.3 dev enp1s0

      - Delete a static route:
          $ sudo ip route del {{destination_ip}} dev {{eth0}}

      - Change or replace a static route:
          $ sudo ip route {{change|replace}} {{destination_ip}} via {{gateway_ip}} dev {{eth0}}
          
          
          $ route -n
            Kernel IP routing table
   Destination     Gateway         Genmask         Flags Metric Ref    
   Use Iface
   0.0.0.0         192.168.1.1     0.0.0.0         UG    600    0        
   0  wlp0s20f3
   169.254.0.0     0.0.0.0         255.255.0.0     U     1000   0        
   0  wlp0s20f3
   192.168.1.0     0.0.0.0         255.255.255.0   U     600    0        
   0  wlp0s20f3
          $ ip route
   default via 192.168.1.1 dev wlp0s20f3 proto dhcp metric 600 
   169.254.0.0/16 dev wlp0s20f3 scope link metric 1000 
   192.168.1.0/24 dev wlp0s20f3 proto kernel scope link src 192.168.1.18
   metric 600 

    
\end{verbatim}


\subsection{DNS..server}
\begin{verbatim}
     https://www.cloudflare.com/en-gb/learning/dns
     /what-is-a-dns-server/
     
     Your programs need to know what IP address corresponds to a given 
     host name (ex. )
     Domain Name Servers (DNS) take care of this.

     You just have to specify the IP address of 1 or more DNS servers
     in your /etc/resolv.conf file:
       nameserver 217.19.192.132
       nameserver 212.27.32.177
     The changes takes effect immediately
     
     
     DNS (Domain Name Server) server is used to translate internet 
     site names (host name) to IP address.
     IP addresses of DNS servers are placed in /etc/resolv.conf file.
     
     To add Google DNS server see following :
     https://easyengine.io/tutorials/linux/google-public-dns/
     
     The Domain Name System (DNS) is the phonebook of the Internet.
     When users type domain names such as ‘google.com’ or 
     ‘nytimes.com’ into web browsers, DNS is responsible for finding
     the correct IP address for those sites. 
     
     Browsers then use those  IP addresses to communicate with
     origin servers or CDN edge servers to access website information.
     This all happens thanks to DNS servers: machines dedicated to
     answering DNS queries.  
   
         • Configure your network settings to use the IP addresses 8.8.8.8 
         and 8.8.4.4 as your DNS servers.

     check Google dns server : 
       $ ping 8.8.8.8
        PING 8.8.8.8 (8.8.8.8) 56(84) bytes of data.
        64 bytes from 8.8.8.8: icmp_seq=1 ttl=118 time=16.4 ms
        64 bytes from 8.8.8.8: icmp_seq=2 ttl=118 time=15.5 ms
        64 bytes from 8.8.8.8: icmp_seq=3 ttl=118 time=14.4 ms
        64 bytes from 8.8.8.8: icmp_seq=4 ttl=118 time=14.4 ms 
      
\end{verbatim}


\subsection{arp : manage Address Resolution Protocol}
\begin{verbatim}
       The arp command in Linux is a network utility tool used to display,    
       add, and remove entries in the Address Resolution Protocol (ARP) 
       cache. This command is crucial for managing network communications in 
       a Linux system
           
        - Show the current ARP table (mac addresses of every machine that with current machine ):
            arp -a

         display the ARP cache. 
         The output shows the IP address (192.168.1.1),
         the corresponding MAC address (ab:cd:ef:gh:ij:kl),
         the network interface (eth0) on which the ARP entry is located.


        - Delete a specific entry:
            arp -d {{address}}

        - Create an entry in the ARP table:
            arp -s {{address}} {{mac_address}}

-a	Show all entries in the ARP cache.	    arp -a
-d	Delete an entry in the ARP cache.	        arp -d 192.168.1.1
-s	Add a new entry to the ARP cache.	       arp -s 192.168.1.1 ab:cd:ef:gh:ij:kl
-v	Enable verbose mode.	                     arp -v -a
-n	Do not resolve names.	                 arp -n -a
-i	Specify the network interface.	         arp -i eth0 -a
-D	Read hardware address from given device.	arp -D eth0 -s 192.168.1.1
-e	Display in default (Linux) style.	         arp -e -a
-H	Specify the hardware type.	             arp -H ether -a
-p	Publish the entry.	                arp -s 192.168.1.1 ab:cd:ef:gh:ij:kl -p
-r	Read new entries from file or standard input.	      arp -r < arp_entries.txt
-t	Specify the hardware type.	            arp -t ether -a       
       
\end{verbatim}




\subsection{ifconfig : InerFace CONFIGurator }
\begin{verbatim}
     
     https://www.baeldung.com/linux/ifconfig-command
     
     
     ifconfig can be installed with command : 
     $ sudo apt install net-tools
     
     Connection names in Linux: 
     lo            ---> loopback
     
     (w)lps20 wifi0  ---> Wireless connection
     
     (e)nps20 eth0   ---> Ethernet cable connection 
     (e)nps21 eth1   
     
     Prints details about all the network interfaces available
     on your system.
     
     It is used to view and change the configuration of the network
     interfaces on your system.
     
     Display information about all network interfaces currently in operation:
       $ ifconfig
     
     To ENABLE an inactive interface, provide ifconfig with the interface
     name followed by the keyword up:
       $ sudo ifconfig eth1 up
     
     You can DISABLE an active network interface using the down keyword.
       $ sudo ifconfig eth1 down
       
     Configuring an interface
     ------------------------
     ifconfig can be used at the command line to configure 
     (or re-configure) a network interface. This is often unnecessary 
     since this configuration is often handled by a script when you
     boot the system. If you'd like to do so manually, you need superuser 
     privileges, so we'll use sudo again when running these commands. 
     
     To assign a static IP address to an interface, specify the interface 
     name and the IP address. 
     For example, to assign the IP address 192.168.7.100 to the interface 
     enp1s0, use the command:
       
       $ sudo ifconfig enp1s0 192.168.7.100
     
     To assign a network mask to an interface, use the keyword "netmask" and 
     the netmask address. For instance, to configure the interface enp1s0 to 
     use a network mask of 255.255.255.0, the command would be:     
       
       $ sudo ifconfig enp1s0 netmask 255.255.255.0
       $ sudo ifconfig enp1s0 10.0.0.100/16     # netmask 
                                                # of 16 bit for subnet
     
     These configurations can combined in a single command. For instance, to    
     configure interface enp1s0 to use the static IP address 192.168.7.100, 
     and the network mask 255.255.255.0:
     
       $ sudo ifconfig enp1s0 192.168.7.100 netmask 255.255.255.0
     
     mtu N	
     This parameter sets the MTU (maximum transfer unit) of an interface. 
     This setting is used to limit the maximum packet size transferred by the 
     interface. If you're unsure about it, you can safely leave this setting 
     alone.  
     
     netmask address	
     Set the IP network mask for this interface. This value defaults to the 
     usual class A, B or C network mask (as derived from the interface IP 
     address), but it can be set to any value.      
     ---------------------------------------------------------------
     
     What about DHCP (Dynamic Host Configuration Protocol)?
     ifconfig can only assign a static IP address to a network interface. To    
     assign a dynamic IP address using DHCP, use the "dhclient" command.
     
     Dynamic Host Configuration Protocol (DHCP) is a client/server protocol 
     that automatically provides an Internet Protocol (IP) host with its IP 
     address and other related configuration information such as the subnet 
     mask and default gateway. RFCs 2131 and 2132 define DHCP as an Internet 
     Engineering Task Force (IETF) standard based on Bootstrap Protocol 
     (BOOTP), a protocol with which DHCP shares many implementation details. 
     DHCP allows hosts to obtain required TCP/IP configuration information 
     from a DHCP server. 
       
     Examples : 
     
     Display details of all interfaces, including disabled interfaces:
     We can see whether an interface is enabled or disabled by looking
     at the flags, in particular the UP and RUNNING keywords.
           $ ifconfig -a
     
     For a quick summary of the interfaces, with some brief network stats,
     we can use the -s option:
           $ ifconfig -s   
     View network settings of an Ethernet adapter enp0s3:
           $ ifconfig enp0s3
           
     Enable enp0s3 interface: 
           $ ifconfig enp0s3 up 
           
     Disable enp0s3 interface:
           $ ifconfig enp0s3 down          
      

\end{verbatim}





\subsection{ping   : ping command}
\begin{verbatim}
    
    
    $ ping rt-ed.co.il # We can see that IP of the site is 212.199.114.183
      
      PING rt-ed.co.il (212.199.114.183) 56(84) bytes of data.
      64 bytes from 212.199.114.183.static.012.net.il (212.199.114.183):   
         icmp_seq=2 ttl=57 time=13.3 ms
      64 bytes from 212.199.114.183.static.012.net.il (212.199.114.183): 
         icmp_seq=3 ttl=57 time=36.1 ms


    $ ping 192.168.1.12
     
     PING 192.168.1.12 (192.168.1.12) 56(84) bytes of data.
     64 bytes from 192.168.1.12: icmp_seq=1 ttl=64 time=129 ms
     64 bytes from 192.168.1.12: icmp_seq=2 ttl=64 time=4.95 ms
     64 bytes from 192.168.1.12: icmp_seq=3 ttl=64 time=4.64 ms

    
    
    Send ICMP ECHO_REQUEST packets to network hosts.
    - Ping host:
        ping {{host}}

    - Ping a host only a specific number of times:
        ping -c {{count}} {{host}}

    - Ping host, specifying the interval in seconds between
      requests (default is 1 second):
        ping -i {{seconds}} {{host}}

    - Ping host without trying to lookup symbolic names for
      addresses:
        ping -n {{host}}

    - Ping host and ring the bell when a packet is received
      (if your terminal supports it):
        ping -a {{host}}

    - Also display a message if no response was received:
        ping -O {{host}}
    
\end{verbatim}

\subsection{hostname : system’s hostname}
\begin{verbatim}
    hostname command in Linux is used to obtain the DNS (Domain Name System)
    name and set the system’s hostname or NIS (Network Information System) 
    domain name.
    A hostname is a name given to a computer and attached to the network.
    Its main purpose is to uniquely identify over a network.
        $ hostname
            amit-vm-mint21
        $ hostname -i
            127.0.1.1
\end{verbatim}

\subsection{wget   : network downloader}
\begin{verbatim}

      GNU Wget is a free utility for non-interactive download of    
      files from the Web.
      It supports HTTP, HTTPS, and FTP protocols, as well as
      retrieval through HTTP proxies.
      Wget is non-interactive, meaning that it can work in the    
      background, while the user is not logged on.
      This allows you to start a retrieval and disconnect from
      the system, letting Wget finish the work.  By contrast,
      most of the Web browsers require constant user's presence,
      which can be a great hindrance when transferring a lot of     
      data.

     Wget provides a number of options allowing you to download 
     multiple files, resume downloads, limit the bandwidth, 
     recursive downloads, download in the background, mirror a 
     website, and much more.

     Mirror a site: 
         wget -m http://lwn.net

     Mirror a site (-r : recusive): 
         wget -r -np http://www.xml.com/ldd/chapter/book/

     #recursively download an online book 
     #for offline access. -np : no parent  
\end{verbatim}







\section{TextToSpeech }

\subsection{TextToSpeech..apps }
\begin{verbatim}
    say : converts text to audible speech using the GNUstep speech engine.
            sudo apt-get install gnustep-gui-runtime
            say "hello"

    festival : General multi-lingual speech synthesis system.
            sudo apt-get install festival
            echo "hello" | festival --tts
            
    spd-say : sends text-to-speech output request to speech-dispatcher
            sudo apt-get install speech-dispatcher
            spd-say "hello"
 
    espeak : is a multi-lingual software speech synthesizer.
            sudo apt-get install espeak
            espeak "hello"
            
      Example to see voices : 
         espeak --voices=variant  # to see avalable languages
         espeak -compile=en-scottish       
            

\end{verbatim}

\subsection{spd-say : text-to-speech output  }
\begin{verbatim}
       See: https://askubuntu.com/questions/501910/
       how-to-text-to-speech-output-using-command-line        
       
       Install : 
       sudo apt-get install speech-dispatcher
       spd-say "hello"
       
       
       See spd-say

       -r, --rate
              Set the rate of the speech (between -100 and +100, default: 0)

       -p, --pitch
              Set the pitch of the speech (between -100 and +100, default: 0)

       -i, --volume
              Set the volume (intensity) of the speech (between -100 and +100,
              default: 0)
              
       -l, --language
              Set the language (ISO code)
       
       -t, --voice-type
              Set  the preferred voice type (male1, male2, male3, female1, fe‐
              male2 female3, child_male, child_female)
              
       -L, --list-synthesis-voices
              Get the list of synthesis voices
              
        spd-say -t Robert "Hello, how do you do "      
        spd-say -t Pablo "Hello, how do you do "
        spd-say -t female2 "Hello, how do you do "
                
\end{verbatim}




\section{Appendix.}

\subsection{LaTeX     : editor}
\begin{verbatim}
  
 Bookmarks addition was based on following explanation : 
      https://michaelgoerz.net/notes/pdf-bookmarks-with-latex.html

sudo apt install texlive-full # install full LaTex editor

 To add bookmarks do following steps:
   1. sudo apt install texlive-full # install LaTex editor 
   2. cd ~/Documents/docs/c_for_embeded_pdf 
   3. touch c_for_embed.tex
   4. geany & c_for_embed.tex     # prepare LaTex document with bookmarks 
      definitions
   5. pdflatex c_for_embed.tex    # this command generates linux_admin.pdf with 
      desired bookmarks :)
 Note : the latex document can be edited with texmaker 
  

  see : https://artofproblemsolving.com/wiki/index.php/LaTeX
     
  To enlarge interface font run following command
  from terminal window :
      $ texmaker -dpiscale 1

     
\end{verbatim}



\subsection{Shell..built..in..commands     : shell built in commands}
\begin{verbatim}
     See : 
     https://www.gnu.org/software/bash/manual/html_node/Shell-Builtin-Commands.html
     
     Shell Builtin Commands
     Builtin commands are contained within the shell itself.
     When the name of a builtin command is used as the first word 
     of a simple command (see Simple Commands), the shell executes the command   
     directly, without invoking another program. 
     Builtin commands are necessary to implement functionality impossible or 
     inconvenient to obtain with separate utilities.
     
     Example: 
       $ type -a time
          time is a shell keyword # time is built in inside the shell
          time is /usr/bin/time   # also time is in this location
          time is /bin/time       # also time is in this location

\end{verbatim}

\subsection{shift+PgUp/PgDn     : page scroll in terminal}
\begin{verbatim}
     Page scroll in terminal
\end{verbatim}

\subsection{niceness  : process priority}
\begin{verbatim}
     To my knowledge, the Linux kernel does not change the niceness
     of a process, and I fail to see why it would since it doesn't have
     to to lower the priority of a process. 
     The niceness is an information given to the kernel, telling it how
     nice that process is willing to be. 
     The kernel scheduler is free to take this information into account
     the way it wants in order to change the priority of a process,
     it doesn't need to change its value.

     On the other hand, in user land, there are daemons like AND whose
     task is to renice processes according to rules set up by the admin.
     Do you have such a daemon installed on your server?

     Niceness values range from -20 (highest priority - most favorable to the    
     process) to 19 (lowest priority - least favorable to the process).

     Niceness can be changed with command renice for the running process. 

     The NI field in ps command shows the nice value of the process.
     The nice value determines the priority of the process. The higher the value, 
     the lower the priority--the "nicer" the process is to other processes. 
     The default nice value is 0 on Linux workstations. 
 
\end{verbatim}





\subsection{VT       : virtual terminal}
\begin{verbatim}
     
     
     Linux systems come with a bunch of virtual terminals,
     or VTs. Your graphical user interface on Ubuntu runs on VT2,
     and VT3 to VT6 allow you to login via command line.     
     
     When the graphics is stuck, you can still access
     the console window.

     To access consoles: [Ctrl][Alt][F3]/[F4]/[F5]
     
     It can be used to work with terminal tty3/tty4/tty5 only.
     In the text console, you can try to kill the guilty application

     Once this is done, you can go back to the graphic session
     by pressing [Ctrl][Alt][F2] (depends on distribution)     
     
\end{verbatim}



\subsection{ctrl+shift+arrowup/down     : scroll in terminal}
\begin{verbatim}
     Scroll in terminal window
\end{verbatim}


\subsection{Additional...shells }
\begin{verbatim}
     Different shells : 
     sh #bourne shell
     #bourne again shell :  bash
     sch # c shell , less recommended
     zsh # another shell 
           
\end{verbatim}

\subsection{Special...files }
\begin{verbatim}

     Runs automation scripts every time the terminal runs:
        .bashrc – remote control
     Runs automation scripts every time before the program run: 
        .profile 
           
\end{verbatim}



\subsection{inode    : Create files or set access/modif. times}
\begin{verbatim}
    https://en.wikipedia.org/wiki/Inode
    • The inode (index node) is a data structure in a Unix-style file 			  
      system that describes a file-system object such as a file or a
      directory. 
    • Each inode stores the attributes and disk block locations of the
      object's data. File-system object attributes may include metadata 
      (times of last change,access, modification), as well as owner and
      permission data
           
\end{verbatim}




\subsection{Disc..Partitions }

\begin{verbatim}

      /       root directory
      
      /bin/   Basic, essential system commands
      
      /boot/  Bootloader, Kernel images, initrd (initial RAM disk) 
              and configuration files
              
      /dev/   Files representing devices
      
      /etc/   System configuration files(as registry)
      
      /home/  User directories
      
      /lib/   Basic system shared libraries

      /lost+found Corrupt files the system tried to recover

      /media  Mount points for removable media 
              (/media/usbdisk, /media/cdrom)
      /mnt/   Mount points for temporarily mounted filesystems
      
      /opt/   Specific tools installed by the sysadmin 
              (for all users)
      /proc/  Access to system information (list of processes)
      
      /root/  root user home directory
      
      /sbin/  superuser only commands
      
      /sys/   System and device controls (cpu frequency,
              device power, etc.
              
      /tmp/   temporary files
      
      /usr/   Regular user tools (not essential to the system)
              /usr/bin/, /usr/lib/, /usr/sbin...
              
      /usr/local/  Specific software installed by the sys-admin
                   (often preferred to /opt/

      /var/   Data used by the system or system servers 
              /var/log/, /var/spool/mail (incoming mail), 
              /var/spool/lpd (print jobs)...

\end{verbatim}


\subsection{Globbing}
\begin{verbatim}
		Expansion : 
			~  : home directory
			[], ? , *
			!
			$

			“
			‘

			*  : any characters
			?  : exactly one character
             
             echo ?.txt     # ? replace exactly one character
             echo ??.txt
             
             ^ is fot not :
             $ echo [^0-9].txt
                 a.txt b.txt g.txt i.txt    
            
\end{verbatim} 
           
           
\end{document}
